\documentclass[UTF8,a4paper,fontset=fandol]{ctexart}
\usepackage[margin=1.7cm]{geometry}
\usepackage{xcolor}
\usepackage{hyperref}
\usepackage{booktabs}
\usepackage{longtable}
\usepackage{fancyhdr}
\usepackage{fvextra}
\hypersetup{colorlinks=true,linkcolor=blue,urlcolor=blue}
\pagestyle{fancy}
\fancyhf{}
\lhead{量化项目全部代码附录}
\rhead{\thepage}
\setlength{\headheight}{14pt}
\setlength{\parindent}{0pt}
\setlength{\parskip}{4pt}
\DefineVerbatimEnvironment{ProjectCode}{Verbatim}{
  fontsize=\scriptsize,
  numbers=left,
  numbersep=5pt,
  frame=single,
  framesep=2mm,
  breaklines=true,
  breakanywhere=true,
  tabsize=4
}
\title{量化项目全部代码附录}
\author{量化作业项目}
\date{2026-08-03}
\begin{document}
\maketitle
\tableofcontents
\clearpage

\section*{使用说明}
\addcontentsline{toc}{section}{使用说明}
本文件是项目源代码的只读汇总副本。生成过程没有改动任何原始代码文件。

文件数量：16；总代码行数：7,342。

推荐使用 XeLaTeX 编译。每个代码章节均记录原始相对路径、物理行数和 SHA-256 校验值。

\section*{文件清单}
\addcontentsline{toc}{section}{文件清单}
\begin{longtable}{r p{8.1cm} r}
\toprule
序号 & 文件 & 行数 \\
\midrule
\endhead
1 & \path{pipeline_code/downsample_tick_data_and_apply_adjustment.py} & 186 \\
2 & \path{pipeline_code/apply_adjustment_factor_to_existing_tables.py} & 116 \\
3 & \path{pipeline_code/construct_factors.py} & 163 \\
4 & \path{pipeline_code/evaluate_factors.py} & 442 \\
5 & \path{scripts/import_trade_data.py} & 186 \\
6 & \path{scripts/adjust_existing_processed.py} & 116 \\
7 & \path{scripts/build_kline_csv.py} & 135 \\
8 & \path{scripts/analyze_factors_backtest.py} & 1,627 \\
9 & \path{scripts/train_lstm_pytorch.py} & 885 \\
10 & \path{scripts/validate_outputs.py} & 409 \\
11 & \path{scripts/build_report.py} & 281 \\
12 & \path{scripts/build_explained_assignment_pdf.py} & 828 \\
13 & \path{scripts/build_final_illustrated_report.py} & 797 \\
14 & \path{scripts/build_factor_optimization_report.py} & 513 \\
15 & \path{scripts/build_factor_formula_report.py} & 376 \\
16 & \path{tests/test_pipeline.py} & 282 \\
\bottomrule
\end{longtable}
\clearpage
\section{1. \texttt{pipeline\_code/downsample\_tick\_data\_and\_apply\_adjustment.py}}
逐笔数据降采样、复权与日频/分钟频输出。

\textbf{物理行数：} 186\\
\textbf{SHA-256：} \texttt{b0ba5872d7a461458fd82a782912a687a7a3842b8c6d5f53d7b759fb02fcb70f}

\begin{ProjectCode}
#!/usr/bin/env python3
"""Downsample tick trades and apply adjustment factors to OHLC prices.

Prices are adjusted as ``raw_price / 100 * adjfactor``.  Turnover amounts use
the actually traded (unadjusted) price, which is the conventional definition.
"""

from __future__ import annotations

import argparse
from pathlib import Path

import numpy as np
import pandas as pd


FIELDS = [
    "open", "high", "low", "close", "volume", "trade_count", "amount",
    "buy_volume", "sell_volume", "buy_amount", "sell_amount",
]
PRICE_FIELDS = ["open", "high", "low", "close"]
ZERO_FIELDS = [field for field in FIELDS if field not in PRICE_FIELDS]


def trading_minutes() -> list[int]:
    """Return the assignment minute grid, including the 15:00 close."""
    minutes = list(range(915, 926))
    for hour, start, end in ((9, 30, 59), (10, 0, 59), (11, 0, 30),
                             (13, 0, 59), (14, 0, 59), (15, 0, 0)):
        minutes.extend(hour * 100 + minute for minute in range(start, end + 1))
    return minutes


def load_adjustment(path: Path) -> pd.DataFrame:
    factors = pd.read_pickle(path).copy()
    factors.index = pd.to_datetime(factors.index).strftime("%Y%m%d")
    factors.columns = [str(column).split(".")[0].strip() for column in factors.columns]
    if factors.columns.duplicated().any():
        raise ValueError("Duplicate six-digit codes in adjustment-factor table")
    return factors.sort_index().ffill().bfill()


def summarize_stock(csv_path: Path, adjustment: float) -> tuple[dict[str, float], pd.DataFrame]:
    df = pd.read_csv(
        csv_path,
        dtype={"Time": "int64", "Price": "int64", "Volume": "int64", "BSFlag": "int8"},
    )
    if df.empty:
        return ({field: np.nan for field in FIELDS},
                pd.DataFrame(index=trading_minutes(), columns=FIELDS, dtype=float))

    required = {"Time", "Price", "Volume", "BSFlag"}
    if set(df.columns) != required:
        raise ValueError(f"Unexpected columns in {csv_path}: {list(df.columns)}")
    if not df["BSFlag"].isin([0, 1, 2]).all():
        raise ValueError(f"Unexpected BSFlag in {csv_path}")

    df = df.sort_values("Time", kind="stable")
    df["minute"] = df["Time"] // 100000
    raw_price = df["Price"].astype(float) / 100.0
    adjusted_price = raw_price * adjustment
    amount = raw_price * df["Volume"]
    buy = df["BSFlag"].eq(0)
    sell = df["BSFlag"].eq(1)

    daily = {
        "open": float(adjusted_price.iloc[0]),
        "high": float(adjusted_price.max()),
        "low": float(adjusted_price.min()),
        "close": float(adjusted_price.iloc[-1]),
        "volume": int(df["Volume"].sum()),
        "trade_count": int(len(df)),
        "amount": float(amount.sum()),
        "buy_volume": int(df.loc[buy, "Volume"].sum()),
        "sell_volume": int(df.loc[sell, "Volume"].sum()),
        "buy_amount": float(amount[buy].sum()),
        "sell_amount": float(amount[sell].sum()),
    }

    groups = df.groupby("minute", sort=True)
    minute = pd.DataFrame({
        "open": groups["Price"].first() / 100.0 * adjustment,
        "high": groups["Price"].max() / 100.0 * adjustment,
        "low": groups["Price"].min() / 100.0 * adjustment,
        "close": groups["Price"].last() / 100.0 * adjustment,
        "volume": groups["Volume"].sum(),
        "trade_count": groups.size(),
        "amount": amount.groupby(df["minute"]).sum(),
        "buy_volume": df["Volume"].where(buy, 0).groupby(df["minute"]).sum(),
        "sell_volume": df["Volume"].where(sell, 0).groupby(df["minute"]).sum(),
        "buy_amount": amount.where(buy, 0).groupby(df["minute"]).sum(),
        "sell_amount": amount.where(sell, 0).groupby(df["minute"]).sum(),
    })
    minute.index.name = "minute"
    return daily, minute


def write_table(table: pd.DataFrame, path: Path) -> None:
    path.parent.mkdir(parents=True, exist_ok=True)
    table.to_csv(path, index_label=table.index.name or "index")


def import_data(input_root: Path, output_root: Path, factor_path: Path,
                limit_dates: int | None = None) -> None:
    trade_root = input_root / "TRADE"
    date_dirs = sorted(path for path in trade_root.iterdir()
                       if path.is_dir() and path.name.isdigit() and len(path.name) == 8)
    if limit_dates is not None:
        date_dirs = date_dirs[:limit_dates]
    if not date_dirs:
        raise SystemExit(f"No YYYYMMDD directories under {trade_root}")

    all_codes = sorted({path.stem for day in date_dirs for path in day.glob("*.csv")})
    factors = load_adjustment(factor_path)
    missing_codes = sorted(set(all_codes) - set(factors.columns))
    missing_dates = sorted({day.name for day in date_dirs} - set(factors.index))
    if missing_codes or missing_dates:
        raise ValueError(f"Adjustment coverage missing codes={missing_codes}, dates={missing_dates}")

    minutes = trading_minutes()
    daily_rows: dict[str, list[pd.Series]] = {field: [] for field in FIELDS}
    previous_close = pd.Series(index=all_codes, dtype=float)

    for day_number, day_dir in enumerate(date_dirs, start=1):
        date = day_dir.name
        day_factor = factors.loc[date, all_codes].astype(float)
        daily_by_code = pd.DataFrame(index=all_codes, columns=FIELDS, dtype=float)
        minute_by_field = {
            field: pd.DataFrame(index=minutes, columns=all_codes, dtype=float)
            for field in FIELDS
        }

        for csv_path in sorted(day_dir.glob("*.csv")):
            code = csv_path.stem
            adjustment = float(day_factor[code])
            if not np.isfinite(adjustment) or adjustment <= 0:
                raise ValueError(f"Invalid adjustment {adjustment} for {date}/{code}")
            daily, minute = summarize_stock(csv_path, adjustment)
            for field, value in daily.items():
                daily_by_code.at[code, field] = value
            for field in FIELDS:
                minute_by_field[field].loc[minute.index, code] = minute[field].to_numpy()

        # A missing daily price is carried from the preceding adjusted close.
        for field in PRICE_FIELDS:
            daily_by_code[field] = daily_by_code[field].fillna(previous_close)
        for field in ZERO_FIELDS:
            daily_by_code[field] = daily_by_code[field].fillna(0)

        # A no-trade minute uses the preceding minute close. Leading empty
        # auction minutes use yesterday's close; only the first day may bfill.
        close_filled = minute_by_field["close"].ffill().fillna(previous_close).bfill()
        for field in PRICE_FIELDS:
            minute_by_field[field] = minute_by_field[field].fillna(close_filled)
        for field in ZERO_FIELDS:
            minute_by_field[field] = minute_by_field[field].fillna(0)

        for field in FIELDS:
            daily_rows[field].append(daily_by_code[field].rename(date))
            table = minute_by_field[field]
            table.index.name = "minute"
            write_table(table, output_root / "minute" / field / f"{date}.csv")

        previous_close = daily_by_code["close"].combine_first(previous_close)
        print(f"[{day_number}/{len(date_dirs)}] imported {date}", flush=True)

    for field, rows in daily_rows.items():
        table = pd.DataFrame(rows)
        table.index.name = "date"
        write_table(table, output_root / "daily" / f"{field}.csv")


def main() -> None:
    project_root = Path(__file__).resolve().parents[1]

    parser = argparse.ArgumentParser(description=__doc__)
    parser.add_argument("--input", type=Path, default=project_root / "data")
    parser.add_argument("--output", type=Path, default=project_root / "processed")
    parser.add_argument("--adjfactor", type=Path, default=project_root / "adjfactor.pkl")
    parser.add_argument("--limit-dates", type=int)
    args = parser.parse_args()
    import_data(args.input, args.output, args.adjfactor, args.limit_dates)


if __name__ == "__main__":
    main()
\end{ProjectCode}
\clearpage
\section{2. \texttt{pipeline\_code/apply\_adjustment\_factor\_to\_existing\_tables.py}}
对已有处理结果应用复权因子并重建相关输出。

\textbf{物理行数：} 116\\
\textbf{SHA-256：} \texttt{bb852264b4f5060ec339f6df084e7d6a0ee52e5baebf00a46dca529f765a3453}

\begin{ProjectCode}
#!/usr/bin/env python3
"""Apply adjustment factors to existing unadjusted daily/minute tables.

This acceleration path is mathematically equivalent to the raw importer.  It
also reconstructs missing minute prices from volume, so leading empty auction
minutes use the prior daily close instead of information from the current day.
"""

from __future__ import annotations

import argparse
import shutil
from pathlib import Path

import numpy as np
import pandas as pd

from downsample_tick_data_and_apply_adjustment import PRICE_FIELDS, ZERO_FIELDS, load_adjustment


def load_wide(path: Path, index_name: str) -> pd.DataFrame:
    table = pd.read_csv(path, dtype={index_name: str})
    table[index_name] = table[index_name].astype(str).str.replace(r"\.0$", "", regex=True)
    return table.set_index(index_name)


def write_wide(table: pd.DataFrame, path: Path) -> None:
    path.parent.mkdir(parents=True, exist_ok=True)
    table.to_csv(path, index_label=table.index.name)


def adjustment_for(factors: pd.DataFrame, dates: pd.Index, codes: pd.Index) -> pd.DataFrame:
    keyed_dates = pd.Index(pd.to_datetime(dates).strftime("%Y%m%d"))
    result = factors.reindex(index=keyed_dates, columns=codes)
    result.index = dates
    if result.isna().any().any():
        raise ValueError("Adjustment factor does not cover processed panel")
    return result


def rebuild(source: Path, output: Path, factor_path: Path) -> None:
    factors = load_adjustment(factor_path)
    output.mkdir(parents=True, exist_ok=True)

    # Copy the seven non-price field trees unchanged: adjustment affects prices,
    # not shares, trade counts, or actually transacted amounts.
    for field in ZERO_FIELDS:
        src_daily = source / "daily" / f"{field}.csv"
        dst_daily = output / "daily" / f"{field}.csv"
        daily = load_wide(src_daily, "date").astype(float).fillna(0)
        daily.index.name = "date"
        write_wide(daily, dst_daily)
        src_minute = source / "minute" / field
        dst_minute = output / "minute" / field
        if dst_minute.exists():
            shutil.rmtree(dst_minute)
        shutil.copytree(src_minute, dst_minute, copy_function=shutil.copy2)
        print(f"copied {field}", flush=True)

    raw_daily: dict[str, pd.DataFrame] = {}
    for field in PRICE_FIELDS:
        table = load_wide(source / "daily" / f"{field}.csv", "date").astype(float)
        raw_daily[field] = table

    raw_daily_volume = load_wide(source / "daily" / "volume.csv", "date").astype(float)
    has_trades = raw_daily_volume.notna()
    factor_matrix = adjustment_for(factors, raw_daily["close"].index,
                                   raw_daily["close"].columns)
    adjusted_daily_close = (raw_daily["close"] * factor_matrix).where(has_trades).ffill()
    adjusted_daily: dict[str, pd.DataFrame] = {}
    for field in PRICE_FIELDS:
        adjusted = (raw_daily[field] * factor_matrix).where(has_trades).fillna(adjusted_daily_close)
        adjusted.index.name = "date"
        adjusted_daily[field] = adjusted
        write_wide(adjusted, output / "daily" / f"{field}.csv")
        print(f"adjusted daily/{field}", flush=True)

    dates = list(raw_daily["close"].index)
    codes = raw_daily["close"].columns
    previous_adjusted_close = pd.Series(index=codes, dtype=float)
    for number, date in enumerate(dates, start=1):
        volume = load_wide(source / "minute" / "volume" / f"{date}.csv", "minute")
        volume = volume.reindex(columns=codes).astype(float)
        traded = volume.gt(0)

        reconstructed: dict[str, pd.DataFrame] = {}
        factor = factors.loc[pd.to_datetime(date).strftime("%Y%m%d"), codes].astype(float)
        if not np.isfinite(factor).all():
            raise ValueError(f"Invalid adjustment factor on {date}")
        raw_close = load_wide(source / "minute" / "close" / f"{date}.csv", "minute")
        adjusted_close = raw_close.reindex(columns=codes).astype(float).where(traded).mul(factor, axis=1)
        close_filled = adjusted_close.ffill().fillna(previous_adjusted_close).bfill()
        reconstructed["close"] = close_filled
        for field in ("open", "high", "low"):
            actual = load_wide(source / "minute" / field / f"{date}.csv", "minute")
            reconstructed[field] = (actual.reindex(columns=codes).astype(float)
                                    .where(traded).mul(factor, axis=1).fillna(close_filled))
        for field, table in reconstructed.items():
            table.index.name = "minute"
            write_wide(table, output / "minute" / field / f"{date}.csv")

        previous_adjusted_close = adjusted_daily["close"].loc[date].combine_first(previous_adjusted_close)
        print(f"[{number}/{len(dates)}] repriced {date}", flush=True)


def main() -> None:
    parser = argparse.ArgumentParser(description=__doc__)
    parser.add_argument("--source", type=Path, required=True)
    parser.add_argument("--output", type=Path, default=Path("processed"))
    parser.add_argument("--adjfactor", type=Path, default=Path("adjfactor.pkl"))
    args = parser.parse_args()
    rebuild(args.source, args.output, args.adjfactor)


if __name__ == "__main__":
    main()
\end{ProjectCode}
\clearpage
\section{3. \texttt{pipeline\_code/construct\_factors.py}}
构建七个日频因子并输出因子长表和元数据。

\textbf{物理行数：} 163\\
\textbf{SHA-256：} \texttt{370fc3ab7acad7e11b71e7a07d240ff35ababdad5e72b2241034ca9a81df2090}

\begin{ProjectCode}
#!/usr/bin/env python3
"""Construct seven daily factors from processed field tables and save CSV files."""

from __future__ import annotations

import argparse
import warnings
from pathlib import Path

import numpy as np
import pandas as pd


warnings.filterwarnings("ignore", category=FutureWarning,
                        message="The previous implementation of stack")

FACTOR_META = {
    "amount_mean_sd_log": {
        "name_zh": "成交额多窗口均值离散度",
        "formula": "log(1 + sd(mean_1d(amount), mean_5d(amount), mean_10d(amount), mean_20d(amount)))",
        "kind": "题目示例因子",
    },
    "buy_sell_imbalance_surprise_10d": {
        "name_zh": "10日主动买卖不平衡冲击",
        "formula": "active_imbalance_t - mean_prev_10d(active_imbalance)",
        "kind": "自建因子",
    },
    "intraday_range_10d": {
        "name_zh": "10日日内振幅",
        "formula": "mean_10d((high - low) / close)",
        "kind": "自建因子",
    },
    "reversal_5d": {
        "name_zh": "5日反转",
        "formula": "-(close / close_lag_5d - 1)",
        "kind": "自建因子",
    },
    "avg_trade_size_surprise_20d": {
        "name_zh": "20日单笔成交规模异常",
        "formula": "zscore_20d(log(1 + amount / max(trade_count, 1)))",
        "kind": "新增因子",
    },
    "price_volume_pressure_10d": {
        "name_zh": "10日价量压力",
        "formula": "mean_10d(return_1d * diff(log(1 + volume)))",
        "kind": "新增因子",
    },
    "gap_intraday_divergence_5d": {
        "name_zh": "5日隔夜-日内收益背离",
        "formula": "mean_5d((open / close_lag_1d - 1) - (close / open - 1))",
        "kind": "新增因子",
    },
}


def load_wide(processed: Path, field: str) -> pd.DataFrame:
    path = processed / "daily" / f"{field}.csv"
    table = pd.read_csv(path, dtype={"date": str})
    table["date"] = pd.to_datetime(table["date"])
    table = table.set_index("date").astype(float).sort_index()
    table.columns = table.columns.astype(str).str.zfill(6)
    return table


def save_wide(table: pd.DataFrame, path: Path) -> None:
    path.parent.mkdir(parents=True, exist_ok=True)
    copy = table.copy()
    copy.index = copy.index.strftime("%Y%m%d")
    copy.index.name = "date"
    copy.to_csv(path)


def build_factors(processed: Path, output: Path) -> dict[str, pd.DataFrame]:
    open_price = load_wide(processed, "open")
    close = load_wide(processed, "close")
    high = load_wide(processed, "high")
    low = load_wide(processed, "low")
    amount = load_wide(processed, "amount")
    volume = load_wide(processed, "volume")
    trade_count = load_wide(processed, "trade_count")
    buy_volume = load_wide(processed, "buy_volume")
    sell_volume = load_wide(processed, "sell_volume")

    amount_means = [amount.rolling(window, min_periods=window).mean()
                    for window in (1, 5, 10, 20)]
    stacked_means = np.stack([table.to_numpy() for table in amount_means], axis=0)
    amount_factor_values = np.full(amount.shape, np.nan, dtype=float)
    valid = np.isfinite(stacked_means).all(axis=0)
    amount_factor_values[valid] = np.log1p(
        np.std(stacked_means[:, valid], axis=0, ddof=1)
    )

    # Exclude auction volume from the denominator because BSFlag=2 is not an
    # aggressive buy or sell. The predictive component is the current order-
    # flow shock relative to the previous ten trading days. A five-day simple
    # average was empirically too smooth and cancelled the short-lived signal.
    active_volume = (buy_volume + sell_volume).replace(0, np.nan)
    imbalance = (buy_volume - sell_volume).div(active_volume)
    imbalance_history = imbalance.shift(1).rolling(10, min_periods=10).mean()
    daily_return = close.pct_change(fill_method=None)
    log_volume_change = np.log1p(volume).diff()
    log_average_trade_size = np.log1p(
        amount.div(trade_count.clip(lower=1))
    )
    average_trade_size_mean = log_average_trade_size.rolling(
        20, min_periods=20
    ).mean()
    average_trade_size_std = log_average_trade_size.rolling(
        20, min_periods=20
    ).std()
    overnight_gap = open_price.div(close.shift(1).replace(0, np.nan)) - 1
    intraday_return = close.div(open_price.replace(0, np.nan)) - 1
    factors = {
        "amount_mean_sd_log": pd.DataFrame(
            amount_factor_values, index=amount.index, columns=amount.columns
        ),
        "buy_sell_imbalance_surprise_10d": imbalance - imbalance_history,
        "intraday_range_10d": (
            (high - low) / close.clip(lower=1e-12)
        ).rolling(10, min_periods=10).mean(),
        "reversal_5d": -(close / close.shift(5) - 1),
        "avg_trade_size_surprise_20d": (
            log_average_trade_size - average_trade_size_mean
        ).div(average_trade_size_std.replace(0, np.nan)),
        "price_volume_pressure_10d": (
            daily_return * log_volume_change
        ).rolling(10, min_periods=10).mean(),
        "gap_intraday_divergence_5d": (
            overnight_gap - intraday_return
        ).rolling(5, min_periods=5).mean(),
    }

    daily_dir = output / "daily"
    for name, table in factors.items():
        save_wide(table, daily_dir / f"{name}.csv")

    long_parts = [table.stack().rename(name)
                  for name, table in factors.items()]
    factor_long = pd.concat(long_parts, axis=1).reset_index(names=["date", "code"])
    factor_long["date"] = factor_long["date"].dt.strftime("%Y%m%d")
    factor_long["code"] = factor_long["code"].astype(str).str.zfill(6)
    output.mkdir(parents=True, exist_ok=True)
    factor_long.to_csv(output / "factor_long.csv", index=False, quoting=1)

    metadata = pd.DataFrame([
        {"factor": factor, **meta} for factor, meta in FACTOR_META.items()
    ])
    metadata.to_csv(output / "factor_metadata.csv", index=False)
    return factors


def main() -> None:
    project_root = Path(__file__).resolve().parents[1]
    parser = argparse.ArgumentParser(description=__doc__)
    parser.add_argument("--processed", type=Path, default=project_root / "processed")
    parser.add_argument("--output", type=Path, default=project_root / "factors")
    args = parser.parse_args()
    factors = build_factors(args.processed.resolve(), args.output.resolve())
    print(f"constructed {len(factors)} factors in {args.output.resolve()}")


if __name__ == "__main__":
    main()
\end{ProjectCode}
\clearpage
\section{4. \texttt{pipeline\_code/evaluate\_factors.py}}
因子清洗、中性化接口、IC/Rank IC、分组收益与因子图。

\textbf{物理行数：} 442\\
\textbf{SHA-256：} \texttt{0df6f9751b9864fe1e636834afd145ab02078a11d814e9abfdf654a64b7eed76}

\begin{ProjectCode}
#!/usr/bin/env python3
"""Evaluate saved daily factors with IC, Rank IC, ICIR, and quantile CSV outputs."""

from __future__ import annotations

import argparse
import json
import math
import os
import tempfile
import warnings
from pathlib import Path

os.environ.setdefault(
    "MPLCONFIGDIR",
    str(Path(tempfile.gettempdir()) / "quant-assignment-matplotlib"),
)

import matplotlib  # noqa: E402
import numpy as np
import pandas as pd
from scipy.stats import spearmanr

matplotlib.use("Agg")
import matplotlib.pyplot as plt  # noqa: E402

from construct_factors import FACTOR_META, load_wide


warnings.filterwarnings("ignore", category=FutureWarning,
                        message="The previous implementation of stack")
TRADING_DAYS = 252


def load_factor_panels(factor_root: Path) -> dict[str, pd.DataFrame]:
    factors: dict[str, pd.DataFrame] = {}
    for name in FACTOR_META:
        path = factor_root / "daily" / f"{name}.csv"
        table = pd.read_csv(path, dtype={"date": str})
        table["date"] = pd.to_datetime(table["date"])
        table = table.set_index("date").astype(float).sort_index()
        table.columns = table.columns.astype(str).str.zfill(6)
        factors[name] = table
    return factors


def load_risk_exposures(path: Path) -> pd.DataFrame:
    """Load true industry and market-cap exposures for neutralization.

    Accepted schemas are either date-specific
    ``date,code,industry,market_cap`` or static
    ``code,industry,market_cap``.  Codes are never inferred because the
    assignment explicitly states that identifiers are obfuscated.
    """
    table = pd.read_csv(path, dtype={"date": str, "code": str, "industry": str})
    required = {"code", "industry", "market_cap"}
    if not required.issubset(table.columns):
        raise ValueError(
            f"{path} must contain code, industry, market_cap"
            " and may optionally contain date"
        )
    table = table.copy()
    table["code"] = table["code"].astype(str).str.replace(r"\.0$", "", regex=True).str.zfill(6)
    table["industry"] = table["industry"].astype(str).str.strip()
    table["market_cap"] = pd.to_numeric(table["market_cap"], errors="coerce")
    if "date" in table:
        table["date"] = pd.to_datetime(table["date"], errors="coerce")
        key = ["date", "code"]
    else:
        key = ["code"]
    if table[key].isna().any().any() or table.duplicated(key).any():
        raise ValueError(f"Invalid or duplicate risk-exposure keys in {path}")
    if table["industry"].eq("").any() or table["market_cap"].le(0).any():
        raise ValueError(f"Industry must be non-empty and market_cap positive in {path}")
    return table


def neutralize_cross_section(
    values: pd.Series,
    industry: pd.Series,
    market_cap: pd.Series,
) -> pd.Series:
    """Residualize one cross-section against industry and log market cap."""
    frame = pd.DataFrame({
        "value": pd.to_numeric(values, errors="coerce"),
        "industry": industry.astype(str),
        "market_cap": pd.to_numeric(market_cap, errors="coerce"),
    }, index=values.index)
    valid = (
        np.isfinite(frame["value"])
        & np.isfinite(frame["market_cap"])
        & frame["market_cap"].gt(0)
        & frame["industry"].ne("")
    )
    frame = frame.loc[valid]
    if len(frame) < 20:
        raise ValueError("Fewer than 20 valid stocks for risk neutralization")
    log_cap = np.log(frame["market_cap"].to_numpy(float))
    log_cap_std = log_cap.std(ddof=1)
    if not np.isfinite(log_cap_std) or log_cap_std == 0:
        raise ValueError("Market-cap exposure has zero cross-sectional variance")
    log_cap = (log_cap - log_cap.mean()) / log_cap_std
    industry_dummies = pd.get_dummies(
        frame["industry"], prefix="industry", drop_first=True, dtype=float
    )
    design = np.column_stack([
        np.ones(len(frame), dtype=float),
        log_cap,
        industry_dummies.to_numpy(float),
    ])
    if len(frame) <= design.shape[1] + 5:
        raise ValueError("Too many industry groups for the available cross-section")
    coefficients, *_ = np.linalg.lstsq(
        design, frame["value"].to_numpy(float), rcond=None
    )
    residual = frame["value"].to_numpy(float) - design @ coefficients
    result = pd.Series(np.nan, index=values.index, dtype=float)
    result.loc[frame.index] = residual
    return result


def merge_day_exposures(day: pd.DataFrame, exposures: pd.DataFrame) -> pd.DataFrame:
    """Attach exposures and require at least 90% coverage for a trading day."""
    if "date" in exposures:
        merged = day.merge(exposures, on=["date", "code"], how="left")
    else:
        merged = day.merge(exposures, on="code", how="left")
    covered = (
        merged["industry"].notna()
        & pd.to_numeric(merged["market_cap"], errors="coerce").gt(0)
    )
    coverage = float(covered.mean()) if len(merged) else 0.0
    if coverage < 0.90:
        date = day["date"].iloc[0] if len(day) else "unknown"
        raise ValueError(
            f"Risk-exposure coverage {coverage:.1%} is below 90% on {date}"
        )
    return merged.loc[covered].copy()


def safe_ratio(mean: float, std: float, annualize: bool = False) -> float:
    if not np.isfinite(std) or std == 0:
        return np.nan
    value = mean / std
    return value * math.sqrt(TRADING_DAYS) if annualize else value


def max_drawdown(nav: pd.Series | np.ndarray) -> float:
    values = np.asarray(nav, dtype=float)
    if not len(values):
        return np.nan
    return float(np.nanmin(values / np.maximum.accumulate(values) - 1.0))


def compute_factor_correlation(panel: pd.DataFrame) -> pd.DataFrame:
    """Return pooled cross-sectional correlations of cleaned factor z-scores."""
    wide = panel.pivot_table(
        index=["date", "code"],
        columns="factor",
        values="factor_z",
        aggfunc="first",
    )
    correlation = wide.corr(min_periods=20)
    correlation.index.name = "factor"
    return correlation


def build_quantile_summary(quantile: pd.DataFrame) -> pd.DataFrame:
    """Aggregate each factor/quantile return path into a compact table."""
    rows = []
    for (factor, quantile_number), group in quantile.groupby(
        ["factor", "quantile"], sort=True
    ):
        group = group.sort_values("date")
        returns = group["quantile_return"].fillna(0.0)
        nav = (1 + returns).cumprod()
        rows.append({
            "factor": factor,
            "quantile": int(quantile_number),
            "n_dates": len(group),
            "mean_daily_return": returns.mean(),
            "annual_return": nav.iloc[-1] ** (TRADING_DAYS / len(group)) - 1,
            "total_return": nav.iloc[-1] - 1,
            "annual_volatility": returns.std(ddof=1) * math.sqrt(TRADING_DAYS),
            "max_drawdown": max_drawdown(nav),
        })
    return pd.DataFrame(rows)


def write_factor_figures(
    ic_series: pd.DataFrame,
    quantile: pd.DataFrame,
    long_short: pd.DataFrame,
    correlation: pd.DataFrame,
    output: Path,
) -> None:
    """Write per-factor diagnostics and combined correlation/long-short plots."""
    figures = output / "figures"
    figures.mkdir(parents=True, exist_ok=True)
    for factor in sorted(quantile["factor"].unique()):
        factor_quantile = quantile.loc[quantile["factor"].eq(factor)].copy()
        factor_ic = ic_series.loc[ic_series["factor"].eq(factor)].sort_values("date")
        figure, axes = plt.subplots(1, 3, figsize=(14, 4.2))

        mean_return = factor_quantile.groupby("quantile")["quantile_return"].mean()
        axes[0].bar(mean_return.index.astype(str), mean_return.to_numpy() * 10_000)
        axes[0].axhline(0, color="0.25", linewidth=0.8)
        axes[0].set_title("Mean daily return by quantile")
        axes[0].set_xlabel("Quantile")
        axes[0].set_ylabel("Basis points")

        for quantile_number, group in factor_quantile.groupby("quantile", sort=True):
            axes[1].plot(
                group.sort_values("date")["date"],
                group.sort_values("date")["nav"],
                label=f"Q{int(quantile_number)}",
                linewidth=1.2,
            )
        axes[1].set_title("Quantile cumulative NAV")
        axes[1].set_ylabel("NAV")
        axes[1].legend(ncol=2, fontsize=8)

        rolling_rank_ic = factor_ic["rank_ic"].rolling(20, min_periods=10).mean()
        axes[2].plot(factor_ic["date"], rolling_rank_ic, linewidth=1.3)
        axes[2].axhline(0, color="0.25", linewidth=0.8)
        axes[2].set_title("20-day rolling Rank IC")
        axes[2].set_ylabel("Rank IC")
        for axis in axes[1:]:
            axis.tick_params(axis="x", labelrotation=30, labelsize=8)
        figure.suptitle(factor)
        figure.tight_layout()
        figure.savefig(figures / f"{factor}_diagnostics.png", dpi=160)
        plt.close(figure)

    figure, axis = plt.subplots(figsize=(7.5, 6.2))
    image = axis.imshow(correlation.to_numpy(), vmin=-1, vmax=1, cmap="coolwarm")
    axis.set_xticks(range(len(correlation.columns)), correlation.columns, rotation=35, ha="right")
    axis.set_yticks(range(len(correlation.index)), correlation.index)
    for row in range(len(correlation.index)):
        for column in range(len(correlation.columns)):
            value = correlation.iloc[row, column]
            axis.text(column, row, f"{value:.2f}", ha="center", va="center", fontsize=9)
    axis.set_title("Factor correlation")
    figure.colorbar(image, ax=axis, shrink=0.82)
    figure.tight_layout()
    figure.savefig(figures / "factor_correlation_heatmap.png", dpi=160)
    plt.close(figure)

    figure, axis = plt.subplots(figsize=(9, 5))
    for factor, group in long_short.groupby("factor", sort=True):
        group = group.sort_values("date")
        axis.plot(group["date"], group["long_short_nav"], label=factor, linewidth=1.2)
    axis.axhline(1, color="0.25", linewidth=0.8)
    axis.set_title("Q5 minus Q1 cumulative NAV")
    axis.set_ylabel("NAV")
    axis.legend(fontsize=8)
    axis.tick_params(axis="x", labelrotation=30, labelsize=8)
    figure.tight_layout()
    figure.savefig(figures / "factor_long_short_nav.png", dpi=160)
    plt.close(figure)


def clean_factor_panel(
    factors: dict[str, pd.DataFrame],
    next_return: pd.DataFrame,
    risk_exposures: pd.DataFrame | None = None,
) -> pd.DataFrame:
    returns_long = next_return.stack().rename("next_ret").reset_index()
    returns_long.columns = ["date", "code", "next_ret"]
    parts: list[pd.DataFrame] = []
    for factor, table in factors.items():
        raw = table.stack().rename("factor_value").reset_index()
        raw.columns = ["date", "code", "factor_value"]
        raw["factor"] = factor
        panel = raw.merge(returns_long, on=["date", "code"], how="left")
        clean_days = []
        for _, day in panel.groupby("date", sort=True):
            day = day.copy()
            if risk_exposures is not None:
                day = merge_day_exposures(day, risk_exposures)
            valid = np.isfinite(day["factor_value"]) & np.isfinite(day["next_ret"])
            if valid.sum() < 20:
                continue
            lo, hi = day.loc[valid, "factor_value"].quantile([0.01, 0.99])
            day.loc[valid, "factor_clean"] = day.loc[valid, "factor_value"].clip(lo, hi)
            if risk_exposures is not None:
                day.loc[valid, "factor_neutral"] = neutralize_cross_section(
                    day.loc[valid, "factor_clean"],
                    day.loc[valid, "industry"],
                    day.loc[valid, "market_cap"],
                )
            else:
                day.loc[valid, "factor_neutral"] = day.loc[valid, "factor_clean"]
            valid &= np.isfinite(day["factor_neutral"])
            mean = day.loc[valid, "factor_neutral"].mean()
            std = day.loc[valid, "factor_neutral"].std(ddof=1)
            if not np.isfinite(std) or std == 0:
                continue
            day.loc[valid, "factor_z"] = (
                day.loc[valid, "factor_neutral"] - mean
            ) / std
            day.loc[valid, "neutralization_applied"] = risk_exposures is not None
            clean_days.append(day.loc[valid])
        if clean_days:
            parts.append(pd.concat(clean_days, ignore_index=True))
    if not parts:
        raise ValueError("No valid factor evaluation rows")
    return pd.concat(parts, ignore_index=True)


def evaluate(panel: pd.DataFrame, output: Path) -> pd.DataFrame:
    output.mkdir(parents=True, exist_ok=True)
    ic_rows = []
    quantile_rows = []
    for (factor, date), day in panel.groupby(["factor", "date"], sort=True):
        x = day["factor_z"].to_numpy(float)
        y = day["next_ret"].to_numpy(float)
        ic = np.corrcoef(x, y)[0, 1] if np.std(x) and np.std(y) else np.nan
        rank_ic = spearmanr(x, y).statistic if np.std(x) and np.std(y) else np.nan
        ic_rows.append({"factor": factor, "date": date, "n_stock": len(day),
                        "ic": ic, "rank_ic": rank_ic})
        ranks = day["factor_z"].rank(method="first", pct=True)
        quantiles = np.ceil(ranks * 5).clip(1, 5).astype(int)
        grouped = day.assign(quantile=quantiles).groupby("quantile")["next_ret"]
        for quantile, values in grouped:
            quantile_rows.append({"factor": factor, "date": date,
                                  "quantile": int(quantile),
                                  "quantile_return": values.mean(),
                                  "n_stock": len(values)})

    ic_series = pd.DataFrame(ic_rows).sort_values(["factor", "date"])
    quantile = pd.DataFrame(quantile_rows).sort_values(["factor", "quantile", "date"])
    quantile["nav"] = quantile.groupby(["factor", "quantile"])["quantile_return"].transform(
        lambda series: (1 + series.fillna(0)).cumprod())
    pivot = quantile.pivot_table(index=["factor", "date"], columns="quantile",
                                 values="quantile_return").reset_index()
    pivot["long_short_return"] = pivot[5] - pivot[1]
    pivot = pivot.sort_values(["factor", "date"])
    pivot["long_short_nav"] = pivot.groupby("factor")["long_short_return"].transform(
        lambda series: (1 + series.fillna(0)).cumprod())

    summaries = []
    for factor, values in ic_series.groupby("factor"):
        long_short = pivot.loc[pivot["factor"].eq(factor), "long_short_return"].dropna()
        ic = values["ic"].dropna()
        rank_ic = values["rank_ic"].dropna()
        ir = safe_ratio(ic.mean(), ic.std(ddof=1))
        rank_ir = safe_ratio(rank_ic.mean(), rank_ic.std(ddof=1))
        long_short_nav = (1 + long_short).cumprod()
        summaries.append({
            "factor": factor,
            "n_dates": len(ic),
            "ic": ic.mean(),
            "ic_std": ic.std(ddof=1),
            "ir": ir,
            "icir": ir * math.sqrt(TRADING_DAYS),
            "ic_positive_rate": (ic > 0).mean(),
            "rank_ic": rank_ic.mean(),
            "rank_ic_std": rank_ic.std(ddof=1),
            "rank_ir": rank_ir,
            "rank_icir": rank_ir * math.sqrt(TRADING_DAYS),
            "rank_ic_positive_rate": (rank_ic > 0).mean(),
            "long_short_ir": safe_ratio(
                long_short.mean(), long_short.std(ddof=1), annualize=True
            ),
            "long_short_total_return": (
                long_short_nav.iloc[-1] - 1 if len(long_short_nav) else np.nan
            ),
            "long_short_max_drawdown": max_drawdown(long_short_nav),
        })
    summary = pd.DataFrame(summaries).sort_values(
        "rank_ic", key=lambda series: series.abs(), ascending=False
    )

    ic_output = ic_series.copy()
    ic_output["date"] = ic_output["date"].dt.strftime("%Y%m%d")
    quantile_output = quantile.copy()
    quantile_output["date"] = quantile_output["date"].dt.strftime("%Y%m%d")
    long_short_output = pivot.copy()
    long_short_output["date"] = long_short_output["date"].dt.strftime("%Y%m%d")
    ic_output.to_csv(output / "factor_ic_series.csv", index=False)
    quantile_output.to_csv(output / "factor_quantile_returns.csv", index=False)
    long_short_output.to_csv(output / "factor_long_short_returns.csv", index=False)
    quantile_summary = build_quantile_summary(quantile)
    quantile_summary.to_csv(output / "factor_quantile_summary.csv", index=False)
    correlation = compute_factor_correlation(panel)
    correlation.to_csv(output / "factor_correlation.csv")
    summary.to_csv(output / "factor_summary.csv", index=False)
    write_factor_figures(ic_series, quantile, pivot, correlation, output)
    return summary


def main() -> None:
    project_root = Path(__file__).resolve().parents[1]
    parser = argparse.ArgumentParser(description=__doc__)
    parser.add_argument("--processed", type=Path, default=project_root / "processed")
    parser.add_argument("--factors", type=Path, default=project_root / "factors")
    parser.add_argument("--output", type=Path, default=project_root / "evaluation")
    parser.add_argument(
        "--risk-exposures",
        type=Path,
        help="Optional CSV with code, industry, market_cap and optional date",
    )
    args = parser.parse_args()

    factors = load_factor_panels(args.factors.resolve())
    close = load_wide(args.processed.resolve(), "close")
    risk_path = args.risk_exposures
    if risk_path is None:
        candidate = project_root / "data" / "risk_exposures.csv"
        risk_path = candidate if candidate.exists() else None
    risk_exposures = load_risk_exposures(risk_path.resolve()) if risk_path else None
    panel = clean_factor_panel(
        factors,
        close.shift(-1) / close - 1,
        risk_exposures=risk_exposures,
    )
    output = args.output.resolve()
    output.mkdir(parents=True, exist_ok=True)
    panel_output = panel.copy()
    panel_output["date"] = panel_output["date"].dt.strftime("%Y%m%d")
    panel_output.to_csv(output / "factor_evaluation_panel.csv", index=False, quoting=1)
    neutralization_status = {
        "status": "APPLIED" if risk_exposures is not None else "SKIPPED_NO_EXPOSURES",
        "method": "daily OLS residual on industry dummies and log market cap",
        "source": str(risk_path.resolve()) if risk_path else None,
        "reason": (
            None
            if risk_exposures is not None
            else "Obfuscated stock codes and no industry/market_cap file in the assignment data"
        ),
    }
    (output / "neutralization_status.json").write_text(
        json.dumps(neutralization_status, ensure_ascii=False, indent=2) + "\n",
        encoding="utf-8",
    )
    summary = evaluate(panel, output)
    print(summary.to_string(index=False))


if __name__ == "__main__":
    main()
\end{ProjectCode}
\clearpage
\section{5. \texttt{scripts/import\_trade\_data.py}}
原始逐笔成交数据导入和基础字段汇总入口。

\textbf{物理行数：} 186\\
\textbf{SHA-256：} \texttt{5472a7ae70d2387bc021d0819abfad638d18db9737c8dc58abeaae80effa1e5f}

\begin{ProjectCode}
#!/usr/bin/env python3
"""Convert tick trades into adjusted daily/minute field tables.

Prices are adjusted as ``raw_price / 100 * adjfactor``.  Turnover amounts use
the actually traded (unadjusted) price, which is the conventional definition.
"""

from __future__ import annotations

import argparse
from pathlib import Path

import numpy as np
import pandas as pd


FIELDS = [
    "open", "high", "low", "close", "volume", "trade_count", "amount",
    "buy_volume", "sell_volume", "buy_amount", "sell_amount",
]
PRICE_FIELDS = ["open", "high", "low", "close"]
ZERO_FIELDS = [field for field in FIELDS if field not in PRICE_FIELDS]


def trading_minutes() -> list[int]:
    """Return the assignment minute grid, including the 15:00 close."""
    minutes = list(range(915, 926))
    for hour, start, end in ((9, 30, 59), (10, 0, 59), (11, 0, 30),
                             (13, 0, 59), (14, 0, 59), (15, 0, 0)):
        minutes.extend(hour * 100 + minute for minute in range(start, end + 1))
    return minutes


def load_adjustment(path: Path) -> pd.DataFrame:
    factors = pd.read_pickle(path).copy()
    factors.index = pd.to_datetime(factors.index).strftime("%Y%m%d")
    factors.columns = [str(column).split(".")[0].strip() for column in factors.columns]
    if factors.columns.duplicated().any():
        raise ValueError("Duplicate six-digit codes in adjustment-factor table")
    return factors.sort_index().ffill().bfill()


def summarize_stock(csv_path: Path, adjustment: float) -> tuple[dict[str, float], pd.DataFrame]:
    df = pd.read_csv(
        csv_path,
        dtype={"Time": "int64", "Price": "int64", "Volume": "int64", "BSFlag": "int8"},
    )
    if df.empty:
        return ({field: np.nan for field in FIELDS},
                pd.DataFrame(index=trading_minutes(), columns=FIELDS, dtype=float))

    required = {"Time", "Price", "Volume", "BSFlag"}
    if set(df.columns) != required:
        raise ValueError(f"Unexpected columns in {csv_path}: {list(df.columns)}")
    if not df["BSFlag"].isin([0, 1, 2]).all():
        raise ValueError(f"Unexpected BSFlag in {csv_path}")

    df = df.sort_values("Time", kind="stable")
    df["minute"] = df["Time"] // 100000
    raw_price = df["Price"].astype(float) / 100.0
    adjusted_price = raw_price * adjustment
    amount = raw_price * df["Volume"]
    buy = df["BSFlag"].eq(0)
    sell = df["BSFlag"].eq(1)

    daily = {
        "open": float(adjusted_price.iloc[0]),
        "high": float(adjusted_price.max()),
        "low": float(adjusted_price.min()),
        "close": float(adjusted_price.iloc[-1]),
        "volume": int(df["Volume"].sum()),
        "trade_count": int(len(df)),
        "amount": float(amount.sum()),
        "buy_volume": int(df.loc[buy, "Volume"].sum()),
        "sell_volume": int(df.loc[sell, "Volume"].sum()),
        "buy_amount": float(amount[buy].sum()),
        "sell_amount": float(amount[sell].sum()),
    }

    groups = df.groupby("minute", sort=True)
    minute = pd.DataFrame({
        "open": groups["Price"].first() / 100.0 * adjustment,
        "high": groups["Price"].max() / 100.0 * adjustment,
        "low": groups["Price"].min() / 100.0 * adjustment,
        "close": groups["Price"].last() / 100.0 * adjustment,
        "volume": groups["Volume"].sum(),
        "trade_count": groups.size(),
        "amount": amount.groupby(df["minute"]).sum(),
        "buy_volume": df["Volume"].where(buy, 0).groupby(df["minute"]).sum(),
        "sell_volume": df["Volume"].where(sell, 0).groupby(df["minute"]).sum(),
        "buy_amount": amount.where(buy, 0).groupby(df["minute"]).sum(),
        "sell_amount": amount.where(sell, 0).groupby(df["minute"]).sum(),
    })
    minute.index.name = "minute"
    return daily, minute


def write_table(table: pd.DataFrame, path: Path) -> None:
    path.parent.mkdir(parents=True, exist_ok=True)
    table.to_csv(path, index_label=table.index.name or "index")


def import_data(input_root: Path, output_root: Path, factor_path: Path,
                limit_dates: int | None = None) -> None:
    trade_root = input_root / "TRADE"
    date_dirs = sorted(path for path in trade_root.iterdir()
                       if path.is_dir() and path.name.isdigit() and len(path.name) == 8)
    if limit_dates is not None:
        date_dirs = date_dirs[:limit_dates]
    if not date_dirs:
        raise SystemExit(f"No YYYYMMDD directories under {trade_root}")

    all_codes = sorted({path.stem for day in date_dirs for path in day.glob("*.csv")})
    factors = load_adjustment(factor_path)
    missing_codes = sorted(set(all_codes) - set(factors.columns))
    missing_dates = sorted({day.name for day in date_dirs} - set(factors.index))
    if missing_codes or missing_dates:
        raise ValueError(f"Adjustment coverage missing codes={missing_codes}, dates={missing_dates}")

    minutes = trading_minutes()
    daily_rows: dict[str, list[pd.Series]] = {field: [] for field in FIELDS}
    previous_close = pd.Series(index=all_codes, dtype=float)

    for day_number, day_dir in enumerate(date_dirs, start=1):
        date = day_dir.name
        day_factor = factors.loc[date, all_codes].astype(float)
        daily_by_code = pd.DataFrame(index=all_codes, columns=FIELDS, dtype=float)
        minute_by_field = {
            field: pd.DataFrame(index=minutes, columns=all_codes, dtype=float)
            for field in FIELDS
        }

        for csv_path in sorted(day_dir.glob("*.csv")):
            code = csv_path.stem
            adjustment = float(day_factor[code])
            if not np.isfinite(adjustment) or adjustment <= 0:
                raise ValueError(f"Invalid adjustment {adjustment} for {date}/{code}")
            daily, minute = summarize_stock(csv_path, adjustment)
            for field, value in daily.items():
                daily_by_code.at[code, field] = value
            for field in FIELDS:
                minute_by_field[field].loc[minute.index, code] = minute[field].to_numpy()

        # A missing daily price is carried from the preceding adjusted close.
        for field in PRICE_FIELDS:
            daily_by_code[field] = daily_by_code[field].fillna(previous_close)
        for field in ZERO_FIELDS:
            daily_by_code[field] = daily_by_code[field].fillna(0)

        # A no-trade minute uses the preceding minute close. Leading empty
        # auction minutes use yesterday's close; only the first day may bfill.
        close_filled = minute_by_field["close"].ffill().fillna(previous_close).bfill()
        for field in PRICE_FIELDS:
            minute_by_field[field] = minute_by_field[field].fillna(close_filled)
        for field in ZERO_FIELDS:
            minute_by_field[field] = minute_by_field[field].fillna(0)

        for field in FIELDS:
            daily_rows[field].append(daily_by_code[field].rename(date))
            table = minute_by_field[field]
            table.index.name = "minute"
            write_table(table, output_root / "minute" / field / f"{date}.csv")

        previous_close = daily_by_code["close"].combine_first(previous_close)
        print(f"[{day_number}/{len(date_dirs)}] imported {date}", flush=True)

    for field, rows in daily_rows.items():
        table = pd.DataFrame(rows)
        table.index.name = "date"
        write_table(table, output_root / "daily" / f"{field}.csv")


def main() -> None:
    project_root = Path(__file__).resolve().parents[1]

    parser = argparse.ArgumentParser(description=__doc__)
    parser.add_argument("--input", type=Path, default=project_root / "data")
    parser.add_argument("--output", type=Path, default=project_root / "processed")
    parser.add_argument("--adjfactor", type=Path, default=project_root / "adjfactor.pkl")
    parser.add_argument("--limit-dates", type=int)
    args = parser.parse_args()
    import_data(args.input, args.output, args.adjfactor, args.limit_dates)


if __name__ == "__main__":
    main()
\end{ProjectCode}
\clearpage
\section{6. \texttt{scripts/adjust\_existing\_processed.py}}
已有日频和分钟频表的复权调整脚本。

\textbf{物理行数：} 116\\
\textbf{SHA-256：} \texttt{c3a2bbae21d4d5a7df311eab243b440c02a6d9cfc848f13438648992cf63a592}

\begin{ProjectCode}
#!/usr/bin/env python3
"""Rebuild adjusted outputs from verified unadjusted field tables.

This acceleration path is mathematically equivalent to the raw importer.  It
also reconstructs missing minute prices from volume, so leading empty auction
minutes use the prior daily close instead of information from the current day.
"""

from __future__ import annotations

import argparse
import shutil
from pathlib import Path

import numpy as np
import pandas as pd

from import_trade_data import PRICE_FIELDS, ZERO_FIELDS, load_adjustment


def load_wide(path: Path, index_name: str) -> pd.DataFrame:
    table = pd.read_csv(path, dtype={index_name: str})
    table[index_name] = table[index_name].astype(str).str.replace(r"\.0$", "", regex=True)
    return table.set_index(index_name)


def write_wide(table: pd.DataFrame, path: Path) -> None:
    path.parent.mkdir(parents=True, exist_ok=True)
    table.to_csv(path, index_label=table.index.name)


def adjustment_for(factors: pd.DataFrame, dates: pd.Index, codes: pd.Index) -> pd.DataFrame:
    keyed_dates = pd.Index(pd.to_datetime(dates).strftime("%Y%m%d"))
    result = factors.reindex(index=keyed_dates, columns=codes)
    result.index = dates
    if result.isna().any().any():
        raise ValueError("Adjustment factor does not cover processed panel")
    return result


def rebuild(source: Path, output: Path, factor_path: Path) -> None:
    factors = load_adjustment(factor_path)
    output.mkdir(parents=True, exist_ok=True)

    # Copy the seven non-price field trees unchanged: adjustment affects prices,
    # not shares, trade counts, or actually transacted amounts.
    for field in ZERO_FIELDS:
        src_daily = source / "daily" / f"{field}.csv"
        dst_daily = output / "daily" / f"{field}.csv"
        daily = load_wide(src_daily, "date").astype(float).fillna(0)
        daily.index.name = "date"
        write_wide(daily, dst_daily)
        src_minute = source / "minute" / field
        dst_minute = output / "minute" / field
        if dst_minute.exists():
            shutil.rmtree(dst_minute)
        shutil.copytree(src_minute, dst_minute, copy_function=shutil.copy2)
        print(f"copied {field}", flush=True)

    raw_daily: dict[str, pd.DataFrame] = {}
    for field in PRICE_FIELDS:
        table = load_wide(source / "daily" / f"{field}.csv", "date").astype(float)
        raw_daily[field] = table

    raw_daily_volume = load_wide(source / "daily" / "volume.csv", "date").astype(float)
    has_trades = raw_daily_volume.notna()
    factor_matrix = adjustment_for(factors, raw_daily["close"].index,
                                   raw_daily["close"].columns)
    adjusted_daily_close = (raw_daily["close"] * factor_matrix).where(has_trades).ffill()
    adjusted_daily: dict[str, pd.DataFrame] = {}
    for field in PRICE_FIELDS:
        adjusted = (raw_daily[field] * factor_matrix).where(has_trades).fillna(adjusted_daily_close)
        adjusted.index.name = "date"
        adjusted_daily[field] = adjusted
        write_wide(adjusted, output / "daily" / f"{field}.csv")
        print(f"adjusted daily/{field}", flush=True)

    dates = list(raw_daily["close"].index)
    codes = raw_daily["close"].columns
    previous_adjusted_close = pd.Series(index=codes, dtype=float)
    for number, date in enumerate(dates, start=1):
        volume = load_wide(source / "minute" / "volume" / f"{date}.csv", "minute")
        volume = volume.reindex(columns=codes).astype(float)
        traded = volume.gt(0)

        reconstructed: dict[str, pd.DataFrame] = {}
        factor = factors.loc[pd.to_datetime(date).strftime("%Y%m%d"), codes].astype(float)
        if not np.isfinite(factor).all():
            raise ValueError(f"Invalid adjustment factor on {date}")
        raw_close = load_wide(source / "minute" / "close" / f"{date}.csv", "minute")
        adjusted_close = raw_close.reindex(columns=codes).astype(float).where(traded).mul(factor, axis=1)
        close_filled = adjusted_close.ffill().fillna(previous_adjusted_close).bfill()
        reconstructed["close"] = close_filled
        for field in ("open", "high", "low"):
            actual = load_wide(source / "minute" / field / f"{date}.csv", "minute")
            reconstructed[field] = (actual.reindex(columns=codes).astype(float)
                                    .where(traded).mul(factor, axis=1).fillna(close_filled))
        for field, table in reconstructed.items():
            table.index.name = "minute"
            write_wide(table, output / "minute" / field / f"{date}.csv")

        previous_adjusted_close = adjusted_daily["close"].loc[date].combine_first(previous_adjusted_close)
        print(f"[{number}/{len(dates)}] repriced {date}", flush=True)


def main() -> None:
    parser = argparse.ArgumentParser(description=__doc__)
    parser.add_argument("--source", type=Path, required=True)
    parser.add_argument("--output", type=Path, default=Path("processed"))
    parser.add_argument("--adjfactor", type=Path, default=Path("adjfactor.pkl"))
    args = parser.parse_args()
    rebuild(args.source, args.output, args.adjfactor)


if __name__ == "__main__":
    main()
\end{ProjectCode}
\clearpage
\section{7. \texttt{scripts/build\_kline\_csv.py}}
生成日频K线长表CSV。

\textbf{物理行数：} 135\\
\textbf{SHA-256：} \texttt{6c7cbd968d9795b9fcc884840b7711ba07ad9d549381b0b9dd69bf30c3cfb8c9}

\begin{ProjectCode}
#!/usr/bin/env python3
"""Merge the daily field panels into one analysis-ready K-line CSV."""

from __future__ import annotations

import argparse
import csv
import json
from pathlib import Path

import numpy as np
import pandas as pd


FIELDS = (
    "open",
    "high",
    "low",
    "close",
    "volume",
    "trade_count",
    "amount",
    "buy_volume",
    "sell_volume",
    "buy_amount",
    "sell_amount",
)
INTEGER_FIELDS = ("volume", "trade_count", "buy_volume", "sell_volume")


def read_panel(path: Path) -> pd.DataFrame:
    """Read one date-by-code field panel without losing identifier zeroes."""
    table = pd.read_csv(path, dtype={"date": str})
    if "date" not in table:
        raise ValueError(f"Missing date column: {path}")
    table["date"] = table["date"].str.strip().str.replace(r"\.0$", "", regex=True)
    table = table.set_index("date")
    table.columns = pd.Index([str(code).zfill(6) for code in table.columns], name="code")
    if table.index.has_duplicates or table.columns.has_duplicates:
        raise ValueError(f"Duplicate date or code in {path}")
    return table.apply(pd.to_numeric, errors="raise")


def validate_kline(kline: pd.DataFrame) -> None:
    """Validate identifiers, finiteness, OHLC relationships, and flow totals."""
    if not kline["date"].str.fullmatch(r"\d{8}").all():
        raise ValueError("Dates must use YYYYMMDD")
    if not kline["code"].str.fullmatch(r"\d{6}").all():
        raise ValueError("Stock codes must be six digits")
    if kline.duplicated(["date", "code"]).any():
        raise ValueError("Duplicate date/code rows in K-line output")

    numeric = kline[list(FIELDS)].to_numpy(dtype=float)
    if not np.isfinite(numeric).all():
        raise ValueError("K-line output contains non-finite values")

    tolerance = 1e-10
    if not (kline["high"] + tolerance >= kline[["open", "close"]].max(axis=1)).all():
        raise ValueError("K-line output contains high < open/close")
    if not (kline["low"] - tolerance <= kline[["open", "close"]].min(axis=1)).all():
        raise ValueError("K-line output contains low > open/close")
    if not (kline["high"] + tolerance >= kline["low"]).all():
        raise ValueError("K-line output contains high < low")

    nonnegative = (
        "volume",
        "trade_count",
        "amount",
        "buy_volume",
        "sell_volume",
        "buy_amount",
        "sell_amount",
    )
    if (kline[list(nonnegative)] < -tolerance).any().any():
        raise ValueError("K-line output contains negative activity values")
    if not (kline["volume"] + tolerance >= kline["buy_volume"] + kline["sell_volume"]).all():
        raise ValueError("buy_volume + sell_volume exceeds total volume")
    if not (kline["amount"] + tolerance >= kline["buy_amount"] + kline["sell_amount"]).all():
        raise ValueError("buy_amount + sell_amount exceeds total amount")


def build_daily_kline(input_dir: Path, output_path: Path) -> dict[str, object]:
    """Build and write a date-major, code-minor daily K-line long table."""
    panels = {field: read_panel(input_dir / f"{field}.csv") for field in FIELDS}
    base = panels["close"]
    dates = pd.Index(base.index.astype(str), name="date")
    codes = pd.Index(sorted(base.columns.astype(str)), name="code")
    if not dates.is_monotonic_increasing:
        raise ValueError("Daily input dates are not sorted")

    for field, panel in panels.items():
        if not panel.index.equals(base.index) or set(panel.columns) != set(codes):
            raise ValueError(f"Daily panel alignment mismatch: {field}")
        panels[field] = panel.reindex(columns=codes)

    index = pd.MultiIndex.from_product([dates, codes], names=["date", "code"])
    kline = index.to_frame(index=False)
    for field in FIELDS:
        kline[field] = panels[field].to_numpy().reshape(-1)

    for field in INTEGER_FIELDS:
        values = kline[field].to_numpy(dtype=float)
        rounded = np.rint(values)
        if not np.allclose(values, rounded, rtol=0, atol=1e-9):
            raise ValueError(f"Expected integer-valued field: {field}")
        kline[field] = rounded.astype(np.int64)

    validate_kline(kline)
    output_path.parent.mkdir(parents=True, exist_ok=True)
    kline.to_csv(output_path, index=False, quoting=csv.QUOTE_NONNUMERIC)

    return {
        "status": "PASS",
        "output": str(output_path),
        "rows": len(kline),
        "columns": list(kline.columns),
        "date_count": len(dates),
        "code_count": len(codes),
        "first_date": dates[0],
        "last_date": dates[-1],
    }


def main() -> None:
    project_root = Path(__file__).resolve().parents[1]
    parser = argparse.ArgumentParser(description=__doc__)
    parser.add_argument("--input", type=Path, default=project_root / "processed" / "daily")
    parser.add_argument("--output", type=Path, default=project_root / "kline_data" / "daily_kline.csv")
    args = parser.parse_args()
    summary = build_daily_kline(args.input.resolve(), args.output.resolve())
    print(json.dumps(summary, ensure_ascii=False, indent=2))


if __name__ == "__main__":
    main()
\end{ProjectCode}
\clearpage
\section{8. \texttt{scripts/analyze\_factors\_backtest.py}}
因子评价、组合回测、成本、换手控制、消融和绘图主脚本。

\textbf{物理行数：} 1,627\\
\textbf{SHA-256：} \texttt{56fe1623f5bff841b0d52d34c97bb7b5e8d7a0ce72f4c520ae1a74f2a3c037c0}

\begin{ProjectCode}
#!/usr/bin/env python3
"""Build/evaluate factors and run execution-aware IC-weighted backtests."""

from __future__ import annotations

import argparse
import json
import math
import os
import sys
import tempfile
from pathlib import Path

os.environ.setdefault(
    "MPLCONFIGDIR",
    str(Path(tempfile.gettempdir()) / "quant-assignment-matplotlib"),
)

import matplotlib  # noqa: E402
import numpy as np
import pandas as pd

matplotlib.use("Agg")
import matplotlib.pyplot as plt  # noqa: E402


ROOT = Path(__file__).resolve().parents[1]
sys.path.insert(0, str(ROOT / "pipeline_code"))

from construct_factors import build_factors, load_wide  # noqa: E402
from evaluate_factors import (  # noqa: E402
    clean_factor_panel,
    evaluate,
    load_risk_exposures,
)


TRADING_DAYS = 252
INITIAL_CAPITAL = 10_000_000.0
SELL_FEE = 0.0005
BASE_SLIPPAGE_BPS = 5.0
IMPACT_COEFFICIENT = 0.01
MAX_IMPACT_RATE = 0.02
LIMIT_THRESHOLD = 0.095
TOP_N = 10
IC_LOOKBACK = 20
NEW_FACTORS = {
    "avg_trade_size_surprise_20d",
    "price_volume_pressure_10d",
    "gap_intraday_divergence_5d",
}
# New signals are deliberately shrunk because they have a shorter research
# history than the four core factors.  The scale affects only the composite
# score; single-factor IC and quantile evaluation remains fully standardized.
NEW_FACTOR_SIGNAL_SCALE = 0.10
COST_SCENARIOS = {
    "optimistic": {
        "base_slippage_bps": 0.0,
        "impact_coefficient": 0.0,
        "max_impact_rate": 0.0,
    },
    "base": {
        "base_slippage_bps": BASE_SLIPPAGE_BPS,
        "impact_coefficient": IMPACT_COEFFICIENT,
        "max_impact_rate": MAX_IMPACT_RATE,
    },
    "pessimistic": {
        "base_slippage_bps": 10.0,
        "impact_coefficient": IMPACT_COEFFICIENT * 1.5,
        "max_impact_rate": MAX_IMPACT_RATE * 1.5,
    },
}
ABLATION_STEPS = [
    ("amount_only", ["amount_mean_sd_log"]),
    (
        "amount_plus_imbalance",
        ["amount_mean_sd_log", "buy_sell_imbalance_surprise_10d"],
    ),
    (
        "amount_plus_imbalance_plus_range",
        [
            "amount_mean_sd_log",
            "buy_sell_imbalance_surprise_10d",
            "intraday_range_10d",
        ],
    ),
    (
        "all_legacy_four",
        [
            "amount_mean_sd_log",
            "buy_sell_imbalance_surprise_10d",
            "intraday_range_10d",
            "reversal_5d",
        ],
    ),
    (
        "plus_trade_size_surprise",
        [
            "amount_mean_sd_log",
            "buy_sell_imbalance_surprise_10d",
            "intraday_range_10d",
            "reversal_5d",
            "avg_trade_size_surprise_20d",
        ],
    ),
    (
        "plus_price_volume_pressure",
        [
            "amount_mean_sd_log",
            "buy_sell_imbalance_surprise_10d",
            "intraday_range_10d",
            "reversal_5d",
            "avg_trade_size_surprise_20d",
            "price_volume_pressure_10d",
        ],
    ),
    (
        "all_seven",
        [
            "amount_mean_sd_log",
            "buy_sell_imbalance_surprise_10d",
            "intraday_range_10d",
            "reversal_5d",
            "avg_trade_size_surprise_20d",
            "price_volume_pressure_10d",
            "gap_intraday_divergence_5d",
        ],
    ),
]
TURNOVER_POLICIES = {
    "daily_top10": {
        "portfolio_size": 10,
        "rebalance_interval": 1,
        "buffer_exit_rank": None,
    },
    "five_day_top10": {
        "portfolio_size": 10,
        "rebalance_interval": 5,
        "buffer_exit_rank": None,
    },
    "daily_buffer20": {
        "portfolio_size": 10,
        "rebalance_interval": 1,
        "buffer_exit_rank": 20,
    },
    "five_day_buffer20": {
        "portfolio_size": 10,
        "rebalance_interval": 5,
        "buffer_exit_rank": 20,
    },
    "five_day_buffer60_top30": {
        "portfolio_size": 30,
        "rebalance_interval": 5,
        "buffer_exit_rank": 60,
    },
}
RECOMMENDED_TURNOVER_POLICY = "five_day_buffer60_top30"


def safe_ratio(mean: float, std: float, annualize: bool = False) -> float:
    if not np.isfinite(std) or std == 0:
        return np.nan
    value = mean / std
    return value * math.sqrt(TRADING_DAYS) if annualize else value


def max_drawdown(nav: pd.Series | np.ndarray) -> float:
    values = np.asarray(nav, dtype=float)
    if not len(values):
        return np.nan
    return float(np.nanmin(values / np.maximum.accumulate(values) - 1.0))


def returns_and_market_state(
    processed: Path,
    mode: str,
) -> tuple[pd.DataFrame, pd.Series, dict[str, pd.DataFrame]]:
    """Return delayed holding-period returns and entry-day trading constraints."""
    close = load_wide(processed, "close")
    open_price = load_wide(processed, "open")
    high = load_wide(processed, "high")
    low = load_wide(processed, "low")
    volume = load_wide(processed, "volume")
    amount = load_wide(processed, "amount")

    if mode == "next_close_to_close":
        entry = close.shift(-1)
        exit_price = close.shift(-2)
    elif mode == "next_open_to_open":
        entry = open_price.shift(-1)
        exit_price = open_price.shift(-2)
    else:
        raise ValueError(mode)

    returns = exit_price / entry - 1
    trade_dates = pd.Series(close.index, index=close.index).shift(-1)
    trade_high = high.shift(-1)
    trade_low = low.shift(-1)
    trade_volume = volume.shift(-1)
    trade_amount = amount.shift(-1)
    one_price = (trade_high - trade_low).abs().le(entry.abs() * 1e-10 + 1e-10)
    entry_change = entry / close - 1
    suspended = (
        trade_volume.le(0)
        | trade_amount.le(0)
        | ~np.isfinite(entry)
    )
    limit_up = one_price & entry_change.ge(LIMIT_THRESHOLD) & ~suspended
    limit_down = one_price & entry_change.le(-LIMIT_THRESHOLD) & ~suspended
    state = {
        "amount": trade_amount,
        "suspended": suspended,
        "limit_up": limit_up,
        "limit_down": limit_down,
    }
    return returns, trade_dates, state


def compute_weights(
    panel: pd.DataFrame,
    returns: pd.DataFrame,
    strategy: str,
) -> pd.DataFrame:
    ret_long = returns.stack().rename("ret").reset_index()
    ret_long.columns = ["date", "code", "ret"]
    merged = panel[["factor", "date", "code", "factor_z"]].merge(
        ret_long, on=["date", "code"], how="left"
    )
    rows = []
    for (factor, date), day in merged.groupby(["factor", "date"], sort=True):
        valid = np.isfinite(day["factor_z"]) & np.isfinite(day["ret"])
        x = day.loc[valid, "factor_z"].to_numpy(float)
        y = day.loc[valid, "ret"].to_numpy(float)
        raw_ic = (
            np.corrcoef(x, y)[0, 1]
            if len(x) >= 20 and np.std(x) and np.std(y)
            else np.nan
        )
        rows.append({"factor": factor, "date": date, "raw_ic": raw_ic})
    weights = pd.DataFrame(rows).sort_values(["factor", "date"])
    if strategy == "leaky_same_day_ic":
        weights["weight"] = weights["raw_ic"]
    elif strategy == "historical_20d_ic":
        weights["weight"] = weights.groupby("factor")["raw_ic"].transform(
            lambda series: series.shift(1).rolling(
                IC_LOOKBACK, min_periods=IC_LOOKBACK
            ).mean()
        )
    else:
        raise ValueError(strategy)
    return weights


def factor_signal_scale(factor: str) -> float:
    """Return the conservative composite-score multiplier for one factor."""
    return NEW_FACTOR_SIGNAL_SCALE if factor in NEW_FACTORS else 1.0


def state_for_codes(
    state_panels: dict[str, pd.DataFrame],
    date: pd.Timestamp,
    codes: set[str],
) -> pd.DataFrame:
    rows = {}
    for code in sorted(codes):
        rows[code] = {
            name: (
                panel.at[date, code]
                if date in panel.index and code in panel.columns
                else np.nan
            )
            for name, panel in state_panels.items()
        }
    state = pd.DataFrame.from_dict(rows, orient="index")
    state.index.name = "code"
    state["suspended"] = state["suspended"].fillna(True).astype(bool)
    state["limit_up"] = state["limit_up"].fillna(False).astype(bool)
    state["limit_down"] = state["limit_down"].fillna(False).astype(bool)
    state["amount"] = pd.to_numeric(state["amount"], errors="coerce").fillna(0.0)
    state["buyable"] = ~state["suspended"] & ~state["limit_up"]
    state["sellable"] = ~state["suspended"] & ~state["limit_down"]
    return state


def equal_weight_benchmark_return(
    returns: pd.DataFrame,
    state_panels: dict[str, pd.DataFrame],
    date: pd.Timestamp,
) -> tuple[float, int]:
    """Return the entry-day tradable stock-pool equal-weight benchmark."""
    if date not in returns.index:
        return math.nan, 0
    day_returns = returns.loc[date]
    codes = set(day_returns.index[day_returns.notna()].astype(str))
    if not codes:
        return math.nan, 0
    state = state_for_codes(state_panels, date, codes)
    valid_codes = [
        code
        for code in sorted(codes)
        if bool(state.at[code, "buyable"])
        and np.isfinite(day_returns.get(code, np.nan))
    ]
    if not valid_codes:
        return math.nan, 0
    return float(day_returns.loc[valid_codes].mean()), len(valid_codes)


def select_buffered_codes(
    ranked_day: pd.DataFrame,
    current_weights: dict[str, float],
    portfolio_size: int = TOP_N,
    buffer_exit_rank: int | None = None,
) -> list[str]:
    """Keep held names inside the exit buffer, then fill by current rank."""
    ranked_day = ranked_day.sort_values(["rank", "code"])
    ranked_codes = ranked_day["code"].astype(str).tolist()
    if buffer_exit_rank is None:
        return ranked_codes[:portfolio_size]
    ranks = ranked_day.set_index("code")["rank"].to_dict()
    retained = sorted(
        [
            code
            for code, weight in current_weights.items()
            if weight > 0 and ranks.get(code, math.inf) <= buffer_exit_rank
        ],
        key=lambda code: (ranks[code], code),
    )[:portfolio_size]
    selected = list(retained)
    for code in ranked_codes:
        if code not in selected:
            selected.append(code)
        if len(selected) == portfolio_size:
            break
    return selected


def build_constrained_target(
    selected_codes: list[str],
    current_weights: dict[str, float],
    state: pd.DataFrame,
    portfolio_size: int = TOP_N,
) -> tuple[dict[str, float], float, dict[str, int]]:
    """Allocate after freezing positions that cannot be traded at the entry."""
    selected = list(dict.fromkeys(selected_codes))
    selected_set = set(selected)
    fixed: dict[str, float] = {}

    for code, weight in current_weights.items():
        if weight <= 0:
            continue
        sellable = bool(state.at[code, "sellable"]) if code in state.index else False
        buyable = bool(state.at[code, "buyable"]) if code in state.index else False
        if not sellable or (code in selected_set and not buyable):
            fixed[code] = float(weight)

    fixed_total = min(sum(fixed.values()), 1.0)
    if sum(fixed.values()) > 1.0 and fixed:
        scale = 1.0 / sum(fixed.values())
        fixed = {code: weight * scale for code, weight in fixed.items()}
    candidates = [
        code
        for code in selected
        if code not in fixed
        and code in state.index
        and bool(state.at[code, "buyable"])
    ]
    available = max(0.0, 1.0 - sum(fixed.values()))
    target = dict(fixed)
    if candidates:
        allocation = available / len(candidates)
        target.update({code: allocation for code in candidates})
    cash_weight = max(0.0, 1.0 - sum(target.values()))

    blocked_buys = sum(
        code not in candidates
        and current_weights.get(code, 0.0) < 1.0 / portfolio_size
        for code in selected
    )
    blocked_sells = sum(
        code not in selected_set
        and weight > 0
        and code in state.index
        and not bool(state.at[code, "sellable"])
        for code, weight in current_weights.items()
    )
    diagnostics = {
        "blocked_buys": int(blocked_buys),
        "blocked_sells": int(blocked_sells),
        "suspended_names": int(state["suspended"].sum()),
        "limit_up_names": int(state["limit_up"].sum()),
        "limit_down_names": int(state["limit_down"].sum()),
    }
    return target, cash_weight, diagnostics


def estimate_trade_costs(
    nav: float,
    current_weights: dict[str, float],
    target_weights: dict[str, float],
    state: pd.DataFrame,
    base_slippage_bps: float = BASE_SLIPPAGE_BPS,
    impact_coefficient: float = IMPACT_COEFFICIENT,
    max_impact_rate: float = MAX_IMPACT_RATE,
) -> tuple[dict[str, float], pd.DataFrame]:
    """Estimate fees, spread/slippage, and square-root participation impact."""
    rows = []
    for code in sorted(set(current_weights) | set(target_weights)):
        current = float(current_weights.get(code, 0.0))
        target = float(target_weights.get(code, 0.0))
        trade_weight = target - current
        if abs(trade_weight) < 1e-12:
            continue
        row = state.loc[code] if code in state.index else pd.Series(dtype=float)
        buyable = bool(row.get("buyable", False))
        sellable = bool(row.get("sellable", False))
        if trade_weight > 0 and not buyable:
            raise AssertionError(f"Attempted blocked buy: {code}")
        if trade_weight < 0 and not sellable:
            raise AssertionError(f"Attempted blocked sell: {code}")
        notional = abs(trade_weight) * nav
        daily_amount = max(float(row.get("amount", 0.0)), 1.0)
        participation = notional / daily_amount
        base_rate = base_slippage_bps / 10_000.0
        impact_rate = min(
            max_impact_rate,
            impact_coefficient * math.sqrt(max(participation, 0.0)),
        )
        base_cost = notional * base_rate
        impact_cost = notional * impact_rate
        sell_fee = notional * SELL_FEE if trade_weight < 0 else 0.0
        rows.append({
            "code": code,
            "current_weight": current,
            "target_weight": target,
            "trade_weight": trade_weight,
            "notional": notional,
            "daily_amount": daily_amount,
            "participation_rate": participation,
            "buyable": buyable,
            "sellable": sellable,
            "suspended": bool(row.get("suspended", True)),
            "limit_up": bool(row.get("limit_up", False)),
            "limit_down": bool(row.get("limit_down", False)),
            "base_slippage_cost": base_cost,
            "impact_cost": impact_cost,
            "sell_fee": sell_fee,
            "total_cost": base_cost + impact_cost + sell_fee,
        })
    trades = pd.DataFrame(rows)
    costs = {
        "base_slippage_cost": float(trades["base_slippage_cost"].sum()) if len(trades) else 0.0,
        "impact_cost": float(trades["impact_cost"].sum()) if len(trades) else 0.0,
        "sell_fee": float(trades["sell_fee"].sum()) if len(trades) else 0.0,
        "total_cost": float(trades["total_cost"].sum()) if len(trades) else 0.0,
        "buy_turnover": float(trades.loc[trades["trade_weight"].gt(0), "trade_weight"].sum()) if len(trades) else 0.0,
        "sell_turnover": float(-trades.loc[trades["trade_weight"].lt(0), "trade_weight"].sum()) if len(trades) else 0.0,
    }
    return costs, trades


def simulate(
    panel: pd.DataFrame,
    returns: pd.DataFrame,
    trade_dates: pd.Series,
    state_panels: dict[str, pd.DataFrame],
    strategy: str,
    mode: str,
    base_slippage_bps: float = BASE_SLIPPAGE_BPS,
    impact_coefficient: float = IMPACT_COEFFICIENT,
    max_impact_rate: float = MAX_IMPACT_RATE,
    portfolio_size: int = TOP_N,
    rebalance_interval: int = 1,
    buffer_exit_rank: int | None = None,
) -> tuple[pd.DataFrame, pd.DataFrame, pd.DataFrame, pd.DataFrame]:
    if portfolio_size < 1:
        raise ValueError("portfolio_size must be at least one")
    if rebalance_interval < 1:
        raise ValueError("rebalance_interval must be at least one")
    if buffer_exit_rank is not None and buffer_exit_rank < portfolio_size:
        raise ValueError("buffer_exit_rank cannot be smaller than the portfolio")
    weights = compute_weights(panel, returns, strategy)
    components = panel[["factor", "date", "code", "factor_z"]].merge(
        weights[["factor", "date", "weight"]],
        on=["factor", "date"],
        how="inner",
    )
    components = components[np.isfinite(components["weight"])]
    components["signal_scale"] = components["factor"].map(factor_signal_scale)
    components["component"] = (
        components["factor_z"]
        * components["weight"]
        * components["signal_scale"]
    )
    signals = components.groupby(["date", "code"], as_index=False).agg(
        signal=("component", "sum"),
        n_factor=("factor", "nunique"),
    )
    ret_long = returns.stack().rename("ret").reset_index()
    ret_long.columns = ["date", "code", "ret"]
    eligible = signals.merge(ret_long, on=["date", "code"], how="inner")
    eligible = eligible[np.isfinite(eligible["ret"])].sort_values(
        ["date", "signal", "code"],
        ascending=[True, False, True],
    )
    eligible["rank"] = eligible.groupby("date").cumcount() + 1
    maximum_rank = max(portfolio_size, buffer_exit_rank or portfolio_size)
    ranked = eligible.loc[eligible["rank"].le(maximum_rank)].copy()

    nav = INITIAL_CAPITAL
    gross_nav = INITIAL_CAPITAL
    benchmark_nav = INITIAL_CAPITAL
    active_nav = INITIAL_CAPITAL
    current_weights: dict[str, float] = {}
    results = []
    holdings = []
    trade_parts = []
    for period_number, (date, desired) in enumerate(
        ranked.groupby("date", sort=True)
    ):
        if len(desired) < maximum_rank:
            continue
        rebalanced = period_number % rebalance_interval == 0 or not current_weights
        desired_codes = select_buffered_codes(
            desired,
            current_weights,
            portfolio_size=portfolio_size,
            buffer_exit_rank=buffer_exit_rank,
        )
        relevant_codes = set(desired_codes) | set(current_weights)
        state = state_for_codes(state_panels, date, relevant_codes)
        if rebalanced:
            target_weights, cash_weight, diagnostics = build_constrained_target(
                desired_codes,
                current_weights,
                state,
                portfolio_size=portfolio_size,
            )
            costs, trades = estimate_trade_costs(
                nav,
                current_weights,
                target_weights,
                state,
                base_slippage_bps=base_slippage_bps,
                impact_coefficient=impact_coefficient,
                max_impact_rate=max_impact_rate,
            )
        else:
            target_weights = dict(current_weights)
            cash_weight = max(0.0, 1.0 - sum(target_weights.values()))
            costs = {
                "base_slippage_cost": 0.0,
                "impact_cost": 0.0,
                "sell_fee": 0.0,
                "total_cost": 0.0,
                "buy_turnover": 0.0,
                "sell_turnover": 0.0,
            }
            trades = pd.DataFrame()
            diagnostics = {
                "blocked_buys": 0,
                "blocked_sells": 0,
                "suspended_names": int(state["suspended"].sum()),
                "limit_up_names": int(state["limit_up"].sum()),
                "limit_down_names": int(state["limit_down"].sum()),
            }
        if costs["total_cost"] >= nav:
            raise RuntimeError(f"Trading costs exceed NAV on {date}")

        returns_by_code = {}
        for code in target_weights:
            value = returns.at[date, code] if code in returns.columns else np.nan
            returns_by_code[code] = float(value) if np.isfinite(value) else 0.0
        gross_return = float(sum(
            weight * returns_by_code.get(code, 0.0)
            for code, weight in target_weights.items()
        ))
        benchmark_return, benchmark_n_stock = equal_weight_benchmark_return(
            returns, state_panels, date
        )
        if not np.isfinite(benchmark_return):
            benchmark_return = 0.0
        nav_end = (nav - costs["total_cost"]) * (1 + gross_return)
        net_return = nav_end / nav - 1
        gross_nav *= 1 + gross_return
        benchmark_nav *= 1 + benchmark_return
        active_return = net_return - benchmark_return
        active_nav *= (1 + net_return) / (1 + benchmark_return)
        trade_date = trade_dates.get(date, pd.NaT)
        results.append({
            "strategy": strategy,
            "mode": mode,
            "signal_date": date,
            "trade_date": trade_date,
            "rebalanced": rebalanced,
            "portfolio_size": portfolio_size,
            "rebalance_interval": rebalance_interval,
            "buffer_exit_rank": buffer_exit_rank,
            "gross_return": gross_return,
            "benchmark_return": benchmark_return,
            "benchmark_n_stock": benchmark_n_stock,
            **costs,
            **diagnostics,
            "cash_weight": cash_weight,
            "net_return": net_return,
            "active_return": active_return,
            "nav": nav_end,
            "gross_nav": gross_nav,
            "benchmark_nav": benchmark_nav,
            "active_nav": active_nav,
            "n_holding": int(sum(weight > 0 for weight in target_weights.values())),
        })

        signal_map = desired.set_index("code")["signal"].to_dict()
        for code, target_weight in target_weights.items():
            row = state.loc[code]
            holdings.append({
                "strategy": strategy,
                "mode": mode,
                "signal_date": date,
                "trade_date": trade_date,
                "rebalanced": rebalanced,
                "portfolio_size": portfolio_size,
                "code": code,
                "signal": signal_map.get(code, np.nan),
                "target_weight": target_weight,
                "suspended": bool(row["suspended"]),
                "limit_up": bool(row["limit_up"]),
                "limit_down": bool(row["limit_down"]),
            })
        if len(trades):
            trades = trades.copy()
            trades.insert(0, "mode", mode)
            trades.insert(0, "strategy", strategy)
            trades.insert(2, "signal_date", date)
            trades.insert(3, "trade_date", trade_date)
            trades.insert(4, "rebalance_interval", rebalance_interval)
            trades.insert(5, "portfolio_size", portfolio_size)
            trades.insert(6, "buffer_exit_rank", buffer_exit_rank)
            trade_parts.append(trades)

        denominator = 1 + gross_return
        current_weights = {
            code: weight * (1 + returns_by_code.get(code, 0.0)) / denominator
            for code, weight in target_weights.items()
            if weight > 0 and denominator > 0
        }
        nav = nav_end

    return (
        pd.DataFrame(results),
        pd.DataFrame(holdings),
        weights,
        pd.concat(trade_parts, ignore_index=True) if trade_parts else pd.DataFrame(),
    )


def summarize_backtest_results(results: pd.DataFrame) -> pd.DataFrame:
    """Summarize gross, net, benchmark, active return, turnover, and costs."""
    rows = []
    for (strategy, mode), group in results.groupby(["strategy", "mode"], sort=True):
        group = group.sort_values("signal_date")
        periods = len(group)
        net_returns = group["net_return"]
        benchmark_returns = group["benchmark_return"]
        active_returns = net_returns - benchmark_returns
        final_nav = group["nav"].iloc[-1]
        total_return = final_nav / INITIAL_CAPITAL - 1
        gross_total_return = group["gross_nav"].iloc[-1] / INITIAL_CAPITAL - 1
        benchmark_total_return = group["benchmark_nav"].iloc[-1] / INITIAL_CAPITAL - 1
        active_total_return = group["active_nav"].iloc[-1] / INITIAL_CAPITAL - 1
        average_one_way_turnover = (
            group["buy_turnover"] + group["sell_turnover"]
        ).mean() / 2
        rows.append({
            "strategy": strategy,
            "mode": mode,
            "n_periods": periods,
            "gross_total_return": gross_total_return,
            "gross_annual_return": (1 + gross_total_return) ** (TRADING_DAYS / periods) - 1,
            "total_return": total_return,
            "annual_return": (1 + total_return) ** (TRADING_DAYS / periods) - 1,
            "benchmark_total_return": benchmark_total_return,
            "benchmark_annual_return": (
                (1 + benchmark_total_return) ** (TRADING_DAYS / periods) - 1
            ),
            "active_total_return": active_total_return,
            "annual_excess_return": (1 + active_total_return) ** (TRADING_DAYS / periods) - 1,
            "tracking_error": active_returns.std(ddof=1) * math.sqrt(TRADING_DAYS),
            "information_ratio": safe_ratio(
                active_returns.mean(), active_returns.std(ddof=1), annualize=True
            ),
            "cost_drag": gross_total_return - total_return,
            "annual_volatility": net_returns.std(ddof=1) * math.sqrt(TRADING_DAYS),
            "sharpe": safe_ratio(
                net_returns.mean(), net_returns.std(ddof=1), annualize=True
            ),
            "max_drawdown": max_drawdown(group["nav"]),
            "win_rate": net_returns.gt(0).mean(),
            "avg_buy_turnover": group["buy_turnover"].mean(),
            "avg_sell_turnover": group["sell_turnover"].mean(),
            "avg_one_way_turnover": average_one_way_turnover,
            "annualized_one_way_turnover": average_one_way_turnover * TRADING_DAYS,
            "avg_cash_weight": group["cash_weight"].mean(),
            "avg_blocked_buys": group["blocked_buys"].mean(),
            "avg_blocked_sells": group["blocked_sells"].mean(),
            "total_sell_fee": group["sell_fee"].sum(),
            "total_slippage_cost": group["base_slippage_cost"].sum(),
            "total_impact_cost": group["impact_cost"].sum(),
            "total_cost": group["total_cost"].sum(),
            "final_nav": final_nav,
        })
    return pd.DataFrame(rows)


def summarize_period_slice(
    group: pd.DataFrame,
    sample: str,
) -> dict[str, object]:
    """Summarize one normalized research/holdout/full result slice."""
    group = group.sort_values("signal_date").copy()
    periods = len(group)
    net_returns = group["net_return"].astype(float)
    gross_returns = group["gross_return"].astype(float)
    benchmark_returns = group["benchmark_return"].astype(float)
    active_returns = net_returns - benchmark_returns
    net_nav = (1 + net_returns).cumprod()
    gross_nav = (1 + gross_returns).cumprod()
    benchmark_nav = (1 + benchmark_returns).cumprod()
    active_nav = ((1 + net_returns) / (1 + benchmark_returns)).cumprod()
    total_return = float(net_nav.iloc[-1] - 1)
    active_total_return = float(active_nav.iloc[-1] - 1)
    normalized_nav_with_start = np.r_[1.0, net_nav.to_numpy(float)]
    average_one_way_turnover = float(
        (group["buy_turnover"] + group["sell_turnover"]).mean() / 2
    )
    return {
        "sample": sample,
        "mode": group["mode"].iloc[0],
        "turnover_policy": group["turnover_policy"].iloc[0],
        "start_date": group["signal_date"].iloc[0],
        "end_date": group["signal_date"].iloc[-1],
        "n_periods": periods,
        "portfolio_size": int(group["portfolio_size"].iloc[0]),
        "rebalance_interval": int(group["rebalance_interval"].iloc[0]),
        "buffer_exit_rank": int(group["buffer_exit_rank"].iloc[0]),
        "gross_total_return": float(gross_nav.iloc[-1] - 1),
        "total_return": total_return,
        "annual_return": (1 + total_return) ** (TRADING_DAYS / periods) - 1,
        "benchmark_total_return": float(benchmark_nav.iloc[-1] - 1),
        "active_total_return": active_total_return,
        "annual_excess_return": (
            (1 + active_total_return) ** (TRADING_DAYS / periods) - 1
        ),
        "annual_volatility": float(net_returns.std(ddof=1) * math.sqrt(TRADING_DAYS)),
        "sharpe": safe_ratio(
            net_returns.mean(), net_returns.std(ddof=1), annualize=True
        ),
        "max_drawdown": max_drawdown(normalized_nav_with_start),
        "win_rate": float(net_returns.gt(0).mean()),
        "tracking_error": float(active_returns.std(ddof=1) * math.sqrt(TRADING_DAYS)),
        "information_ratio": safe_ratio(
            active_returns.mean(), active_returns.std(ddof=1), annualize=True
        ),
        "avg_one_way_turnover": average_one_way_turnover,
        "annualized_one_way_turnover": average_one_way_turnover * TRADING_DAYS,
        "total_sell_fee": float(group["sell_fee"].sum()),
        "total_slippage_cost": float(group["base_slippage_cost"].sum()),
        "total_impact_cost": float(group["impact_cost"].sum()),
        "total_cost": float(group["total_cost"].sum()),
        "normalized_final_nav": float(net_nav.iloc[-1]),
    }


def build_optimized_strategy_metrics(
    turnover_results: pd.DataFrame,
) -> tuple[pd.DataFrame, pd.DataFrame]:
    """Return full/research80/holdout20 metrics and the final strategy path."""
    optimized = turnover_results.loc[
        turnover_results["turnover_policy"].eq(RECOMMENDED_TURNOVER_POLICY)
    ].copy()
    rows = []
    for _mode, group in optimized.groupby("mode", sort=True):
        group = group.sort_values("signal_date").reset_index(drop=True)
        split = int(len(group) * 0.8)
        for sample, sample_rows in (
            ("research80", group.iloc[:split]),
            ("holdout20", group.iloc[split:]),
            ("full", group),
        ):
            rows.append(summarize_period_slice(sample_rows, sample))
    return pd.DataFrame(rows), optimized


def composite_signal_statistics(
    panel: pd.DataFrame,
    returns: pd.DataFrame,
) -> dict[str, float]:
    """Evaluate historical-IC composite Rank IC and gross Q5-Q1 returns."""
    weights = compute_weights(panel, returns, "historical_20d_ic")
    components = panel[["factor", "date", "code", "factor_z"]].merge(
        weights[["factor", "date", "weight"]],
        on=["factor", "date"],
        how="inner",
    )
    components = components[np.isfinite(components["weight"])].copy()
    components["signal_scale"] = components["factor"].map(factor_signal_scale)
    components["component"] = (
        components["factor_z"]
        * components["weight"]
        * components["signal_scale"]
    )
    signal = components.groupby(["date", "code"], as_index=False)["component"].sum()
    signal = signal.rename(columns={"component": "signal"})
    ret_long = returns.stack().rename("ret").reset_index()
    ret_long.columns = ["date", "code", "ret"]
    merged = signal.merge(ret_long, on=["date", "code"], how="inner")
    rank_ics = []
    long_short_returns = []
    for _date, day in merged.groupby("date", sort=True):
        day = day.loc[np.isfinite(day["signal"]) & np.isfinite(day["ret"])].copy()
        if len(day) < 20 or day["signal"].nunique() < 2 or day["ret"].nunique() < 2:
            continue
        rank_ics.append(day["signal"].corr(day["ret"], method="spearman"))
        quantile = np.ceil(day["signal"].rank(method="first", pct=True) * 5).clip(1, 5)
        grouped = day.assign(quantile=quantile).groupby("quantile")["ret"].mean()
        if 1 in grouped.index and 5 in grouped.index:
            long_short_returns.append(float(grouped.loc[5] - grouped.loc[1]))
    long_short_nav = np.cumprod(1 + np.asarray(long_short_returns, dtype=float))
    return {
        "composite_rank_ic": float(np.nanmean(rank_ics)) if rank_ics else math.nan,
        "composite_long_short_total_return": (
            float(long_short_nav[-1] - 1) if len(long_short_nav) else math.nan
        ),
    }


def run_cost_sensitivity(
    panel: pd.DataFrame,
    processed: Path,
) -> pd.DataFrame:
    rows = []
    strategy_parameters = TURNOVER_POLICIES[RECOMMENDED_TURNOVER_POLICY]
    for mode in ("next_close_to_close", "next_open_to_open"):
        returns, trade_dates, state = returns_and_market_state(processed, mode)
        for scenario, parameters in COST_SCENARIOS.items():
            result, _holdings, _weights, _trades = simulate(
                panel,
                returns,
                trade_dates,
                state,
                "historical_20d_ic",
                mode,
                **strategy_parameters,
                **parameters,
            )
            summary = summarize_backtest_results(result).iloc[0].to_dict()
            rows.append({
                "scenario": scenario,
                **parameters,
                **summary,
            })
    return pd.DataFrame(rows)


def run_factor_ablation(
    panel: pd.DataFrame,
    processed: Path,
) -> pd.DataFrame:
    rows = []
    strategy_parameters = TURNOVER_POLICIES[RECOMMENDED_TURNOVER_POLICY]
    for mode in ("next_close_to_close", "next_open_to_open"):
        returns, trade_dates, state = returns_and_market_state(processed, mode)
        for step, factors in ABLATION_STEPS:
            subset = panel.loc[panel["factor"].isin(factors)].copy()
            statistics = composite_signal_statistics(subset, returns)
            result, _holdings, _weights, _trades = simulate(
                subset,
                returns,
                trade_dates,
                state,
                "historical_20d_ic",
                mode,
                **strategy_parameters,
            )
            summary = summarize_backtest_results(result).iloc[0].to_dict()
            rows.append({
                "ablation_step": step,
                "n_factors": len(factors),
                "factors": "|".join(factors),
                **statistics,
                **summary,
            })
    return pd.DataFrame(rows)


def run_turnover_control(
    panel: pd.DataFrame,
    processed: Path,
) -> tuple[pd.DataFrame, pd.DataFrame, pd.DataFrame, pd.DataFrame]:
    """Compare rebalance intervals and rank buffers under base costs."""
    summary_rows = []
    result_parts = []
    holding_parts = []
    trade_parts = []
    for mode in ("next_close_to_close", "next_open_to_open"):
        returns, trade_dates, state = returns_and_market_state(processed, mode)
        for policy, parameters in TURNOVER_POLICIES.items():
            results, holdings, _weights, trades = simulate(
                panel,
                returns,
                trade_dates,
                state,
                "historical_20d_ic",
                mode,
                **parameters,
            )
            for table in (results, holdings, trades):
                if len(table):
                    table.insert(0, "turnover_policy", policy)
            summary = summarize_backtest_results(results).iloc[0].to_dict()
            summary_rows.append({
                "turnover_policy": policy,
                "recommended": policy == RECOMMENDED_TURNOVER_POLICY,
                **parameters,
                **summary,
            })
            result_parts.append(results)
            holding_parts.append(holdings)
            if len(trades):
                trade_parts.append(trades)
    return (
        pd.DataFrame(summary_rows),
        pd.concat(result_parts, ignore_index=True),
        pd.concat(holding_parts, ignore_index=True),
        pd.concat(trade_parts, ignore_index=True),
    )


def write_backtest_figures(
    results: pd.DataFrame,
    summary: pd.DataFrame,
    sensitivity: pd.DataFrame,
    ablation: pd.DataFrame,
    turnover_summary: pd.DataFrame,
    turnover_results: pd.DataFrame,
    optimized_metrics: pd.DataFrame,
    output: Path,
) -> None:
    figures = output / "figures"
    figures.mkdir(parents=True, exist_ok=True)
    formal_results = results.loc[results["strategy"].eq("historical_20d_ic")].copy()

    figure, axes = plt.subplots(2, 1, figsize=(10, 8), sharex=False)
    for axis, (mode, group) in zip(
        axes, formal_results.groupby("mode", sort=True), strict=True
    ):
        group = group.sort_values("signal_date")
        axis.plot(group["signal_date"], group["nav"] / INITIAL_CAPITAL, label="strategy net")
        axis.plot(group["signal_date"], group["gross_nav"] / INITIAL_CAPITAL, label="strategy gross")
        axis.plot(group["signal_date"], group["benchmark_nav"] / INITIAL_CAPITAL, label="equal-weight benchmark")
        axis.plot(group["signal_date"], group["active_nav"] / INITIAL_CAPITAL, label="active NAV")
        axis.axhline(1, color="0.25", linewidth=0.8)
        axis.set_title(mode)
        axis.set_ylabel("NAV")
        axis.legend(ncol=2, fontsize=8)
        axis.tick_params(axis="x", labelrotation=30, labelsize=8)
    figure.tight_layout()
    figure.savefig(figures / "strategy_benchmark_excess.png", dpi=160)
    plt.close(figure)

    formal_summary = summary.loc[summary["strategy"].eq("historical_20d_ic")].copy()
    positions = np.arange(len(formal_summary))
    figure, axes = plt.subplots(1, 2, figsize=(12, 4.6))
    bottoms = np.zeros(len(formal_summary))
    for column, label in (
        ("total_sell_fee", "sell fee"),
        ("total_slippage_cost", "base slippage"),
        ("total_impact_cost", "market impact"),
    ):
        values = formal_summary[column].to_numpy() / 1_000_000
        axes[0].bar(positions, values, bottom=bottoms, label=label)
        bottoms += values
    axes[0].set_xticks(positions, formal_summary["mode"], rotation=20, ha="right")
    axes[0].set_ylabel("Cost (million CNY)")
    axes[0].set_title("Transaction-cost decomposition")
    axes[0].legend(fontsize=8)
    width = 0.25
    for offset, column, label in (
        (-width, "gross_total_return", "gross"),
        (0, "benchmark_total_return", "benchmark"),
        (width, "total_return", "net"),
    ):
        axes[1].bar(positions + offset, formal_summary[column] * 100, width, label=label)
    axes[1].axhline(0, color="0.25", linewidth=0.8)
    axes[1].set_xticks(positions, formal_summary["mode"], rotation=20, ha="right")
    axes[1].set_ylabel("Total return (%)")
    axes[1].set_title("Gross, benchmark, and net return")
    axes[1].legend(fontsize=8)
    figure.tight_layout()
    figure.savefig(figures / "cost_and_return_breakdown.png", dpi=160)
    plt.close(figure)

    scenario_order = list(COST_SCENARIOS)
    mode_order = ["next_close_to_close", "next_open_to_open"]
    figure, axis = plt.subplots(figsize=(9, 4.8))
    width = 0.24
    positions = np.arange(len(mode_order))
    for index, scenario in enumerate(scenario_order):
        values = [
            sensitivity.loc[
                sensitivity["scenario"].eq(scenario) & sensitivity["mode"].eq(mode),
                "total_return",
            ].iloc[0] * 100
            for mode in mode_order
        ]
        axis.bar(positions + (index - 1) * width, values, width, label=scenario)
    axis.axhline(0, color="0.25", linewidth=0.8)
    axis.set_xticks(positions, mode_order, rotation=15, ha="right")
    axis.set_ylabel("Net total return (%)")
    axis.set_title("Cost sensitivity")
    axis.legend()
    figure.tight_layout()
    figure.savefig(figures / "cost_sensitivity.png", dpi=160)
    plt.close(figure)

    figure, axis = plt.subplots(figsize=(9.5, 4.8))
    ablation_order = [step for step, _factors in ABLATION_STEPS]
    positions = np.arange(len(ablation_order))
    for index, mode in enumerate(mode_order):
        values = [
            ablation.loc[
                ablation["ablation_step"].eq(step) & ablation["mode"].eq(mode),
                "total_return",
            ].iloc[0] * 100
            for step in ablation_order
        ]
        axis.bar(positions + (index - 0.5) * 0.36, values, 0.36, label=mode)
    axis.axhline(0, color="0.25", linewidth=0.8)
    axis.set_xticks(positions, ablation_order, rotation=20, ha="right")
    axis.set_ylabel("Net total return (%)")
    axis.set_title("Sequential factor ablation")
    axis.legend(fontsize=8)
    figure.tight_layout()
    figure.savefig(figures / "factor_ablation.png", dpi=160)
    plt.close(figure)

    diagnostic = summary.copy()
    figure, axis = plt.subplots(figsize=(8.5, 4.8))
    positions = np.arange(len(mode_order))
    for index, strategy in enumerate(("historical_20d_ic", "leaky_same_day_ic")):
        values = [
            diagnostic.loc[
                diagnostic["strategy"].eq(strategy) & diagnostic["mode"].eq(mode),
                "total_return",
            ].iloc[0] * 100
            for mode in mode_order
        ]
        axis.bar(positions + (index - 0.5) * 0.36, values, 0.36, label=strategy)
    axis.axhline(0, color="0.25", linewidth=0.8)
    axis.set_xticks(positions, mode_order, rotation=15, ha="right")
    axis.set_ylabel("Net total return (%)")
    axis.set_title("Look-ahead bias diagnostic")
    axis.legend(fontsize=8)
    figure.tight_layout()
    figure.savefig(figures / "lookahead_diagnostic.png", dpi=160)
    plt.close(figure)

    policy_order = list(TURNOVER_POLICIES)
    figure, axes = plt.subplots(1, 2, figsize=(12.5, 4.8))
    positions = np.arange(len(policy_order))
    width = 0.36
    for index, mode in enumerate(mode_order):
        policy_rows = turnover_summary.loc[
            turnover_summary["mode"].eq(mode)
        ].set_index("turnover_policy").loc[policy_order]
        axes[0].bar(
            positions + (index - 0.5) * width,
            policy_rows["avg_one_way_turnover"] * 100,
            width,
            label=mode,
        )
        axes[1].bar(
            positions + (index - 0.5) * width,
            policy_rows["total_return"] * 100,
            width,
            label=mode,
        )
    for axis in axes:
        axis.set_xticks(positions, policy_order, rotation=20, ha="right")
        axis.axhline(0, color="0.25", linewidth=0.8)
        axis.legend(fontsize=8)
    axes[0].set_title("Average one-way turnover")
    axes[0].set_ylabel("Turnover per day (%)")
    axes[1].set_title("Net total return after costs")
    axes[1].set_ylabel("Total return (%)")
    figure.tight_layout()
    figure.savefig(figures / "turnover_control_comparison.png", dpi=160)
    plt.close(figure)

    recommended = turnover_results.loc[
        turnover_results["turnover_policy"].isin(
            ["daily_top10", RECOMMENDED_TURNOVER_POLICY]
        )
    ].copy()
    figure, axes = plt.subplots(2, 1, figsize=(10, 8))
    for axis, mode in zip(axes, mode_order, strict=True):
        mode_rows = recommended.loc[recommended["mode"].eq(mode)]
        for policy, group in mode_rows.groupby("turnover_policy", sort=True):
            group = group.sort_values("signal_date")
            axis.plot(
                group["signal_date"],
                group["nav"] / INITIAL_CAPITAL,
                label=policy,
            )
        benchmark = mode_rows.loc[
            mode_rows["turnover_policy"].eq(RECOMMENDED_TURNOVER_POLICY)
        ].sort_values("signal_date")
        axis.plot(
            benchmark["signal_date"],
            benchmark["benchmark_nav"] / INITIAL_CAPITAL,
            label="equal-weight benchmark",
        )
        axis.axhline(1, color="0.25", linewidth=0.8)
        axis.set_title(mode)
        axis.set_ylabel("NAV")
        axis.legend(fontsize=8)
        axis.tick_params(axis="x", labelrotation=30, labelsize=8)
    figure.tight_layout()
    figure.savefig(figures / "turnover_control_nav.png", dpi=160)
    plt.close(figure)

    optimized = turnover_results.loc[
        turnover_results["turnover_policy"].eq(RECOMMENDED_TURNOVER_POLICY)
    ].copy()
    figure, axes = plt.subplots(2, 2, figsize=(12, 8), sharex="col")
    for row_index, mode in enumerate(mode_order):
        group = optimized.loc[optimized["mode"].eq(mode)].sort_values("signal_date")
        net_nav = group["nav"] / INITIAL_CAPITAL
        benchmark_nav = group["benchmark_nav"] / INITIAL_CAPITAL
        drawdown = net_nav / net_nav.cummax() - 1
        axes[row_index, 0].plot(group["signal_date"], net_nav, label="optimized strategy")
        axes[row_index, 0].plot(
            group["signal_date"], benchmark_nav, label="equal-weight benchmark"
        )
        axes[row_index, 0].axhline(1, color="0.25", linewidth=0.8)
        axes[row_index, 0].set_title(f"{mode}: NAV")
        axes[row_index, 0].set_ylabel("NAV")
        axes[row_index, 0].legend(fontsize=8)
        axes[row_index, 1].fill_between(
            group["signal_date"], drawdown * 100, 0, color="#B64B4B", alpha=0.75
        )
        axes[row_index, 1].set_title(f"{mode}: drawdown")
        axes[row_index, 1].set_ylabel("Drawdown (%)")
        for axis in axes[row_index]:
            axis.tick_params(axis="x", labelrotation=30, labelsize=8)
    figure.tight_layout()
    figure.savefig(figures / "optimized_strategy_nav_drawdown.png", dpi=170)
    plt.close(figure)

    optimized_summary = turnover_summary.loc[
        turnover_summary["turnover_policy"].eq(RECOMMENDED_TURNOVER_POLICY)
    ].set_index("mode").loc[mode_order]
    positions = np.arange(len(mode_order))
    figure, axes = plt.subplots(1, 3, figsize=(15.0, 4.8))
    width = 0.26
    for offset, column, label in (
        (-width, "total_return", "total return"),
        (0, "annual_return", "annual return"),
        (width, "annual_excess_return", "annual excess"),
    ):
        axes[0].bar(
            positions + offset,
            optimized_summary[column].to_numpy(float) * 100,
            width,
            label=label,
        )
    axes[0].axhline(25, color="#B64B4B", linestyle="--", linewidth=1.0, label="25% target")
    axes[0].set_xticks(positions, mode_order, rotation=15, ha="right")
    axes[0].set_ylabel("Return (%)")
    axes[0].set_title("Optimized strategy return metrics")
    axes[0].legend(fontsize=8)
    axes[1].bar(
        positions - width / 2,
        optimized_summary["annual_volatility"].to_numpy(float) * 100,
        width,
        label="Annual volatility (%)",
    )
    axes[1].bar(
        positions + width / 2,
        -optimized_summary["max_drawdown"].to_numpy(float) * 100,
        width,
        label="Max drawdown (%)",
    )
    axes[1].set_xticks(positions, mode_order, rotation=15, ha="right")
    axes[1].set_ylabel("Risk (%)")
    axes[1].set_title("Volatility and drawdown")
    axes[1].legend(fontsize=8)
    axes[2].bar(
        positions - width / 2,
        optimized_summary["sharpe"].to_numpy(float),
        width,
        label="Sharpe",
    )
    axes[2].bar(
        positions + width / 2,
        optimized_summary["information_ratio"].to_numpy(float),
        width,
        label="Information ratio",
    )
    axes[2].axhline(0, color="0.25", linewidth=0.8)
    axes[2].set_xticks(positions, mode_order, rotation=15, ha="right")
    axes[2].set_ylabel("Ratio")
    axes[2].set_title("Risk-adjusted ratios")
    axes[2].legend(fontsize=8)
    figure.tight_layout()
    figure.savefig(figures / "optimized_strategy_risk_metrics.png", dpi=170)
    plt.close(figure)

    period_metrics = optimized_metrics.loc[
        optimized_metrics["mode"].eq("next_open_to_open")
        & optimized_metrics["sample"].isin(["research80", "holdout20"])
    ].set_index("sample").loc[["research80", "holdout20"]]
    positions = np.arange(len(period_metrics))
    figure, axis = plt.subplots(figsize=(8.4, 4.8))
    for offset, column, label in (
        (-width, "total_return", "strategy net return"),
        (0, "benchmark_total_return", "benchmark return"),
        (width, "active_total_return", "active return"),
    ):
        axis.bar(
            positions + offset,
            period_metrics[column].to_numpy(float) * 100,
            width,
            label=label,
        )
    axis.axhline(0, color="0.25", linewidth=0.8)
    axis.set_xticks(positions, ["Research 80%", "Holdout 20%"])
    axis.set_ylabel("Cumulative return (%)")
    axis.set_title("Optimized next-open strategy: research and holdout periods")
    axis.legend(fontsize=8)
    figure.tight_layout()
    figure.savefig(figures / "optimized_period_split.png", dpi=170)
    plt.close(figure)


def run_backtests(
    panel: pd.DataFrame,
    processed: Path,
    output: Path,
) -> tuple[pd.DataFrame, pd.DataFrame, pd.DataFrame, pd.DataFrame]:
    output.mkdir(parents=True, exist_ok=True)
    results_parts = []
    holding_parts = []
    weight_parts = []
    trade_parts = []
    for mode in ("next_close_to_close", "next_open_to_open"):
        returns, trade_dates, state = returns_and_market_state(processed, mode)
        for strategy in ("leaky_same_day_ic", "historical_20d_ic"):
            results, holdings, weights, trades = simulate(
                panel, returns, trade_dates, state, strategy, mode
            )
            results_parts.append(results)
            holding_parts.append(holdings)
            weights = weights.copy()
            weights["strategy"] = strategy
            weights["mode"] = mode
            weight_parts.append(weights)
            trade_parts.append(trades)

    results = pd.concat(results_parts, ignore_index=True)
    holdings = pd.concat(holding_parts, ignore_index=True)
    weights = pd.concat(weight_parts, ignore_index=True)
    trades = pd.concat(trade_parts, ignore_index=True)
    summary = summarize_backtest_results(results)
    sensitivity = run_cost_sensitivity(panel, processed)
    ablation = run_factor_ablation(panel, processed)
    (
        turnover_summary,
        turnover_results,
        turnover_holdings,
        turnover_trades,
    ) = run_turnover_control(panel, processed)
    optimized_metrics, optimized_results = build_optimized_strategy_metrics(
        turnover_results
    )
    lookahead = summary.copy()
    lookahead.insert(
        2,
        "validity",
        np.where(
            lookahead["strategy"].eq("historical_20d_ic"),
            "formal_tradable_signal",
            "invalid_lookahead_diagnostic",
        ),
    )
    lookahead.insert(
        3,
        "is_formal_result",
        lookahead["strategy"].eq("historical_20d_ic"),
    )
    benchmark_excess = results.loc[
        results["strategy"].eq("historical_20d_ic"),
        [
            "mode",
            "signal_date",
            "trade_date",
            "gross_return",
            "net_return",
            "benchmark_return",
            "active_return",
            "gross_nav",
            "nav",
            "benchmark_nav",
            "active_nav",
            "benchmark_n_stock",
        ],
    ].copy()
    write_backtest_figures(
        results,
        summary,
        sensitivity,
        ablation,
        turnover_summary,
        turnover_results,
        optimized_metrics,
        output,
    )

    for table in (
        results,
        holdings,
        weights,
        trades,
        benchmark_excess,
        turnover_results,
        turnover_holdings,
        turnover_trades,
        optimized_metrics,
        optimized_results,
    ):
        for column in (
            "signal_date",
            "trade_date",
            "date",
            "start_date",
            "end_date",
        ):
            if column in table:
                table[column] = pd.to_datetime(table[column]).dt.strftime("%Y%m%d")
    results.to_csv(output / "backtest_results.csv", index=False)
    holdings.to_csv(output / "backtest_holdings.csv", index=False, quoting=1)
    weights.to_csv(output / "factor_ic_weights.csv", index=False)
    trades.to_csv(output / "backtest_trades.csv", index=False, quoting=1)
    summary.to_csv(output / "backtest_summary.csv", index=False)
    benchmark_excess.to_csv(output / "benchmark_excess_returns.csv", index=False)
    sensitivity.to_csv(output / "cost_sensitivity.csv", index=False)
    ablation.to_csv(output / "factor_ablation.csv", index=False)
    lookahead.to_csv(output / "lookahead_comparison.csv", index=False)
    turnover_summary.to_csv(output / "turnover_control_comparison.csv", index=False)
    turnover_results.to_csv(output / "turnover_control_results.csv", index=False)
    turnover_holdings.to_csv(
        output / "turnover_control_holdings.csv", index=False, quoting=1
    )
    turnover_trades.to_csv(
        output / "turnover_control_trades.csv", index=False, quoting=1
    )
    optimized_metrics.to_csv(
        output / "optimized_strategy_metrics.csv", index=False
    )
    optimized_results.to_csv(
        output / "optimized_strategy_results.csv", index=False
    )
    assumptions = {
        "signal_execution_delay": "signal at t, trade at t+1 open or close",
        "sell_fee": SELL_FEE,
        "base_slippage_bps_each_side": BASE_SLIPPAGE_BPS,
        "impact_model": "min(max_rate, coefficient * sqrt(order_notional / daily_amount))",
        "impact_coefficient": IMPACT_COEFFICIENT,
        "max_impact_rate": MAX_IMPACT_RATE,
        "benchmark": (
            "daily equal-weight return of stocks buyable at the entry; "
            "no benchmark transaction costs"
        ),
        "cost_scenarios": COST_SCENARIOS,
        "turnover_policies": TURNOVER_POLICIES,
        "recommended_turnover_policy": RECOMMENDED_TURNOVER_POLICY,
        "new_factors": sorted(NEW_FACTORS),
        "new_factor_signal_scale": NEW_FACTOR_SIGNAL_SCALE,
        "suspension_rule": "entry-day volume <= 0 or amount <= 0 blocks trading",
        "limit_rule": (
            "one-price day and entry/previous-close move >= 9.5%; "
            "limit-up blocks buys, limit-down blocks sells"
        ),
    }
    (output / "execution_assumptions.json").write_text(
        json.dumps(assumptions, ensure_ascii=False, indent=2) + "\n",
        encoding="utf-8",
    )
    return summary, sensitivity, ablation, turnover_summary


def write_analysis_report(
    factor_summary: pd.DataFrame,
    backtest_summary: pd.DataFrame,
    sensitivity: pd.DataFrame,
    ablation: pd.DataFrame,
    turnover_summary: pd.DataFrame,
    neutralization_status: dict,
    output: Path,
) -> None:
    def pct(value: float) -> str:
        return f"{value:.2%}" if np.isfinite(value) else "NA"

    lines = [
        "# 因子评价与执行约束回测结果",
        "",
        "## 口径",
        "",
        "- 因子目标为下一交易日收盘到收盘收益；每日横截面先做 1%/99% 缩尾和 z-score。",
        (
            "- 行业/市值中性化状态："
            f"`{neutralization_status['status']}`。"
            + (
                "已按日对行业哑变量和对数市值回归取残差。"
                if neutralization_status["status"] == "APPLIED"
                else "题目数据缺少真实行业和市值，代码不会用混淆代码伪造暴露。"
            )
        ),
        "- 信号在 t 日形成，只允许在 t+1 开盘或收盘成交，不再使用同一收盘价成交。",
        "- 停牌不可交易；一字涨停不可买入；一字跌停不可卖出。",
        "- 成本包含卖出万五、双边 5bp 基础滑点和基于成交额参与率的平方根冲击成本。",
        (
            "- 新增三个因子在综合信号中的乘数为 "
            f"{NEW_FACTOR_SIGNAL_SCALE:.2f}，用于抑制探索性因子对成熟因子的过度替代。"
        ),
        "",
        "## 因子汇总",
        "",
        "| 因子 | IC | ICIR | Rank IC | Rank ICIR | Q5-Q1总收益 |",
        "|---|---:|---:|---:|---:|---:|",
    ]
    for row in factor_summary.itertuples(index=False):
        lines.append(
            f"| {row.factor} | {row.ic:.4f} | {row.icir:.3f} | "
            f"{row.rank_ic:.4f} | {row.rank_icir:.3f} | "
            f"{pct(row.long_short_total_return)} |"
        )
    lines += [
        "",
        "## 正式组合回测、基准与超额收益",
        "",
        "基准为股票池内入场日可以买入股票的每日等权收益，不扣除基准交易成本。",
        "",
        "| 成交口径 | 毛收益 | 基准收益 | 净收益 | 年化超额 | 信息比率 | 日均单边换手 | 年化单边换手 |",
        "|---|---:|---:|---:|---:|---:|---:|---:|",
    ]
    formal = backtest_summary.loc[
        backtest_summary["strategy"].eq("historical_20d_ic")
    ]
    for row in formal.itertuples(index=False):
        lines.append(
            f"| {row.mode} | {pct(row.gross_total_return)} | "
            f"{pct(row.benchmark_total_return)} | {pct(row.total_return)} | "
            f"{pct(row.annual_excess_return)} | {row.information_ratio:.3f} | "
            f"{pct(row.avg_one_way_turnover)} | "
            f"{row.annualized_one_way_turnover:.1f}x |"
        )
    lines += [
        "",
        "## 交易成本拆分",
        "",
        "| 成交口径 | 手续费 | 基础滑点 | 冲击成本 | 总成本 | 成本拖累 |",
        "|---|---:|---:|---:|---:|---:|",
    ]
    for row in formal.itertuples(index=False):
        lines.append(
            f"| {row.mode} | {row.total_sell_fee:,.0f} | "
            f"{row.total_slippage_cost:,.0f} | {row.total_impact_cost:,.0f} | "
            f"{row.total_cost:,.0f} | {pct(row.cost_drag)} |"
        )
    lines += [
        "",
        "## 成本敏感性",
        "",
        "| 情景 | 成交口径 | 净收益 | 夏普 | 总成本 |",
        "|---|---|---:|---:|---:|",
    ]
    for row in sensitivity.itertuples(index=False):
        lines.append(
            f"| {row.scenario} | {row.mode} | {pct(row.total_return)} | "
            f"{row.sharpe:.3f} | {row.total_cost:,.0f} |"
        )
    lines += [
        "",
        "## 换手控制",
        "",
        "| 政策 | 口径 | 持仓数 | 调仓间隔 | 退出缓冲 | 日均单边换手 | 净收益 | 年化收益 | 最大回撤 | 夏普 |",
        "|---|---|---:|---:|---:|---:|---:|---:|---:|---:|",
    ]
    for row in turnover_summary.itertuples(index=False):
        buffer_label = int(row.buffer_exit_rank) if np.isfinite(row.buffer_exit_rank) else "无"
        lines.append(
            f"| {row.turnover_policy} | {row.mode} | {int(row.portfolio_size)} | "
            f"{row.rebalance_interval}日 | "
            f"{buffer_label} | {pct(row.avg_one_way_turnover)} | "
            f"{pct(row.total_return)} | {pct(row.annual_return)} | "
            f"{pct(row.max_drawdown)} | {row.sharpe:.3f} |"
        )
    recommended_rows = turnover_summary.loc[
        turnover_summary["turnover_policy"].eq(RECOMMENDED_TURNOVER_POLICY)
    ]
    lines += [
        "",
        f"推荐政策为 `{RECOMMENDED_TURNOVER_POLICY}`：每 5 个交易日调仓，"
        "已持有股票只要仍位于前 60 名就继续保留，再从当前排名中补足 30 只。",
        "",
        "## 最终优化策略完整指标",
        "",
        "| 成交口径 | 成本前收益 | 成本后收益 | 年化收益 | 基准收益 | 累计主动收益 | 年化超额 | 年化波动 | 最大回撤 | 夏普 | 信息比率 | 总成本 |",
        "|---|---:|---:|---:|---:|---:|---:|---:|---:|---:|---:|---:|",
    ]
    for row in recommended_rows.itertuples(index=False):
        lines.append(
            f"| {row.mode} | {pct(row.gross_total_return)} | "
            f"{pct(row.total_return)} | {pct(row.annual_return)} | "
            f"{pct(row.benchmark_total_return)} | {pct(row.active_total_return)} | "
            f"{pct(row.annual_excess_return)} | {pct(row.annual_volatility)} | "
            f"{pct(row.max_drawdown)} | {row.sharpe:.3f} | "
            f"{row.information_ratio:.3f} | {row.total_cost:,.0f} |"
        )
    lines += [
        "",
        "## 因子逐步组合消融",
        "",
        "| 组合 | 成交口径 | 因子数 | 综合Rank IC | 多空收益 | 策略净收益 |",
        "|---|---|---:|---:|---:|---:|",
    ]
    for row in ablation.itertuples(index=False):
        lines.append(
            f"| {row.ablation_step} | {row.mode} | {row.n_factors} | "
            f"{row.composite_rank_ic:.4f} | "
            f"{pct(row.composite_long_short_total_return)} | "
            f"{pct(row.total_return)} |"
        )
    lines += [
        "",
        "## 前视偏差诊断",
        "",
        "`leaky_same_day_ic` 使用未来收益，只用于数据泄漏诊断；"
        "正式统计和超额收益结论使用 `historical_20d_ic`。",
        "",
        "图表保存在 `evaluation/figures/` 与 `backtest/figures/`。",
    ]
    output.mkdir(parents=True, exist_ok=True)
    (output / "因子评价与回测.md").write_text(
        "\n".join(lines) + "\n", encoding="utf-8"
    )


def main() -> None:
    parser = argparse.ArgumentParser(description=__doc__)
    parser.add_argument("--processed", type=Path, default=ROOT / "processed")
    parser.add_argument("--factors", type=Path, default=ROOT / "factors")
    parser.add_argument("--evaluation", type=Path, default=ROOT / "evaluation")
    parser.add_argument("--backtest", type=Path, default=ROOT / "backtest")
    parser.add_argument("--reports", type=Path, default=ROOT / "reports")
    parser.add_argument(
        "--risk-exposures",
        type=Path,
        help="Optional true code/industry/market_cap exposure CSV",
    )
    args = parser.parse_args()

    factors = build_factors(args.processed.resolve(), args.factors.resolve())
    close = load_wide(args.processed.resolve(), "close")
    risk_path = args.risk_exposures
    if risk_path is None:
        candidate = ROOT / "data" / "risk_exposures.csv"
        risk_path = candidate if candidate.exists() else None
    risk_exposures = load_risk_exposures(risk_path.resolve()) if risk_path else None
    panel = clean_factor_panel(
        factors,
        close.shift(-1) / close - 1,
        risk_exposures=risk_exposures,
    )
    args.evaluation.mkdir(parents=True, exist_ok=True)
    panel_out = panel.copy()
    panel_out["date"] = panel_out["date"].dt.strftime("%Y%m%d")
    panel_out.to_csv(
        args.evaluation / "factor_evaluation_panel.csv",
        index=False,
        quoting=1,
    )
    neutralization_status = {
        "status": "APPLIED" if risk_exposures is not None else "SKIPPED_NO_EXPOSURES",
        "method": "daily OLS residual on industry dummies and log market cap",
        "source": str(risk_path.resolve()) if risk_path else None,
        "reason": (
            None
            if risk_exposures is not None
            else "Obfuscated stock codes and no industry/market_cap file in assignment data"
        ),
    }
    (args.evaluation / "neutralization_status.json").write_text(
        json.dumps(neutralization_status, ensure_ascii=False, indent=2) + "\n",
        encoding="utf-8",
    )
    factor_summary = evaluate(panel, args.evaluation.resolve())
    backtest_summary, sensitivity, ablation, turnover_summary = run_backtests(
        panel, args.processed.resolve(), args.backtest.resolve()
    )
    write_analysis_report(
        factor_summary,
        backtest_summary,
        sensitivity,
        ablation,
        turnover_summary,
        neutralization_status,
        args.reports.resolve(),
    )
    print(factor_summary.to_string(index=False), flush=True)
    print(backtest_summary.to_string(index=False), flush=True)


if __name__ == "__main__":
    main()
\end{ProjectCode}
\clearpage
\section{9. \texttt{scripts/train\_lstm\_pytorch.py}}
多股票Walk-forward PyTorch LSTM及逻辑回归基线。

\textbf{物理行数：} 885\\
\textbf{SHA-256：} \texttt{3b20d753b654a03810ea37f57c593f6a311a562228abb8a647ea4cd6db6af292}

\begin{ProjectCode}
#!/usr/bin/env python3
"""Train a multi-stock PyTorch LSTM with strict walk-forward evaluation."""

from __future__ import annotations

import argparse
import json
import math
import os
import random
import tempfile
from pathlib import Path

os.environ.setdefault(
    "MPLCONFIGDIR",
    str(Path(tempfile.gettempdir()) / "quant-assignment-matplotlib"),
)

import matplotlib  # noqa: E402
import numpy as np
import pandas as pd
import torch
from sklearn.linear_model import LogisticRegression
from sklearn.metrics import (
    accuracy_score,
    average_precision_score,
    balanced_accuracy_score,
    confusion_matrix,
    f1_score,
    precision_score,
    recall_score,
    roc_auc_score,
)
from torch import nn
from torch.utils.data import DataLoader, TensorDataset

matplotlib.use("Agg")
import matplotlib.pyplot as plt  # noqa: E402


FEATURES = [
    "log_return",
    "range",
    "log_volume",
    "buy_sell_imbalance",
    "time_position",
]


class TorchLSTM(nn.Module):
    """A single-layer LSTM followed by a binary classification head."""

    def __init__(self, input_size: int, hidden_size: int) -> None:
        super().__init__()
        self.input_size = input_size
        self.hidden_size = hidden_size
        self.lstm = nn.LSTM(
            input_size=input_size,
            hidden_size=hidden_size,
            batch_first=True,
        )
        self.classifier = nn.Linear(hidden_size, 1)

    def forward(self, x: torch.Tensor) -> torch.Tensor:
        sequence, _state = self.lstm(x)
        return self.classifier(sequence[:, -1]).squeeze(-1)


def seed_everything(seed: int) -> None:
    random.seed(seed)
    np.random.seed(seed)
    torch.manual_seed(seed)
    if torch.cuda.is_available():
        torch.cuda.manual_seed_all(seed)


def resolve_device(requested: str) -> torch.device:
    if requested == "auto":
        if torch.backends.mps.is_available():
            return torch.device("mps")
        if torch.cuda.is_available():
            return torch.device("cuda")
        return torch.device("cpu")
    if requested == "mps" and not torch.backends.mps.is_available():
        raise RuntimeError("MPS was requested but is not available")
    if requested == "cuda" and not torch.cuda.is_available():
        raise RuntimeError("CUDA was requested but is not available")
    return torch.device(requested)


def read_field(path: Path, codes: list[str]) -> pd.DataFrame:
    table = pd.read_csv(path, usecols=["minute", *codes])
    table["minute"] = table["minute"].astype(int)
    return table.set_index("minute")[codes].astype(float)


def select_codes(
    first_close_file: Path,
    requested_codes: list[str] | None,
    max_codes: int,
) -> list[str]:
    available = pd.read_csv(first_close_file, nrows=0).columns.tolist()[1:]
    available = [str(code).zfill(6) for code in available]
    if requested_codes:
        codes = [str(code).zfill(6) for code in requested_codes]
        missing = sorted(set(codes) - set(available))
        if missing:
            raise ValueError(f"Codes not present in minute tables: {missing}")
    else:
        codes = available[:max_codes]
    if len(codes) < 2:
        raise ValueError("Walk-forward LSTM requires at least two stocks")
    return codes


def make_dataset(
    processed: Path,
    requested_codes: list[str] | None,
    max_codes: int,
    max_dates: int,
    sequence_length: int,
) -> tuple[np.ndarray, np.ndarray, pd.DataFrame, list[str]]:
    close_files = sorted((processed / "minute" / "close").glob("*.csv"))
    if max_dates > 0:
        close_files = close_files[:max_dates]
    if not close_files:
        raise FileNotFoundError("No minute close tables")
    codes = select_codes(close_files[0], requested_codes, max_codes)

    x_parts: list[np.ndarray] = []
    y_parts: list[np.ndarray] = []
    metadata_parts: list[pd.DataFrame] = []
    for number, close_file in enumerate(close_files, start=1):
        date = close_file.stem
        fields = {
            "close": read_field(close_file, codes),
            "high": read_field(
                processed / "minute" / "high" / f"{date}.csv", codes
            ),
            "low": read_field(
                processed / "minute" / "low" / f"{date}.csv", codes
            ),
            "volume": read_field(
                processed / "minute" / "volume" / f"{date}.csv", codes
            ),
            "buy": read_field(
                processed / "minute" / "buy_volume" / f"{date}.csv", codes
            ),
            "sell": read_field(
                processed / "minute" / "sell_volume" / f"{date}.csv", codes
            ),
        }
        minutes = fields["close"].index.to_numpy(int)
        for code in codes:
            close = fields["close"][code]
            features = pd.DataFrame(index=close.index)
            features["log_return"] = np.log(close.clip(lower=1e-12)).diff().fillna(0)
            features["range"] = (
                fields["high"][code] - fields["low"][code]
            ) / close.clip(lower=1e-12)
            features["log_volume"] = np.log1p(fields["volume"][code])
            features["buy_sell_imbalance"] = (
                fields["buy"][code] - fields["sell"][code]
            ) / fields["volume"][code].clip(lower=1)
            features["time_position"] = np.linspace(0, 1, len(features))
            values = (
                features[FEATURES]
                .replace([np.inf, -np.inf], np.nan)
                .fillna(0)
                .to_numpy(float)
            )
            prices = close.to_numpy(float)
            ends = np.arange(sequence_length - 1, len(values) - 1)
            if not len(ends):
                continue
            sequences = np.stack([
                values[end - sequence_length + 1:end + 1]
                for end in ends
            ])
            labels = (prices[ends + 1] > prices[ends]).astype(np.float32)
            x_parts.append(sequences.astype(np.float32))
            y_parts.append(labels)
            metadata_parts.append(pd.DataFrame({
                "date": date,
                "code": code,
                "minute": minutes[ends + 1],
                "current_close": prices[ends],
                "next_close": prices[ends + 1],
            }))
        if number % 20 == 0 or number == len(close_files):
            print(
                f"prepared {number}/{len(close_files)} dates for {len(codes)} stocks",
                flush=True,
            )
    if not x_parts:
        raise ValueError("No LSTM samples were created")
    return (
        np.concatenate(x_parts).astype(np.float32),
        np.concatenate(y_parts).astype(np.float32),
        pd.concat(metadata_parts, ignore_index=True),
        codes,
    )


def make_walk_forward_folds(
    dates: list[str],
    min_train_dates: int,
    validation_dates: int,
    test_dates: int,
    step_dates: int,
) -> list[dict[str, object]]:
    unique_dates = sorted(dict.fromkeys(str(date) for date in dates))
    folds = []
    train_end = min_train_dates
    fold_number = 1
    while train_end + validation_dates + test_dates <= len(unique_dates):
        validation_end = train_end + validation_dates
        test_end = validation_end + test_dates
        folds.append({
            "fold": fold_number,
            "train_dates": unique_dates[:train_end],
            "validation_dates": unique_dates[train_end:validation_end],
            "test_dates": unique_dates[validation_end:test_end],
        })
        fold_number += 1
        train_end += step_dates
    if len(folds) < 2:
        raise ValueError(
            "Walk-forward configuration creates fewer than two test folds; "
            "increase --max-dates or reduce fold windows"
        )
    return folds


def predict_proba(
    model: TorchLSTM,
    x: np.ndarray,
    device: torch.device,
    batch_size: int,
) -> np.ndarray:
    loader = DataLoader(
        TensorDataset(torch.from_numpy(x)),
        batch_size=batch_size,
        shuffle=False,
    )
    pieces = []
    model.eval()
    with torch.inference_mode():
        for (batch_x,) in loader:
            probability = torch.sigmoid(model(batch_x.to(device)))
            pieces.append(probability.cpu().numpy())
    return np.concatenate(pieces)


def metrics_from_probability(
    y: np.ndarray,
    probability: np.ndarray,
) -> dict[str, float]:
    prediction = probability >= 0.5
    clipped = probability.clip(1e-8, 1 - 1e-8)
    loss = float(np.mean(
        -(y * np.log(clipped) + (1 - y) * np.log(1 - clipped))
    ))
    auc = (
        roc_auc_score(y, probability)
        if len(np.unique(y)) > 1
        else math.nan
    )
    pr_auc = (
        average_precision_score(y, probability)
        if len(np.unique(y)) > 1
        else math.nan
    )
    tn, fp, fn, tp = confusion_matrix(
        y.astype(int), prediction.astype(int), labels=[0, 1]
    ).ravel()
    return {
        "positive_rate": float(y.mean()),
        "majority_baseline_accuracy": max(float(y.mean()), 1 - float(y.mean())),
        "loss": loss,
        "accuracy": float(accuracy_score(y, prediction)),
        "balanced_accuracy": float(balanced_accuracy_score(y, prediction)),
        "precision": float(precision_score(y, prediction, zero_division=0)),
        "recall": float(recall_score(y, prediction, zero_division=0)),
        "f1": float(f1_score(y, prediction, zero_division=0)),
        "auc": float(auc),
        "pr_auc": float(pr_auc),
        "tn": int(tn),
        "fp": int(fp),
        "fn": int(fn),
        "tp": int(tp),
    }


def sequence_summary_features(x: np.ndarray) -> np.ndarray:
    """Create a simple 15-column sequence summary for logistic regression."""
    if x.ndim != 3:
        raise ValueError("Expected [sample, time, feature] sequence input")
    return np.concatenate([x[:, -1], x.mean(axis=1), x.std(axis=1)], axis=1)


def write_lstm_figures(
    comparison: pd.DataFrame,
    lstm_fold_metrics: pd.DataFrame,
    baseline_fold_metrics: pd.DataFrame,
    output: Path,
) -> None:
    figures = output / "figures"
    figures.mkdir(parents=True, exist_ok=True)
    metrics = ["balanced_accuracy", "auc", "pr_auc", "f1"]
    positions = np.arange(len(metrics))
    width = 0.24
    figure, axis = plt.subplots(figsize=(9.5, 4.8))
    for index, row in enumerate(comparison.itertuples(index=False)):
        values = [getattr(row, metric) for metric in metrics]
        axis.bar(positions + (index - 1) * width, values, width, label=row.model)
    axis.set_xticks(positions, metrics, rotation=15)
    axis.set_ylim(0, 1)
    axis.set_title("Walk-forward out-of-sample model comparison")
    axis.legend(fontsize=8)
    figure.tight_layout()
    figure.savefig(figures / "model_comparison.png", dpi=160)
    plt.close(figure)

    # Accuracy needs a dedicated diagnostic because the up/down labels are
    # imbalanced.  Overall accuracy alone can look high even when a model
    # mostly predicts the majority (down) class, so show balanced accuracy and
    # the fold-level majority baseline beside it.
    model_labels = {
        "majority_class": "Majority class",
        "logistic_regression": "Logistic regression",
        "lstm": "LSTM",
    }
    ordered_models = ["majority_class", "logistic_regression", "lstm"]
    overall_accuracy = comparison.set_index("model").loc[ordered_models]
    test_lstm = lstm_fold_metrics.loc[
        lstm_fold_metrics["split"].eq("test")
    ].copy()
    test_logistic = baseline_fold_metrics.loc[
        baseline_fold_metrics["split"].eq("test")
        & baseline_fold_metrics["model"].eq("logistic_regression")
    ].copy()
    majority_by_fold = test_lstm[["fold", "majority_baseline_accuracy"]].rename(
        columns={"majority_baseline_accuracy": "accuracy"}
    )
    majority_by_fold["model"] = "majority_class"
    fold_accuracy = pd.concat(
        [
            majority_by_fold[["fold", "model", "accuracy"]],
            test_logistic[["fold", "model", "accuracy"]],
            test_lstm.assign(model="lstm")[["fold", "model", "accuracy"]],
        ],
        ignore_index=True,
    )

    figure, axes = plt.subplots(1, 2, figsize=(10.5, 4.6))
    overall_metrics = ["accuracy", "balanced_accuracy"]
    positions = np.arange(len(overall_metrics))
    width = 0.24
    for index, model in enumerate(ordered_models):
        row = overall_accuracy.loc[model]
        values = [float(row[metric]) for metric in overall_metrics]
        bars = axes[0].bar(
            positions + (index - 1) * width,
            values,
            width,
            label=model_labels[model],
        )
        axes[0].bar_label(bars, labels=[f"{value:.2%}" for value in values], fontsize=8)
    axes[0].set_xticks(positions, ["Accuracy", "Balanced accuracy"])
    axes[0].set_ylim(0, 0.80)
    axes[0].set_ylabel("Score")
    axes[0].set_title("Combined out-of-sample accuracy")
    axes[0].legend(fontsize=8, loc="lower left")

    for model in ordered_models:
        group = fold_accuracy.loc[fold_accuracy["model"].eq(model)].sort_values("fold")
        axes[1].plot(
            group["fold"],
            group["accuracy"],
            marker="o",
            linewidth=1.8,
            label=model_labels[model],
        )
        # The three curves are separated by only a few basis points.  Label
        # the LSTM curve only so the exact values remain readable instead of
        # stacking nine nearly identical annotations on top of one another.
        if model == "lstm":
            for fold, value in zip(group["fold"], group["accuracy"], strict=True):
                axes[1].annotate(
                    f"LSTM {value:.2%}",
                    (fold, value),
                    xytext=(0, 7),
                    textcoords="offset points",
                    ha="center",
                    fontsize=8,
                    color="tab:green",
                )
    axes[1].set_xticks(sorted(fold_accuracy["fold"].unique()))
    axes[1].set_ylim(0.66, 0.73)
    axes[1].set_xlabel("Walk-forward test fold")
    axes[1].set_ylabel("Accuracy (zoomed scale)")
    axes[1].set_title("Test accuracy by fold")
    axes[1].legend(fontsize=8, loc="upper right")
    figure.suptitle("LSTM accuracy diagnostic", fontsize=12)
    figure.tight_layout()
    figure.savefig(figures / "lstm_accuracy_diagnostics.png", dpi=180)
    plt.close(figure)

    lstm_row = comparison.loc[comparison["model"].eq("lstm")].iloc[0]
    matrix = np.array([[lstm_row.tn, lstm_row.fp], [lstm_row.fn, lstm_row.tp]], dtype=int)
    figure, axis = plt.subplots(figsize=(5.4, 4.8))
    image = axis.imshow(matrix, cmap="Blues")
    axis.set_xticks([0, 1], ["pred down", "pred up"])
    axis.set_yticks([0, 1], ["actual down", "actual up"])
    for row in range(2):
        for column in range(2):
            axis.text(column, row, f"{matrix[row, column]:,}", ha="center", va="center")
    axis.set_title("LSTM out-of-sample confusion matrix")
    figure.colorbar(image, ax=axis, shrink=0.82)
    figure.tight_layout()
    figure.savefig(figures / "lstm_confusion_matrix.png", dpi=160)
    plt.close(figure)

    fold_table = pd.concat([
        lstm_fold_metrics.loc[lstm_fold_metrics["split"].eq("test")].assign(model="lstm"),
        baseline_fold_metrics.loc[baseline_fold_metrics["split"].eq("test")],
    ], ignore_index=True)
    figure, axes = plt.subplots(1, 2, figsize=(10, 4.3), sharex=True)
    for axis, metric in zip(axes, ("auc", "pr_auc"), strict=True):
        for model, group in fold_table.groupby("model", sort=True):
            axis.plot(group["fold"], group[metric], marker="o", label=model)
        axis.set_title(f"Test {metric} by fold")
        axis.set_xlabel("Fold")
        axis.set_ylim(0, 1)
        axis.set_xticks(sorted(fold_table["fold"].unique()))
    axes[0].legend(fontsize=8)
    figure.tight_layout()
    figure.savefig(figures / "fold_model_metrics.png", dpi=160)
    plt.close(figure)


def train_fold(
    fold: dict[str, object],
    x: np.ndarray,
    y: np.ndarray,
    metadata: pd.DataFrame,
    hidden_size: int,
    epochs: int,
    learning_rate: float,
    batch_size: int,
    seed: int,
    device: torch.device,
) -> tuple[
    list[dict[str, float]],
    list[dict[str, float]],
    list[dict[str, float]],
    pd.DataFrame,
    dict[str, object],
]:
    fold_number = int(fold["fold"])
    date_values = metadata["date"].astype(str)
    masks = {
        "train": date_values.isin(fold["train_dates"]).to_numpy(),
        "validation": date_values.isin(fold["validation_dates"]).to_numpy(),
        "test": date_values.isin(fold["test_dates"]).to_numpy(),
    }
    split_arrays = {
        split: (x[mask], y[mask], metadata.loc[mask].reset_index(drop=True))
        for split, mask in masks.items()
    }
    x_train, y_train, _ = split_arrays["train"]
    mean = x_train.reshape(-1, x.shape[-1]).mean(axis=0)
    std = x_train.reshape(-1, x.shape[-1]).std(axis=0)
    std[std < 1e-10] = 1
    standardized = {
        split: (
            ((split_x - mean) / std).astype(np.float32),
            split_y,
            split_meta,
        )
        for split, (split_x, split_y, split_meta) in split_arrays.items()
    }

    seed_everything(seed + fold_number)
    model = TorchLSTM(x.shape[-1], hidden_size).to(device)
    optimizer = torch.optim.Adam(model.parameters(), lr=learning_rate)
    criterion = nn.BCEWithLogitsLoss()
    generator = torch.Generator().manual_seed(seed + fold_number)
    train_loader = DataLoader(
        TensorDataset(
            torch.from_numpy(standardized["train"][0]),
            torch.from_numpy(standardized["train"][1]),
        ),
        batch_size=batch_size,
        shuffle=True,
        generator=generator,
    )

    history = []
    best_val_loss = float("inf")
    best_epoch = 0
    best_state: dict[str, torch.Tensor] | None = None
    print(
        f"fold={fold_number} device={device.type} "
        f"train={len(standardized['train'][0])} "
        f"validation={len(standardized['validation'][0])} "
        f"test={len(standardized['test'][0])}",
        flush=True,
    )
    for epoch in range(1, epochs + 1):
        model.train()
        batch_losses = []
        for batch_x, batch_y in train_loader:
            batch_x = batch_x.to(device)
            batch_y = batch_y.to(device)
            optimizer.zero_grad(set_to_none=True)
            loss = criterion(model(batch_x), batch_y)
            loss.backward()
            nn.utils.clip_grad_norm_(model.parameters(), max_norm=5.0)
            optimizer.step()
            batch_losses.append(float(loss.detach().cpu()))
        train_probability = predict_proba(
            model, standardized["train"][0], device, batch_size
        )
        validation_probability = predict_proba(
            model, standardized["validation"][0], device, batch_size
        )
        train_metrics = metrics_from_probability(
            standardized["train"][1], train_probability
        )
        validation_metrics = metrics_from_probability(
            standardized["validation"][1], validation_probability
        )
        history.append({
            "fold": fold_number,
            "epoch": epoch,
            "batch_loss": float(np.mean(batch_losses)),
            "train_loss": train_metrics["loss"],
            "train_accuracy": train_metrics["accuracy"],
            "val_loss": validation_metrics["loss"],
            "val_accuracy": validation_metrics["accuracy"],
            "val_auc": validation_metrics["auc"],
        })
        if validation_metrics["loss"] < best_val_loss:
            best_val_loss = validation_metrics["loss"]
            best_epoch = epoch
            best_state = {
                key: value.detach().cpu().clone()
                for key, value in model.state_dict().items()
            }
        print(
            f"fold {fold_number} epoch {epoch:02d}: "
            f"train_loss={train_metrics['loss']:.4f} "
            f"val_loss={validation_metrics['loss']:.4f} "
            f"val_auc={validation_metrics['auc']:.4f}",
            flush=True,
        )
    if best_state is None:
        raise RuntimeError(f"Fold {fold_number} produced no checkpoint")
    model.load_state_dict(best_state)

    metrics_rows = []
    probabilities = {}
    for split, (split_x, split_y, _split_meta) in standardized.items():
        probability = predict_proba(model, split_x, device, batch_size)
        probabilities[split] = probability
        metrics_rows.append({
            "fold": fold_number,
            "model": "lstm",
            "split": split,
            "n_samples": len(split_y),
            **metrics_from_probability(split_y, probability),
        })

    logistic = LogisticRegression(
        max_iter=500,
        solver="lbfgs",
        random_state=seed + fold_number,
    )
    logistic.fit(
        sequence_summary_features(standardized["train"][0]),
        standardized["train"][1].astype(int),
    )
    baseline_rows = []
    logistic_probabilities = {}
    for split, (split_x, split_y, _split_meta) in standardized.items():
        probability = logistic.predict_proba(sequence_summary_features(split_x))[:, 1]
        logistic_probabilities[split] = probability
        baseline_rows.append({
            "fold": fold_number,
            "model": "logistic_regression",
            "split": split,
            "n_samples": len(split_y),
            **metrics_from_probability(split_y, probability),
        })
    test_x, test_y, test_meta = standardized["test"]
    predictions = test_meta.copy()
    predictions.insert(0, "fold", fold_number)
    predictions["label_up"] = test_y.astype(int)
    predictions["probability_up"] = probabilities["test"]
    predictions["prediction_up"] = (
        probabilities["test"] >= 0.5
    ).astype(int)
    predictions["logistic_probability_up"] = logistic_probabilities["test"]
    predictions["logistic_prediction_up"] = (
        logistic_probabilities["test"] >= 0.5
    ).astype(int)

    checkpoint = {
        "model_type": "torch.nn.LSTM",
        "model_state_dict": best_state,
        "input_size": len(FEATURES),
        "hidden_size": hidden_size,
        "feature_mean": torch.from_numpy(mean.copy()),
        "feature_std": torch.from_numpy(std.copy()),
        "features": FEATURES,
        "sequence_length": int(x.shape[1]),
        "fold": fold_number,
        "best_epoch": best_epoch,
        "train_start": str(fold["train_dates"][0]),
        "train_end": str(fold["train_dates"][-1]),
        "validation_start": str(fold["validation_dates"][0]),
        "validation_end": str(fold["validation_dates"][-1]),
        "test_start": str(fold["test_dates"][0]),
        "test_end": str(fold["test_dates"][-1]),
        "device_used": device.type,
    }
    return history, metrics_rows, baseline_rows, predictions, checkpoint


def train_walk_forward(
    processed: Path,
    output: Path,
    requested_codes: list[str] | None,
    max_codes: int,
    max_dates: int,
    sequence_length: int,
    hidden_size: int,
    epochs: int,
    learning_rate: float,
    batch_size: int,
    seed: int,
    requested_device: str,
    min_train_dates: int,
    validation_dates: int,
    test_dates: int,
    step_dates: int,
) -> pd.DataFrame:
    seed_everything(seed)
    device = resolve_device(requested_device)
    x, y, metadata, codes = make_dataset(
        processed,
        requested_codes,
        max_codes,
        max_dates,
        sequence_length,
    )
    unique_dates = sorted(metadata["date"].astype(str).unique().tolist())
    folds = make_walk_forward_folds(
        unique_dates,
        min_train_dates,
        validation_dates,
        test_dates,
        step_dates,
    )
    output.mkdir(parents=True, exist_ok=True)
    fold_dir = output / "walk_forward"
    fold_dir.mkdir(parents=True, exist_ok=True)

    history_parts = []
    metric_rows = []
    baseline_metric_rows = []
    prediction_parts = []
    checkpoints = []
    for fold in folds:
        history, metrics, baseline_metrics, predictions, checkpoint = train_fold(
            fold,
            x,
            y,
            metadata,
            hidden_size,
            epochs,
            learning_rate,
            batch_size,
            seed,
            device,
        )
        history_parts.extend(history)
        metric_rows.extend(metrics)
        baseline_metric_rows.extend(baseline_metrics)
        prediction_parts.append(predictions)
        checkpoints.append(checkpoint)
        torch.save(
            checkpoint,
            fold_dir / f"fold_{int(fold['fold']):02d}.pt",
        )

    history_table = pd.DataFrame(history_parts)
    fold_metrics = pd.DataFrame(metric_rows)
    baseline_fold_metrics = pd.DataFrame(baseline_metric_rows)
    predictions = pd.concat(prediction_parts, ignore_index=True)
    history_table.to_csv(output / "training_history.csv", index=False)
    fold_metrics.to_csv(output / "walk_forward_metrics.csv", index=False)
    baseline_fold_metrics.to_csv(
        output / "baseline_walk_forward_metrics.csv", index=False
    )
    predictions.to_csv(output / "test_predictions.csv", index=False)

    overall = metrics_from_probability(
        predictions["label_up"].to_numpy(float),
        predictions["probability_up"].to_numpy(float),
    )
    overall_table = pd.DataFrame([{
        "split": "walk_forward_oos",
        "n_folds": len(folds),
        "n_codes": len(codes),
        "n_dates": len(unique_dates),
        "n_samples": len(predictions),
        **overall,
    }])
    overall_table.to_csv(output / "metrics.csv", index=False)

    labels = predictions["label_up"].to_numpy(float)
    majority_probability = np.full(len(labels), labels.mean(), dtype=float)
    comparison = pd.DataFrame([
        {
            "model": "majority_class",
            "n_samples": len(labels),
            **metrics_from_probability(labels, majority_probability),
        },
        {
            "model": "logistic_regression",
            "n_samples": len(labels),
            **metrics_from_probability(
                labels,
                predictions["logistic_probability_up"].to_numpy(float),
            ),
        },
        {
            "model": "lstm",
            "n_samples": len(labels),
            **overall,
        },
    ])
    comparison.to_csv(output / "model_comparison.csv", index=False)
    comparison[["model", "tn", "fp", "fn", "tp"]].to_csv(
        output / "confusion_matrix.csv", index=False
    )
    write_lstm_figures(comparison, fold_metrics, baseline_fold_metrics, output)

    final_checkpoint = dict(checkpoints[-1])
    final_checkpoint.update({
        "walk_forward": True,
        "codes": codes,
        "n_dates": len(unique_dates),
        "fold_count": len(folds),
    })
    torch.save(final_checkpoint, output / "model.pt")
    metadata_doc = {
        key: value
        for key, value in final_checkpoint.items()
        if key not in {"model_state_dict", "feature_mean", "feature_std"}
    }
    metadata_doc["feature_mean"] = final_checkpoint["feature_mean"].numpy().tolist()
    metadata_doc["feature_std"] = final_checkpoint["feature_std"].numpy().tolist()
    metadata_doc["date_start"] = unique_dates[0]
    metadata_doc["date_end"] = unique_dates[-1]
    metadata_doc["folds"] = [
        {
            "fold": int(fold["fold"]),
            "train_start": fold["train_dates"][0],
            "train_end": fold["train_dates"][-1],
            "validation_start": fold["validation_dates"][0],
            "validation_end": fold["validation_dates"][-1],
            "test_start": fold["test_dates"][0],
            "test_end": fold["test_dates"][-1],
        }
        for fold in folds
    ]
    (output / "model_metadata.json").write_text(
        json.dumps(metadata_doc, ensure_ascii=False, indent=2) + "\n",
        encoding="utf-8",
    )

    result = overall_table.iloc[0]
    logistic_result = comparison.loc[
        comparison["model"].eq("logistic_regression")
    ].iloc[0]
    report = f"""# 多股票分钟线 LSTM Walk-forward 拓展

- 股票数：{len(codes)}（{', '.join(codes)}）
- 日期范围：{unique_dates[0]} 至 {unique_dates[-1]}，共 {len(unique_dates)} 个交易日
- 输入窗口：过去 {sequence_length} 分钟
- 特征：{', '.join(FEATURES)}
- 模型：PyTorch ``torch.nn.LSTM``（hidden={hidden_size}）+ 线性二分类头
- 验证方式：{len(folds)} 个 expanding-window walk-forward 折
- 每折：至少 {min_train_dates} 个训练日、{validation_dates} 个验证日、{test_dates} 个不重叠测试日
- 标准化和模型训练均在每折内重新估计，不使用测试期信息

## 合并样本外测试

- 样本数：{int(result['n_samples'])}
- 上涨占比：{result['positive_rate']:.2%}
- 多数类基线准确率：{result['majority_baseline_accuracy']:.2%}
- LSTM 准确率：{result['accuracy']:.2%}
- 平衡准确率：{result['balanced_accuracy']:.2%}
- Precision：{result['precision']:.2%}
- Recall：{result['recall']:.2%}
- F1：{result['f1']:.2%}
- ROC AUC：{result['auc']:.4f}
- PR AUC：{result['pr_auc']:.4f}
- 混淆矩阵：TN={int(result['tn'])}, FP={int(result['fp'])}, FN={int(result['fn'])}, TP={int(result['tp'])}

## 简单模型基线

- 逻辑回归准确率：{logistic_result['accuracy']:.2%}
- 逻辑回归平衡准确率：{logistic_result['balanced_accuracy']:.2%}
- 逻辑回归 ROC AUC：{logistic_result['auc']:.4f}
- 逻辑回归 PR AUC：{logistic_result['pr_auc']:.4f}
- 逻辑回归 F1：{logistic_result['f1']:.2%}

这是多股票、长时间和严格滚动样本外结果。是否可交易仍需结合阈值、换手、滑点和延迟进行策略级验证。
"""
    (output / "LSTM实验报告.md").write_text(
        report.replace("`", "\x60"),
        encoding="utf-8",
    )
    return overall_table


def main() -> None:
    parser = argparse.ArgumentParser(description=__doc__)
    parser.add_argument("--processed", type=Path, default=Path("processed"))
    parser.add_argument("--output", type=Path, default=Path("lstm"))
    parser.add_argument(
        "--codes",
        help="Comma-separated stock codes; defaults to the first --max-codes columns",
    )
    parser.add_argument("--max-codes", type=int, default=6)
    parser.add_argument("--max-dates", type=int, default=140)
    parser.add_argument("--sequence-length", type=int, default=20)
    parser.add_argument("--hidden-size", type=int, default=16)
    parser.add_argument("--epochs", type=int, default=3)
    parser.add_argument("--learning-rate", type=float, default=0.003)
    parser.add_argument("--batch-size", type=int, default=512)
    parser.add_argument("--seed", type=int, default=42)
    parser.add_argument(
        "--device",
        choices=["auto", "cpu", "mps", "cuda"],
        default="auto",
    )
    parser.add_argument("--min-train-dates", type=int, default=60)
    parser.add_argument("--validation-dates", type=int, default=20)
    parser.add_argument("--test-dates", type=int, default=20)
    parser.add_argument("--step-dates", type=int, default=20)
    args = parser.parse_args()
    requested_codes = (
        [code.strip() for code in args.codes.split(",") if code.strip()]
        if args.codes
        else None
    )
    metrics = train_walk_forward(
        args.processed,
        args.output,
        requested_codes,
        args.max_codes,
        args.max_dates,
        args.sequence_length,
        args.hidden_size,
        args.epochs,
        args.learning_rate,
        args.batch_size,
        args.seed,
        args.device,
        args.min_train_dates,
        args.validation_dates,
        args.test_dates,
        args.step_dates,
    )
    print(metrics.to_string(index=False))


if __name__ == "__main__":
    main()
\end{ProjectCode}
\clearpage
\section{10. \texttt{scripts/validate\_outputs.py}}
对数据、因子、回测、图表和LSTM产物进行完整审计。

\textbf{物理行数：} 409\\
\textbf{SHA-256：} \texttt{afbe8f8f27ee8e9ed2fb730d5c3570b4f0991aebab5e1fc6e8ccb8bc2293a828}

\begin{ProjectCode}
#!/usr/bin/env python3
"""Validate generated tables and analysis artifacts; write a compact audit."""

from __future__ import annotations

import argparse
import json
from pathlib import Path

import numpy as np
import pandas as pd
import torch

from import_trade_data import FIELDS, PRICE_FIELDS, trading_minutes


def read_wide(path: Path, index: str) -> pd.DataFrame:
    table = pd.read_csv(path)
    return table.set_index(index).astype(float)


def validate(root: Path) -> dict:
    processed = root / "processed"
    daily = {field: read_wide(processed / "daily" / f"{field}.csv", "date")
             for field in FIELDS}
    shapes = {table.shape for table in daily.values()}
    if shapes != {(302, 300)}:
        raise AssertionError(f"Unexpected daily shapes: {shapes}")
    base_index, base_columns = daily["close"].index, daily["close"].columns
    for field, table in daily.items():
        if not table.index.equals(base_index) or not table.columns.equals(base_columns):
            raise AssertionError(f"Daily alignment mismatch: {field}")
        if not np.isfinite(table.to_numpy()).all():
            raise AssertionError(f"Non-finite daily values: {field}")
    if not (daily["high"] + 1e-10 >= daily["open"]).all().all():
        raise AssertionError("daily high < open")
    if not (daily["high"] + 1e-10 >= daily["close"]).all().all():
        raise AssertionError("daily high < close")
    if not (daily["low"] - 1e-10 <= daily["open"]).all().all():
        raise AssertionError("daily low > open")
    if not (daily["low"] - 1e-10 <= daily["close"]).all().all():
        raise AssertionError("daily low > close")
    for position in range(1, len(base_index)):
        zero_volume = daily["volume"].iloc[position].eq(0)
        expected = daily["close"].iloc[position - 1]
        for field in PRICE_FIELDS:
            if not np.allclose(daily[field].iloc[position].where(zero_volume).to_numpy(),
                               expected.where(zero_volume).to_numpy(), equal_nan=True):
                raise AssertionError(f"Daily prior-close fill mismatch: {field} at {base_index[position]}")

    kline_path = root / "kline_data" / "daily_kline.csv"
    kline = pd.read_csv(kline_path, dtype={"date": str, "code": str})
    expected_kline_columns = ["date", "code", *FIELDS]
    if kline.columns.tolist() != expected_kline_columns:
        raise AssertionError(f"Unexpected K-line columns: {kline.columns.tolist()}")
    if kline.shape != (302 * 300, len(expected_kline_columns)):
        raise AssertionError(f"Unexpected K-line shape: {kline.shape}")
    if kline[["date", "code"]].duplicated().any():
        raise AssertionError("Duplicate date/code rows in K-line CSV")
    if not kline["code"].str.fullmatch(r"\d{6}").all():
        raise AssertionError("K-line codes are not six-digit strings")
    if not (kline["high"] + 1e-10 >= kline[["open", "close"]].max(axis=1)).all():
        raise AssertionError("K-line high < open/close")
    if not (kline["low"] - 1e-10 <= kline[["open", "close"]].min(axis=1)).all():
        raise AssertionError("K-line low > open/close")

    expected_minutes = trading_minutes()
    minute_file_count = 0
    no_trade_checks = 0
    for date_position, date in enumerate(base_index.astype(int).astype(str)):
        tables = {}
        for field in FIELDS:
            path = processed / "minute" / field / f"{date}.csv"
            if not path.exists():
                raise AssertionError(f"Missing {path}")
            table = read_wide(path, "minute")
            minute_file_count += 1
            if table.shape != (253, 300) or table.index.astype(int).tolist() != expected_minutes:
                raise AssertionError(f"Bad minute table: {path} {table.shape}")
            if not np.isfinite(table.to_numpy()).all():
                raise AssertionError(f"Non-finite minute values: {path}")
            tables[field] = table
        if not (tables["high"] + 1e-10 >= tables["open"]).all().all():
            raise AssertionError(f"minute high < open on {date}")
        if not (tables["low"] - 1e-10 <= tables["close"]).all().all():
            raise AssertionError(f"minute low > close on {date}")
        no_trade = tables["volume"].eq(0)
        actual_close = tables["close"].where(~no_trade)
        expected_close = actual_close.ffill()
        if date_position:
            expected_close = expected_close.fillna(daily["close"].iloc[date_position - 1])
        expected_close = expected_close.bfill()
        if not np.allclose(tables["close"].to_numpy(), expected_close.to_numpy(), equal_nan=True):
            raise AssertionError(f"Minute prior-close fill mismatch on {date}")
        for field in PRICE_FIELDS[:-1]:
            if not np.allclose(tables[field].where(no_trade).to_numpy(),
                               tables["close"].where(no_trade).to_numpy(), equal_nan=True):
                raise AssertionError(f"No-trade OHLC fill mismatch {field} on {date}")
        no_trade_checks += int(no_trade.sum().sum())

    factor_summary = pd.read_csv(root / "evaluation" / "factor_summary.csv")
    backtest_summary = pd.read_csv(root / "backtest" / "backtest_summary.csv")
    lstm_metrics = pd.read_csv(root / "lstm" / "metrics.csv")
    walk_forward_metrics = pd.read_csv(root / "lstm" / "walk_forward_metrics.csv")
    if len(factor_summary) != 7 or len(backtest_summary) != 4 or len(lstm_metrics) != 1:
        raise AssertionError("Analysis output row count mismatch")
    expected_modes = {"next_close_to_close", "next_open_to_open"}
    if set(backtest_summary["mode"]) != expected_modes:
        raise AssertionError(f"Unexpected execution modes: {set(backtest_summary['mode'])}")
    required_cost_columns = {
        "total_sell_fee", "total_slippage_cost", "total_impact_cost", "total_cost"
    }
    if not required_cost_columns.issubset(backtest_summary.columns):
        raise AssertionError("Backtest summary is missing execution-cost columns")
    if (backtest_summary[list(required_cost_columns)] < 0).any().any():
        raise AssertionError("Backtest costs must be nonnegative")
    required_benchmark_columns = {
        "gross_total_return",
        "benchmark_total_return",
        "active_total_return",
        "annual_excess_return",
        "information_ratio",
        "avg_one_way_turnover",
        "annualized_one_way_turnover",
        "cost_drag",
    }
    if not required_benchmark_columns.issubset(backtest_summary.columns):
        raise AssertionError("Backtest summary is missing benchmark/excess or turnover columns")

    benchmark_excess = pd.read_csv(root / "backtest" / "benchmark_excess_returns.csv")
    if benchmark_excess.empty or set(benchmark_excess["mode"]) != expected_modes:
        raise AssertionError("Benchmark/excess daily table is incomplete")
    if not np.allclose(
        benchmark_excess["active_return"],
        benchmark_excess["net_return"] - benchmark_excess["benchmark_return"],
    ):
        raise AssertionError("Active returns do not equal strategy minus benchmark")
    if benchmark_excess["benchmark_n_stock"].lt(20).any():
        raise AssertionError("Equal-weight benchmark has too few tradable stocks")

    sensitivity = pd.read_csv(root / "backtest" / "cost_sensitivity.csv")
    if len(sensitivity) != 6 or set(sensitivity["scenario"]) != {
        "optimistic", "base", "pessimistic"
    }:
        raise AssertionError("Cost-sensitivity scenarios are incomplete")
    for _mode, group in sensitivity.groupby("mode"):
        ordered = group.set_index("scenario").loc[["optimistic", "base", "pessimistic"]]
        if not np.all(np.diff(ordered["total_cost"].to_numpy()) >= -1e-8):
            raise AssertionError("Cost-sensitivity costs are not monotonic")
        if not np.all(np.diff(ordered["total_return"].to_numpy()) <= 1e-8):
            raise AssertionError("Cost-sensitivity returns are not monotonic")

    ablation = pd.read_csv(root / "backtest" / "factor_ablation.csv")
    if len(ablation) != 14 or set(ablation["n_factors"]) != set(range(1, 8)):
        raise AssertionError("Factor ablation table is incomplete")
    correlation = pd.read_csv(
        root / "evaluation" / "factor_correlation.csv", index_col="factor"
    )
    if correlation.shape != (7, 7) or not np.allclose(
        correlation.to_numpy(), correlation.to_numpy().T, equal_nan=True
    ):
        raise AssertionError("Factor-correlation matrix is invalid")
    if not np.allclose(np.diag(correlation), 1.0):
        raise AssertionError("Factor-correlation diagonal is not one")

    turnover_control = pd.read_csv(
        root / "backtest" / "turnover_control_comparison.csv"
    )
    expected_turnover_policies = {
        "daily_top10",
        "five_day_top10",
        "daily_buffer20",
        "five_day_buffer20",
        "five_day_buffer60_top30",
    }
    if len(turnover_control) != 10 or set(turnover_control["turnover_policy"]) != expected_turnover_policies:
        raise AssertionError("Turnover-control comparison is incomplete")
    recommended_rows = turnover_control.loc[turnover_control["recommended"].astype(bool)]
    if len(recommended_rows) != 2 or not recommended_rows["turnover_policy"].eq(
        "five_day_buffer60_top30"
    ).all():
        raise AssertionError("Recommended turnover policy is not identified consistently")
    for mode in expected_modes:
        mode_rows = turnover_control.loc[turnover_control["mode"].eq(mode)].set_index(
            "turnover_policy"
        )
        if not (
            mode_rows.at["five_day_buffer60_top30", "avg_one_way_turnover"]
            < mode_rows.at["daily_top10", "avg_one_way_turnover"]
        ):
            raise AssertionError("Recommended policy did not lower turnover")
        if not (
            mode_rows.at["five_day_buffer60_top30", "total_cost"]
            < mode_rows.at["daily_top10", "total_cost"]
        ):
            raise AssertionError("Recommended policy did not lower total costs")
    turnover_results = pd.read_csv(root / "backtest" / "turnover_control_results.csv")
    expected_turnover_rows = int(turnover_control["n_periods"].sum())
    if len(turnover_results) != expected_turnover_rows:
        raise AssertionError("Turnover-control daily results have unexpected row count")
    five_day_results = turnover_results.loc[
        turnover_results["turnover_policy"].str.startswith("five_day")
    ]
    rebalance_share = five_day_results.groupby(["turnover_policy", "mode"])[
        "rebalanced"
    ].mean()
    if not rebalance_share.between(0.19, 0.21).all():
        raise AssertionError("Five-day policies are not rebalancing at the expected frequency")
    turnover_trades = pd.read_csv(
        root / "backtest" / "turnover_control_trades.csv", dtype={"code": str}
    )
    if (
        turnover_trades["trade_weight"].gt(0)
        & ~turnover_trades["buyable"].astype(bool)
    ).any():
        raise AssertionError("Turnover-control audit contains a blocked buy")
    if (
        turnover_trades["trade_weight"].lt(0)
        & ~turnover_trades["sellable"].astype(bool)
    ).any():
        raise AssertionError("Turnover-control audit contains a blocked sell")

    optimized_metrics = pd.read_csv(root / "backtest" / "optimized_strategy_metrics.csv")
    if len(optimized_metrics) != 6 or set(optimized_metrics["sample"]) != {
        "research80", "holdout20", "full"
    }:
        raise AssertionError("Optimized strategy split metrics are incomplete")
    optimized_full_open = optimized_metrics.loc[
        optimized_metrics["sample"].eq("full")
        & optimized_metrics["mode"].eq("next_open_to_open")
    ]
    if len(optimized_full_open) != 1:
        raise AssertionError("Missing optimized next-open full-sample metrics")
    if optimized_full_open["total_return"].iloc[0] <= 0.25:
        raise AssertionError("Optimized next-open total return did not exceed 25%")
    if optimized_full_open["annual_return"].iloc[0] <= 0.25:
        raise AssertionError("Optimized next-open annual return did not exceed 25%")

    neutralization = json.loads(
        (root / "evaluation" / "neutralization_status.json").read_text(encoding="utf-8")
    )
    if neutralization["status"] not in {"APPLIED", "SKIPPED_NO_EXPOSURES"}:
        raise AssertionError("Unexpected neutralization status")

    trades = pd.read_csv(root / "backtest" / "backtest_trades.csv", dtype={"code": str})
    if trades.empty:
        raise AssertionError("Backtest trade audit is empty")
    if (trades["trade_weight"].gt(0) & ~trades["buyable"].astype(bool)).any():
        raise AssertionError("A blocked buy was executed")
    if (trades["trade_weight"].lt(0) & ~trades["sellable"].astype(bool)).any():
        raise AssertionError("A blocked sell was executed")
    trade_cost_sum = (
        trades["base_slippage_cost"] + trades["impact_cost"] + trades["sell_fee"]
    )
    if not np.allclose(trades["total_cost"], trade_cost_sum):
        raise AssertionError("Trade cost components do not add to total_cost")

    backtest_results = pd.read_csv(root / "backtest" / "backtest_results.csv")
    signal_dates = pd.to_datetime(backtest_results["signal_date"].astype(str))
    trade_dates = pd.to_datetime(backtest_results["trade_date"].astype(str))
    if not trade_dates.gt(signal_dates).all():
        raise AssertionError("Backtest contains same-day signal execution")

    model_path = root / "lstm" / "model.pt"
    if not model_path.exists():
        raise AssertionError("Missing PyTorch LSTM checkpoint: lstm/model.pt")
    checkpoint = torch.load(model_path, map_location="cpu", weights_only=True)
    required_checkpoint_keys = {
        "model_type", "model_state_dict", "input_size", "hidden_size",
        "feature_mean", "feature_std", "features", "code", "sequence_length",
        "best_epoch", "device_used", "walk_forward", "codes", "n_dates",
        "fold_count",
    }
    required_checkpoint_keys.discard("code")
    if not required_checkpoint_keys.issubset(checkpoint):
        raise AssertionError(f"Missing LSTM checkpoint keys: {sorted(required_checkpoint_keys - set(checkpoint))}")
    if checkpoint["model_type"] != "torch.nn.LSTM" or checkpoint["features"] != [
            "log_return", "range", "log_volume", "buy_sell_imbalance", "time_position"]:
        raise AssertionError("Unexpected PyTorch LSTM model metadata")
    if not checkpoint["walk_forward"] or len(checkpoint["codes"]) < 2:
        raise AssertionError("LSTM is not a multi-stock walk-forward model")
    if checkpoint["n_dates"] <= 60 or checkpoint["fold_count"] < 2:
        raise AssertionError("LSTM time span or fold count is insufficient")
    test_fold_rows = walk_forward_metrics["split"].eq("test")
    if test_fold_rows.sum() != checkpoint["fold_count"]:
        raise AssertionError("Walk-forward fold metrics do not match checkpoint")
    predictions = pd.read_csv(
        root / "lstm" / "test_predictions.csv",
        dtype={"date": str, "code": str},
    )
    if predictions["fold"].nunique() < 2 or predictions["code"].nunique() < 2:
        raise AssertionError("Walk-forward predictions lack multiple folds or stocks")
    if predictions.groupby("date")["fold"].nunique().gt(1).any():
        raise AssertionError("Walk-forward OOS date appears in multiple test folds")
    if predictions.duplicated(["date", "code", "minute"]).any():
        raise AssertionError("Duplicate walk-forward OOS predictions")
    required_prediction_columns = {
        "logistic_probability_up", "logistic_prediction_up"
    }
    if not required_prediction_columns.issubset(predictions.columns):
        raise AssertionError("LSTM predictions are missing logistic baseline outputs")
    comparison = pd.read_csv(root / "lstm" / "model_comparison.csv")
    if set(comparison["model"]) != {
        "majority_class", "logistic_regression", "lstm"
    }:
        raise AssertionError("LSTM model-comparison baselines are incomplete")
    required_classification_columns = {
        "precision", "recall", "f1", "auc", "pr_auc", "tn", "fp", "fn", "tp"
    }
    if not required_classification_columns.issubset(comparison.columns):
        raise AssertionError("Classification metrics are incomplete")
    if not np.allclose(
        comparison[["tn", "fp", "fn", "tp"]].sum(axis=1),
        comparison["n_samples"],
    ):
        raise AssertionError("Confusion-matrix counts do not match sample totals")
    baseline_fold_metrics = pd.read_csv(
        root / "lstm" / "baseline_walk_forward_metrics.csv"
    )
    if len(baseline_fold_metrics) != checkpoint["fold_count"] * 3:
        raise AssertionError("Per-fold logistic baseline metrics are incomplete")
    holdings = pd.read_csv(root / "backtest" / "backtest_holdings.csv", dtype={"code": str})
    weight_sums = holdings.groupby(["strategy", "mode", "signal_date"])["target_weight"].sum()
    if weight_sums.gt(1 + 1e-10).any() or holdings["target_weight"].le(0).any():
        raise AssertionError("Constrained holdings have invalid weights")

    required_figures = [
        root / "evaluation" / "figures" / "factor_correlation_heatmap.png",
        root / "evaluation" / "figures" / "factor_long_short_nav.png",
        root / "evaluation" / "figures" / "avg_trade_size_surprise_20d_diagnostics.png",
        root / "evaluation" / "figures" / "price_volume_pressure_10d_diagnostics.png",
        root / "evaluation" / "figures" / "gap_intraday_divergence_5d_diagnostics.png",
        root / "backtest" / "figures" / "strategy_benchmark_excess.png",
        root / "backtest" / "figures" / "cost_and_return_breakdown.png",
        root / "backtest" / "figures" / "cost_sensitivity.png",
        root / "backtest" / "figures" / "factor_ablation.png",
        root / "backtest" / "figures" / "lookahead_diagnostic.png",
        root / "backtest" / "figures" / "turnover_control_comparison.png",
        root / "backtest" / "figures" / "turnover_control_nav.png",
        root / "backtest" / "figures" / "optimized_strategy_nav_drawdown.png",
        root / "backtest" / "figures" / "optimized_strategy_risk_metrics.png",
        root / "backtest" / "figures" / "optimized_period_split.png",
        root / "lstm" / "figures" / "model_comparison.png",
        root / "lstm" / "figures" / "lstm_confusion_matrix.png",
        root / "lstm" / "figures" / "fold_model_metrics.png",
    ]
    if any(not path.exists() or path.stat().st_size == 0 for path in required_figures):
        raise AssertionError("One or more required analysis figures are missing")

    audit = {
        "status": "PASS",
        "daily_fields": len(FIELDS),
        "daily_shape_each": [302, 300],
        "daily_kline_rows": len(kline),
        "daily_kline_columns": len(kline.columns),
        "minute_files": minute_file_count,
        "minute_shape_each": [253, 300],
        "minute_includes_1500": True,
        "no_trade_cells_checked": no_trade_checks,
        "factor_count": len(factor_summary),
        "backtest_variants": len(backtest_summary),
        "backtest_execution_delay": "t+1",
        "backtest_cost_model": "fee + slippage + square-root impact",
        "backtest_trade_rows": len(trades),
        "benchmark": "tradable stock-pool daily equal weight",
        "cost_sensitivity_scenarios": len(sensitivity),
        "factor_ablation_variants": len(ablation),
        "turnover_control_variants": len(turnover_control),
        "recommended_turnover_policy": "five_day_buffer60_top30",
        "recommended_daily_turnover_close": float(
            recommended_rows.loc[
                recommended_rows["mode"].eq("next_close_to_close"),
                "avg_one_way_turnover",
            ].iloc[0]
        ),
        "recommended_daily_turnover_open": float(
            recommended_rows.loc[
                recommended_rows["mode"].eq("next_open_to_open"),
                "avg_one_way_turnover",
            ].iloc[0]
        ),
        "factor_figures": len(list((root / "evaluation" / "figures").glob("*.png"))),
        "backtest_figures": len(list((root / "backtest" / "figures").glob("*.png"))),
        "neutralization_status": neutralization["status"],
        "lstm_walk_forward_folds": int(checkpoint["fold_count"]),
        "lstm_codes": len(checkpoint["codes"]),
        "lstm_dates": int(checkpoint["n_dates"]),
        "lstm_oos_samples": len(predictions),
        "lstm_comparison_models": comparison["model"].tolist(),
        "lstm_figures": len(list((root / "lstm" / "figures").glob("*.png"))),
        "lstm_framework": checkpoint["model_type"],
        "lstm_model_file": "lstm/model.pt",
    }
    reports = root / "reports"
    reports.mkdir(parents=True, exist_ok=True)
    (reports / "validation.json").write_text(json.dumps(audit, ensure_ascii=False, indent=2) + "\n",
                                              encoding="utf-8")
    return audit


def main() -> None:
    parser = argparse.ArgumentParser(description=__doc__)
    parser.add_argument("--root", type=Path, default=Path("."))
    args = parser.parse_args()
    print(json.dumps(validate(args.root.resolve()), ensure_ascii=False, indent=2))


if __name__ == "__main__":
    main()
\end{ProjectCode}
\clearpage
\section{11. \texttt{scripts/build\_report.py}}
根据最新CSV结果生成项目Markdown报告。

\textbf{物理行数：} 281\\
\textbf{SHA-256：} \texttt{23a317a80a23115eae8362e4c6e14b1d3a1eab24b22b8882160b2d88a820e5fd}

\begin{ProjectCode}
#!/usr/bin/env python3
"""Assemble current Chinese reports from generated analysis tables."""

from __future__ import annotations

import json
from pathlib import Path

import numpy as np
import pandas as pd


ROOT = Path(__file__).resolve().parents[1]


def pct(value: float) -> str:
    return f"{value:.2%}" if np.isfinite(value) else "NA"


def main() -> None:
    factors = pd.read_csv(ROOT / "evaluation" / "factor_summary.csv")
    correlation = pd.read_csv(
        ROOT / "evaluation" / "factor_correlation.csv", index_col="factor"
    )
    backtests = pd.read_csv(ROOT / "backtest" / "backtest_summary.csv")
    sensitivity = pd.read_csv(ROOT / "backtest" / "cost_sensitivity.csv")
    ablation = pd.read_csv(ROOT / "backtest" / "factor_ablation.csv")
    turnover = pd.read_csv(ROOT / "backtest" / "turnover_control_comparison.csv")
    optimized_metrics = pd.read_csv(
        ROOT / "backtest" / "optimized_strategy_metrics.csv"
    )
    comparison = pd.read_csv(ROOT / "lstm" / "model_comparison.csv")
    lstm = pd.read_csv(ROOT / "lstm" / "metrics.csv").iloc[0]
    metadata = pd.read_csv(ROOT / "factors" / "factor_metadata.csv")
    neutralization = json.loads(
        (ROOT / "evaluation" / "neutralization_status.json").read_text(encoding="utf-8")
    )
    assumptions = json.loads(
        (ROOT / "backtest" / "execution_assumptions.json").read_text(encoding="utf-8")
    )
    recommended_policy = assumptions["recommended_turnover_policy"]
    model_metadata = json.loads(
        (ROOT / "lstm" / "model_metadata.json").read_text(encoding="utf-8")
    )
    reports = ROOT / "reports"
    reports.mkdir(parents=True, exist_ok=True)

    field_doc = """# 数据、字段与风险暴露说明

## 原始数据

`Time, Price, Volume, BSFlag`；`Price` 为乘 100 后的整数。`BSFlag=0/1/2` 分别代表主买、主卖、集合竞价。

## 复权与统计口径

- OHLC：`Price / 100 * adjfactor`。
- `volume`：成交量之和；`trade_count`：逐笔记录数；`amount`：`Price / 100 * Volume` 之和。
- 日频无成交价格以前一日收盘填充；分钟无成交 OHLC 以前一分钟收盘填充；非价格字段填 0。

## 行业与市值中性化

中性化接口接受 `date,code,industry,market_cap`，或不含日期的静态暴露文件。按日使用行业哑变量和对数市值回归，残差再标准化。题目股票代码已混淆，且原始数据没有行业和市值，所以当前运行状态为 `SKIPPED_NO_EXPOSURES`；代码不会用代码前缀或成交额伪造真实暴露。

## 输出结构

- 合并日 K：`kline_data/daily_kline.csv`，90,600 行。
- 日频：`processed/daily/<field>.csv`，每张 302 x 300。
- 分钟频：`processed/minute/<field>/<YYYYMMDD>.csv`，每张 253 x 300，包含 15:00。
"""
    (reports / "数据与字段说明.md").write_text(field_doc, encoding="utf-8")

    lines = [
        "# 逐笔成交量化项目总报告",
        "",
        "## 摘要",
        "",
        "项目覆盖 302 个交易日、300 只股票和 11 个基础字段。"
        "研究流程包括七因子评价、股票池等权基准、超额收益、执行约束回测、"
        "成本拆分与敏感性分析、因子相关性与消融，以及多股票 walk-forward LSTM。",
        "",
        "## 一、因子评价",
        "",
        f"- 中性化状态：`{neutralization['status']}`。"
        + (
            "已使用真实行业与市值暴露。"
            if neutralization["status"] == "APPLIED"
            else "数据未提供真实行业和市值，接口可用但本次不伪造风险暴露。"
        ),
        "",
        "| 因子 | 类型 | 公式 |",
        "|---|---|---|",
    ]
    for row in metadata.itertuples(index=False):
        lines.append(f"| {row.factor}（{row.name_zh}） | {row.kind} | `{row.formula}` |")
    lines += [
        "",
        "| 因子 | IC | ICIR | Rank IC | Rank ICIR | Q5-Q1总收益 |",
        "|---|---:|---:|---:|---:|---:|",
    ]
    for row in factors.itertuples(index=False):
        lines.append(
            f"| {row.factor} | {row.ic:.4f} | {row.icir:.3f} | "
            f"{row.rank_ic:.4f} | {row.rank_icir:.3f} | "
            f"{pct(row.long_short_total_return)} |"
        )
    off_diagonal = correlation.where(~np.eye(len(correlation), dtype=bool)).abs()
    maximum_pair = off_diagonal.stack().idxmax()
    maximum_correlation = correlation.loc[maximum_pair[0], maximum_pair[1]]
    lines += [
        "",
        "因子相关性绝对值最大的一组为 "
        f"`{maximum_pair[0]}` 与 `{maximum_pair[1]}`，相关系数 {maximum_correlation:.3f}。"
        "七个因子并非完全重复；新增因子使用保守信号乘数，避免探索性信号主导组合。",
        "",
        "## 二、基准、超额收益与成本",
        "",
        "基准为入场日可以买入股票的每日等权收益，不扣除基准交易成本。"
        "正式组合使用截至 t-1 日的 20 日历史 IC 均值，在 t 日形成信号并于 t+1 成交。",
        "",
        "| 成交口径 | 毛收益 | 基准收益 | 净收益 | 年化超额 | 信息比率 | 日均单边换手 | 年化单边换手 |",
        "|---|---:|---:|---:|---:|---:|---:|---:|",
    ]
    formal = backtests.loc[backtests["strategy"].eq("historical_20d_ic")]
    for row in formal.itertuples(index=False):
        lines.append(
            f"| {row.mode} | {pct(row.gross_total_return)} | "
            f"{pct(row.benchmark_total_return)} | {pct(row.total_return)} | "
            f"{pct(row.annual_excess_return)} | {row.information_ratio:.3f} | "
            f"{pct(row.avg_one_way_turnover)} | {row.annualized_one_way_turnover:.1f}x |"
        )
    lines += [
        "",
        "| 成交口径 | 卖出手续费 | 基础滑点 | 冲击成本 | 总成本 | 成本拖累 |",
        "|---|---:|---:|---:|---:|---:|",
    ]
    for row in formal.itertuples(index=False):
        lines.append(
            f"| {row.mode} | {row.total_sell_fee:,.0f} | "
            f"{row.total_slippage_cost:,.0f} | {row.total_impact_cost:,.0f} | "
            f"{row.total_cost:,.0f} | {pct(row.cost_drag)} |"
        )
    lines += [
        "",
        "组合在成本前跑赢等权基准，但约 52%—53% 的日均单边换手使冲击成本成为最大成本项，"
        "基准成本假设下净收益和超额收益均为负。",
        "",
        "## 三、成本敏感性",
        "",
        "乐观情景只收卖出手续费；基准情景使用双边 5bp 滑点和平方根冲击；"
        "悲观情景使用双边 10bp 滑点，并把冲击系数和上限提高 50%。",
        "",
        "| 情景 | 成交口径 | 净收益 | 夏普 | 总成本 |",
        "|---|---|---:|---:|---:|",
    ]
    for row in sensitivity.itertuples(index=False):
        lines.append(
            f"| {row.scenario} | {row.mode} | {pct(row.total_return)} | "
            f"{row.sharpe:.3f} | {row.total_cost:,.0f} |"
        )
    lines += [
        "",
        "成本敏感性使用最终推荐政策。乐观、基准和悲观情景用于判断收益是否依赖过低的交易摩擦假设。",
        "",
        "## 四、换手控制",
        "",
        "比较五种政策，包括原 Top10 方案和最终的 Top30、5 日调仓、Top60 退出缓冲区。"
        "缓冲区规则保留仍在退出排名以内的原持仓，再按最新得分补足目标股票数。",
        "",
        "| 政策 | 口径 | 持仓数 | 日均单边换手 | 总成本 | 净收益 | 年化收益 | 最大回撤 | 夏普 | 年化超额 |",
        "|---|---|---:|---:|---:|---:|---:|---:|---:|---:|",
    ]
    for row in turnover.itertuples(index=False):
        lines.append(
            f"| {row.turnover_policy} | {row.mode} | {int(row.portfolio_size)} | "
            f"{pct(row.avg_one_way_turnover)} | "
            f"{row.total_cost:,.0f} | {pct(row.total_return)} | "
            f"{pct(row.annual_return)} | {pct(row.max_drawdown)} | "
            f"{row.sharpe:.3f} | {pct(row.annual_excess_return)} |"
        )
    daily_close = turnover.loc[
        turnover["turnover_policy"].eq("daily_top10")
        & turnover["mode"].eq("next_close_to_close")
    ].iloc[0]
    daily_open = turnover.loc[
        turnover["turnover_policy"].eq("daily_top10")
        & turnover["mode"].eq("next_open_to_open")
    ].iloc[0]
    controlled_close = turnover.loc[
        turnover["turnover_policy"].eq(recommended_policy)
        & turnover["mode"].eq("next_close_to_close")
    ].iloc[0]
    controlled_open = turnover.loc[
        turnover["turnover_policy"].eq(recommended_policy)
        & turnover["mode"].eq("next_open_to_open")
    ].iloc[0]
    lines += [
        "",
        f"推荐使用 `{recommended_policy}`。该政策把日均单边换手从 "
        f"{pct(daily_close.avg_one_way_turnover)} 和 {pct(daily_open.avg_one_way_turnover)} "
        f"降到 {pct(controlled_close.avg_one_way_turnover)} 和 "
        f"{pct(controlled_open.avg_one_way_turnover)}；总成本分别从 "
        f"{daily_close.total_cost:,.0f} 和 {daily_open.total_cost:,.0f} 降到 "
        f"{controlled_close.total_cost:,.0f} 和 {controlled_open.total_cost:,.0f}。",
        "",
        "次日收盘口径净收益为 "
        f"{pct(controlled_close.total_return)}，累计主动收益 "
        f"{pct(controlled_close.active_total_return)}；次日开盘口径净收益 "
        f"{pct(controlled_open.total_return)}，累计主动收益 "
        f"{pct(controlled_open.active_total_return)}，年化超额收益 "
        f"{pct(controlled_open.annual_excess_return)}。",
        "",
        "### 研究区间与后20%检查",
        "",
        "| 样本 | 口径 | 区间 | 累计收益 | 年化收益 | 基准收益 | 累计主动收益 | 最大回撤 | 夏普 |",
        "|---|---|---|---:|---:|---:|---:|---:|---:|",
    ]
    for row in optimized_metrics.itertuples(index=False):
        lines.append(
            f"| {row.sample} | {row.mode} | {row.start_date}-{row.end_date} | "
            f"{pct(row.total_return)} | {pct(row.annual_return)} | "
            f"{pct(row.benchmark_total_return)} | {pct(row.active_total_return)} | "
            f"{pct(row.max_drawdown)} | {row.sharpe:.3f} |"
        )
    lines += [
        "",
        "## 五、因子组合消融",
        "",
        "| 组合 | 成交口径 | 因子数 | 综合Rank IC | 多空收益 | 策略净收益 |",
        "|---|---|---:|---:|---:|---:|",
    ]
    for row in ablation.itertuples(index=False):
        lines.append(
            f"| {row.ablation_step} | {row.mode} | {row.n_factors} | "
            f"{row.composite_rank_ic:.4f} | "
            f"{pct(row.composite_long_short_total_return)} | {pct(row.total_return)} |"
        )
    lines += [
        "",
        "消融结果用于检查三个新增因子的边际贡献。新增因子采用 0.10 的信号乘数，"
        "防止短样本中的强 IC 直接造成组合权重过度集中。",
        "",
        "## 六、前视偏差诊断",
        "",
        "`leaky_same_day_ic` 使用尚未发生的收益确定当日权重，只作为数据泄漏诊断。"
        "项目正式统计、基准比较和超额收益均使用 `historical_20d_ic`。",
        "",
        "## 七、多股票 Walk-forward LSTM",
        "",
        f"- 股票数：{int(lstm.n_codes)}；日期数：{int(lstm.n_dates)}；"
        f"样本外折数：{int(lstm.n_folds)}；样本外样本：{int(lstm.n_samples):,}。",
        f"- 日期范围：{model_metadata['date_start']} 至 {model_metadata['date_end']}。"
        "每折只用训练期估计标准化参数，验证期选模，测试窗口互不重叠。",
        "",
        "| 模型 | Accuracy | Balanced Accuracy | Precision | Recall | F1 | ROC AUC | PR AUC |",
        "|---|---:|---:|---:|---:|---:|---:|---:|",
    ]
    for row in comparison.itertuples(index=False):
        lines.append(
            f"| {row.model} | {pct(row.accuracy)} | {pct(row.balanced_accuracy)} | "
            f"{pct(row.precision)} | {pct(row.recall)} | {pct(row.f1)} | "
            f"{row.auc:.4f} | {row.pr_auc:.4f} |"
        )
    lines += [
        "",
        "LSTM 的 ROC AUC、PR AUC 和 F1 均高于逻辑回归，说明序列模型提供了额外排序信息。"
        f"但在 0.5 阈值下 Recall 仅 {pct(lstm.recall)}，不能仅凭 Accuracy 判断模型有效性。",
        "",
        "## 八、结论",
        "",
        "- 七因子组合存在成本前选股信息；新增因子覆盖单笔成交规模、价量压力和隔夜-日内背离。",
        "- Top30、5 日调仓加 Top60 缓冲区显著降低个股噪声、换手和成本；"
        "次日开盘口径的成本后累计收益与年化收益均超过 25%。",
        "- LSTM 在严格样本外窗口中优于逻辑回归基线，但阈值分类的上涨召回率较低，"
        "还需要阈值、概率分组和交易成本层面的策略检验。",
        "- 因子分组、相关性、净值、成本、消融和模型比较图保存在各输出目录的 `figures/` 文件夹。",
    ]
    (reports / "项目总报告.md").write_text("\n".join(lines) + "\n", encoding="utf-8")


if __name__ == "__main__":
    main()
\end{ProjectCode}
\clearpage
\section{12. \texttt{scripts/build\_explained\_assignment\_pdf.py}}
生成带代码解释的作业PDF。

\textbf{物理行数：} 828\\
\textbf{SHA-256：} \texttt{cccf65bccbd297aa725627f7a285dbe2dcbdc1bfaaa8421cae3c093b0db5b26b}

\begin{ProjectCode}
#!/usr/bin/env python3
"""Render the R-Markdown-style assignment walkthrough to a polished PDF."""

from __future__ import annotations

import argparse
import html
import math
import re
from pathlib import Path

import numpy as np
import pandas as pd
from reportlab.lib import colors
from reportlab.lib.enums import TA_CENTER, TA_LEFT
from reportlab.lib.pagesizes import A4
from reportlab.lib.styles import ParagraphStyle, getSampleStyleSheet
from reportlab.lib.units import mm
from reportlab.pdfbase import pdfmetrics
from reportlab.pdfbase.ttfonts import TTFont
from reportlab.platypus import (
    BaseDocTemplate,
    CondPageBreak,
    Flowable,
    Frame,
    PageBreak,
    PageTemplate,
    Paragraph,
    Preformatted,
    Spacer,
    Table,
    TableStyle,
)
from reportlab.platypus.tableofcontents import TableOfContents


ROOT = Path(__file__).resolve().parents[1]
FONT_PATH = Path("/System/Library/Fonts/Supplemental/Arial Unicode.ttf")
PAGE_WIDTH, PAGE_HEIGHT = A4
NAVY = colors.HexColor("#17324D")
BLUE = colors.HexColor("#286B8F")
CYAN = colors.HexColor("#44A6B5")
TEAL = colors.HexColor("#2F8F83")
ORANGE = colors.HexColor("#E88A3D")
INK = colors.HexColor("#263238")
MUTED = colors.HexColor("#64727B")
LIGHT = colors.HexColor("#F3F6F8")
GRID = colors.HexColor("#D8E0E5")


def register_fonts() -> None:
    if not FONT_PATH.exists():
        raise FileNotFoundError(f"Chinese font not found: {FONT_PATH}")
    pdfmetrics.registerFont(TTFont("CJK", str(FONT_PATH)))


def make_styles() -> dict[str, ParagraphStyle]:
    base = getSampleStyleSheet()
    return {
        "Title": ParagraphStyle(
            "Title", parent=base["Title"], fontName="CJK", fontSize=27,
            leading=38, textColor=NAVY, alignment=TA_CENTER, spaceAfter=8,
        ),
        "Subtitle": ParagraphStyle(
            "Subtitle", parent=base["Normal"], fontName="CJK", fontSize=12,
            leading=19, textColor=BLUE, alignment=TA_CENTER,
        ),
        "CoverMeta": ParagraphStyle(
            "CoverMeta", parent=base["Normal"], fontName="CJK", fontSize=9.5,
            leading=15, textColor=MUTED, alignment=TA_CENTER,
        ),
        "H1": ParagraphStyle(
            "H1", parent=base["Heading1"], fontName="CJK", fontSize=18,
            leading=25, textColor=NAVY, spaceBefore=15, spaceAfter=9,
            keepWithNext=True, wordWrap="CJK",
        ),
        "H2": ParagraphStyle(
            "H2", parent=base["Heading2"], fontName="CJK", fontSize=13.5,
            leading=20, textColor=BLUE, spaceBefore=12, spaceAfter=7,
            keepWithNext=True, wordWrap="CJK",
        ),
        "H3": ParagraphStyle(
            "H3", parent=base["Heading3"], fontName="CJK", fontSize=11,
            leading=17, textColor=TEAL, spaceBefore=9, spaceAfter=5,
            keepWithNext=True, wordWrap="CJK",
        ),
        "Body": ParagraphStyle(
            "Body", parent=base["BodyText"], fontName="CJK", fontSize=9.2,
            leading=15.2, textColor=INK, alignment=TA_LEFT, spaceAfter=6,
            wordWrap="CJK",
        ),
        "Quote": ParagraphStyle(
            "Quote", parent=base["BodyText"], fontName="CJK", fontSize=9.1,
            leading=15, textColor=BLUE, leftIndent=12, rightIndent=8,
            borderPadding=8, backColor=colors.HexColor("#EAF5F7"),
            spaceAfter=8, wordWrap="CJK",
        ),
        "Bullet": ParagraphStyle(
            "Bullet", parent=base["BodyText"], fontName="CJK", fontSize=9.1,
            leading=14.8, textColor=INK, leftIndent=15,
            bulletIndent=2, spaceAfter=3, wordWrap="CJK",
        ),
        "Code": ParagraphStyle(
            "Code", parent=base["Code"], fontName="Courier", fontSize=6.9,
            leading=9.3, textColor=colors.HexColor("#243746"), leftIndent=6,
            rightIndent=6, borderPadding=8, borderColor=GRID,
            borderWidth=0.5, backColor=colors.HexColor("#F6F8FA"),
            spaceBefore=3, spaceAfter=8,
        ),
        "Table": ParagraphStyle(
            "Table", parent=base["BodyText"], fontName="CJK", fontSize=7.2,
            leading=9.6, textColor=INK, wordWrap="CJK",
        ),
        "TableHeader": ParagraphStyle(
            "TableHeader", parent=base["BodyText"], fontName="CJK",
            fontSize=7.2, leading=9.6, textColor=colors.white,
            wordWrap="CJK",
        ),
        "Small": ParagraphStyle(
            "Small", parent=base["BodyText"], fontName="CJK", fontSize=7.8,
            leading=11.5, textColor=MUTED, wordWrap="CJK",
        ),
        "TOCTitle": ParagraphStyle(
            "TOCTitle", parent=base["Heading1"], fontName="CJK",
            fontSize=20, leading=28, textColor=NAVY, spaceAfter=12,
        ),
    }


class ReportDocTemplate(BaseDocTemplate):
    def __init__(self, filename: str, styles: dict[str, ParagraphStyle]) -> None:
        super().__init__(
            filename,
            pagesize=A4,
            leftMargin=20 * mm,
            rightMargin=20 * mm,
            topMargin=19 * mm,
            bottomMargin=18 * mm,
            title="逐笔成交量化作业：代码讲解与结果复盘",
            author="量化作业项目",
            subject="R Markdown 风格代码讲解、结果与验证",
        )
        self.styles = styles
        frame = Frame(
            self.leftMargin,
            self.bottomMargin,
            self.width,
            self.height,
            id="body",
            leftPadding=0,
            rightPadding=0,
            topPadding=0,
            bottomPadding=0,
        )
        self.addPageTemplates(
            PageTemplate(id="main", frames=[frame], onPage=self._draw_page)
        )
        self._bookmark_counter = 0

    def _draw_page(self, canvas, doc) -> None:
        if doc.page == 1:
            return
        canvas.saveState()
        canvas.setStrokeColor(GRID)
        canvas.setLineWidth(0.5)
        canvas.line(
            20 * mm,
            PAGE_HEIGHT - 13 * mm,
            PAGE_WIDTH - 20 * mm,
            PAGE_HEIGHT - 13 * mm,
        )
        canvas.setFont("CJK", 7.5)
        canvas.setFillColor(MUTED)
        canvas.drawString(
            20 * mm,
            PAGE_HEIGHT - 10 * mm,
            "逐笔成交量化作业｜代码讲解与结果复盘",
        )
        canvas.drawRightString(
            PAGE_WIDTH - 20 * mm,
            10 * mm,
            f"第 {doc.page - 1} 页",
        )
        canvas.restoreState()

    def afterFlowable(self, flowable) -> None:
        if not isinstance(flowable, Paragraph):
            return
        style = flowable.style.name
        if style not in {"H1", "H2", "H3"}:
            return
        level = {"H1": 0, "H2": 1, "H3": 2}[style]
        text = flowable.getPlainText()
        key = getattr(flowable, "_bookmark_name", None)
        if key is None:
            key = f"heading-{self._bookmark_counter}"
            flowable._bookmark_name = key
            self._bookmark_counter += 1
        self.canv.bookmarkPage(key)
        self.canv.addOutlineEntry(text, key, level=level, closed=False)
        if level == 0:
            self.notify("TOCEntry", (level, text, self.page - 1, key))


class PipelineChart(Flowable):
    def __init__(self, width: float) -> None:
        super().__init__()
        self.width = width
        self.height = 84

    def draw(self) -> None:
        c = self.canv
        nodes = [
            ("原始逐笔", "Time / Price\nVolume / BSFlag", BLUE),
            ("降采样复权", "日频 11 表\n分钟 3322 表", CYAN),
            ("因子评价", "4 因子\nIC / 分组", TEAL),
            ("回测与 LSTM", "Top 10 组合\n下一分钟涨跌", ORANGE),
            ("报告与审计", "CSV / model.pt\nvalidation PASS", NAVY),
        ]
        gap = 9
        box_w = (self.width - gap * (len(nodes) - 1)) / len(nodes)
        box_h = 58
        y = 12
        for i, (title, detail, color) in enumerate(nodes):
            x = i * (box_w + gap)
            c.setFillColor(colors.white)
            c.setStrokeColor(color)
            c.setLineWidth(1.2)
            c.roundRect(x, y, box_w, box_h, 5, stroke=1, fill=1)
            c.setFillColor(color)
            c.roundRect(x, y + box_h - 18, box_w, 18, 5, stroke=0, fill=1)
            c.rect(x, y + box_h - 18, box_w, 9, stroke=0, fill=1)
            c.setFont("CJK", 8.3)
            c.setFillColor(colors.white)
            c.drawCentredString(x + box_w / 2, y + box_h - 13, title)
            c.setFillColor(INK)
            c.setFont("CJK", 7.1)
            for j, line in enumerate(detail.splitlines()):
                c.drawCentredString(
                    x + box_w / 2,
                    y + 25 - j * 11,
                    line,
                )
            if i < len(nodes) - 1:
                arrow_x = x + box_w + 1
                arrow_y = y + box_h / 2
                c.setStrokeColor(MUTED)
                c.setFillColor(MUTED)
                c.line(arrow_x, arrow_y, arrow_x + gap - 3, arrow_y)
                c.line(
                    arrow_x + gap - 6,
                    arrow_y + 3,
                    arrow_x + gap - 3,
                    arrow_y,
                )
                c.line(
                    arrow_x + gap - 6,
                    arrow_y - 3,
                    arrow_x + gap - 3,
                    arrow_y,
                )


class FactorBarChart(Flowable):
    def __init__(self, width: float) -> None:
        super().__init__()
        self.width = width
        self.height = 205
        self.labels = ["成交额多窗口", "日内振幅", "5日反转", "不平衡冲击"]
        self.ic = [-0.077411, -0.039620, 0.030271, -0.031413]
        self.rank_ic = [-0.106980, -0.078314, 0.055783, -0.025137]

    def draw(self) -> None:
        c = self.canv
        left, right, bottom, top = 42, 12, 39, 24
        plot_w = self.width - left - right
        plot_h = self.height - bottom - top
        lo, hi = -0.12, 0.08

        def ymap(value: float) -> float:
            return bottom + (value - lo) / (hi - lo) * plot_h

        c.setFont("CJK", 7)
        for tick in np.arange(lo, hi + 0.001, 0.04):
            y = ymap(float(tick))
            c.setStrokeColor(GRID)
            c.setLineWidth(0.4)
            c.line(left, y, left + plot_w, y)
            c.setFillColor(MUTED)
            c.drawRightString(left - 5, y - 2.5, f"{tick:.2f}")
        c.setStrokeColor(NAVY)
        c.setLineWidth(0.8)
        c.line(left, ymap(0), left + plot_w, ymap(0))
        group_w = plot_w / len(self.labels)
        bar_w = min(18, group_w * 0.27)
        for i, label in enumerate(self.labels):
            cx = left + group_w * (i + 0.5)
            pairs = ((self.ic[i], BLUE), (self.rank_ic[i], ORANGE))
            for j, (value, color) in enumerate(pairs):
                x = cx + (j - 0.5) * bar_w
                y0, y1 = ymap(0), ymap(value)
                c.setFillColor(color)
                c.rect(
                    x - bar_w / 2,
                    min(y0, y1),
                    bar_w,
                    abs(y1 - y0),
                    stroke=0,
                    fill=1,
                )
            c.setFillColor(INK)
            c.setFont("CJK", 6.8)
            c.drawCentredString(cx, 23, label)
        c.setFont("CJK", 8)
        c.setFillColor(NAVY)
        c.drawString(left, self.height - 12, "因子平均 IC 与 Rank IC")
        legend_x = self.width - 142
        for j, (name, color) in enumerate((("IC", BLUE), ("Rank IC", ORANGE))):
            x = legend_x + j * 70
            c.setFillColor(color)
            c.rect(x, self.height - 17, 9, 6, stroke=0, fill=1)
            c.setFillColor(INK)
            c.setFont("CJK", 6.8)
            c.drawString(x + 13, self.height - 17, name)


class LineChart(Flowable):
    def __init__(
        self,
        width: float,
        title: str,
        series: list[tuple[str, list[float], colors.Color]],
        x_labels: tuple[str, str, str],
        y_format: str = "{:.2f}",
        force_range: tuple[float, float] | None = None,
    ) -> None:
        super().__init__()
        self.width = width
        self.height = 205
        self.title = title
        self.series = series
        self.x_labels = x_labels
        self.y_format = y_format
        self.force_range = force_range

    @staticmethod
    def _nice_range(values: list[float]) -> tuple[float, float]:
        lo, hi = min(values), max(values)
        if math.isclose(lo, hi):
            return lo - 0.1, hi + 0.1
        pad = (hi - lo) * 0.12
        return lo - pad, hi + pad

    def draw(self) -> None:
        c = self.canv
        left, right, bottom, top = 45, 14, 35, 29
        plot_w = self.width - left - right
        plot_h = self.height - bottom - top
        all_values = [
            value
            for _, values, _ in self.series
            for value in values
            if np.isfinite(value)
        ]
        lo, hi = self.force_range or self._nice_range(all_values)

        def ymap(value: float) -> float:
            return bottom + (value - lo) / (hi - lo) * plot_h

        c.setFont("CJK", 7)
        for i in range(5):
            value = lo + (hi - lo) * i / 4
            y = ymap(value)
            c.setStrokeColor(GRID)
            c.setLineWidth(0.4)
            c.line(left, y, left + plot_w, y)
            c.setFillColor(MUTED)
            c.drawRightString(
                left - 5,
                y - 2.5,
                self.y_format.format(value),
            )
        c.setStrokeColor(NAVY)
        c.setLineWidth(0.7)
        c.line(left, bottom, left, bottom + plot_h)
        c.line(left, bottom, left + plot_w, bottom)
        for _, values, color in self.series:
            if len(values) < 2:
                continue
            c.setStrokeColor(color)
            c.setLineWidth(1.5)
            points = []
            for i, value in enumerate(values):
                x = left + plot_w * i / (len(values) - 1)
                points.append((x, ymap(value)))
            path = c.beginPath()
            path.moveTo(*points[0])
            for point in points[1:]:
                path.lineTo(*point)
            c.drawPath(path, stroke=1, fill=0)
        c.setFillColor(NAVY)
        c.setFont("CJK", 8)
        c.drawString(left, self.height - 12, self.title)
        legend_x = self.width - 210
        for i, (name, _, color) in enumerate(self.series):
            x = legend_x + i * 105
            c.setStrokeColor(color)
            c.setLineWidth(2.2)
            c.line(x, self.height - 15, x + 13, self.height - 15)
            c.setFillColor(INK)
            c.setFont("CJK", 6.7)
            c.drawString(x + 17, self.height - 18, name)
        c.setFillColor(MUTED)
        c.setFont("CJK", 6.7)
        c.drawString(left, 20, self.x_labels[0])
        c.drawCentredString(
            left + plot_w / 2,
            20,
            self.x_labels[1],
        )
        c.drawRightString(left + plot_w, 20, self.x_labels[2])


def backtest_chart(width: float) -> Flowable:
    table = pd.read_csv(
        ROOT / "backtest" / "backtest_results.csv",
        dtype={"signal_date": str},
    )
    table = table[table["strategy"].eq("historical_20d_ic")].copy()
    series = []
    x_labels = None
    for mode, name, color in (
        ("close_to_close", "历史 IC｜收盘-收盘", BLUE),
        ("next_open_to_open", "历史 IC｜次日开盘", ORANGE),
    ):
        part = table[table["mode"].eq(mode)].sort_values("signal_date")
        nav = (part["nav"] / 10_000_000.0).tolist()
        series.append((name, nav, color))
        dates = part["signal_date"].astype(str).tolist()
        if dates:
            x_labels = (
                dates[0],
                dates[len(dates) // 2],
                dates[-1],
            )
    return LineChart(
        width,
        "可交易修正版净值（初始净值 = 1）",
        series,
        x_labels or ("开始", "中间", "结束"),
        y_format="{:.2f}",
        force_range=(0.88, 1.72),
    )


def lstm_chart(width: float) -> Flowable:
    history = pd.read_csv(ROOT / "lstm" / "training_history.csv")
    labels = history["epoch"].astype(int).astype(str).tolist()
    return LineChart(
        width,
        "LSTM 训练与验证交叉熵",
        [
            ("训练损失", history["train_loss"].tolist(), BLUE),
            ("验证损失", history["val_loss"].tolist(), ORANGE),
        ],
        (labels[0], labels[len(labels) // 2], labels[-1]),
        y_format="{:.3f}",
        force_range=(0.635, 0.678),
    )


def inline_markup(text: str) -> str:
    marker = chr(96)
    token_re = re.compile(
        "(" + re.escape(marker) + "[^" + re.escape(marker) + "]+"
        + re.escape(marker) + r"|\*\*[^*]+\*\*)"
    )
    pieces: list[str] = []
    pos = 0
    for match in token_re.finditer(text):
        pieces.append(html.escape(text[pos:match.start()]))
        token = match.group(0)
        if token.startswith(marker):
            pieces.append(
                '<font name="CJK" color="#286B8F" size="8">'
                + html.escape(token[1:-1])
                + "</font>"
            )
        else:
            pieces.append("<b>" + html.escape(token[2:-2]) + "</b>")
        pos = match.end()
    pieces.append(html.escape(text[pos:]))
    return "".join(pieces)


def table_widths(n_cols: int, available: float) -> list[float]:
    if n_cols == 2:
        ratios = [0.32, 0.68]
    elif n_cols == 3:
        ratios = [0.26, 0.34, 0.40]
    elif n_cols == 4:
        ratios = [0.30, 0.28, 0.27, 0.15]
    elif n_cols == 5:
        ratios = [0.32, 0.17, 0.17, 0.17, 0.17]
    elif n_cols == 6:
        ratios = [0.30, 0.14, 0.14, 0.14, 0.14, 0.14]
    elif n_cols == 7:
        ratios = [0.25, 0.18, 0.114, 0.114, 0.114, 0.114, 0.114]
    else:
        ratios = [1 / n_cols] * n_cols
    return [available * ratio for ratio in ratios]


def parse_table_row(line: str) -> list[str]:
    return [cell.strip() for cell in line.strip().strip("|").split("|")]


def make_table(
    raw_rows: list[list[str]],
    styles: dict[str, ParagraphStyle],
    available: float,
) -> Table:
    header, *body = raw_rows
    rows = [
        [
            Paragraph(inline_markup(cell), styles["TableHeader"])
            for cell in header
        ],
        *[
            [
                Paragraph(inline_markup(cell), styles["Table"])
                for cell in row
            ]
            for row in body
        ],
    ]
    n_cols = len(header)
    font_size = 6.5 if n_cols >= 7 else 7.2
    table = Table(
        rows,
        colWidths=table_widths(n_cols, available),
        repeatRows=1,
        hAlign="LEFT",
        splitByRow=1,
    )
    table.setStyle(TableStyle([
        ("BACKGROUND", (0, 0), (-1, 0), NAVY),
        ("TEXTCOLOR", (0, 0), (-1, 0), colors.white),
        ("FONTNAME", (0, 0), (-1, -1), "CJK"),
        ("FONTSIZE", (0, 0), (-1, -1), font_size),
        ("LEADING", (0, 0), (-1, -1), font_size + 2.2),
        ("VALIGN", (0, 0), (-1, -1), "TOP"),
        ("GRID", (0, 0), (-1, -1), 0.35, GRID),
        ("ROWBACKGROUNDS", (0, 1), (-1, -1), [colors.white, LIGHT]),
        ("LEFTPADDING", (0, 0), (-1, -1), 4),
        ("RIGHTPADDING", (0, 0), (-1, -1), 4),
        ("TOPPADDING", (0, 0), (-1, -1), 4),
        ("BOTTOMPADDING", (0, 0), (-1, -1), 4),
    ]))
    return table


def parse_markdown(
    source: Path,
    styles: dict[str, ParagraphStyle],
    available: float,
) -> list[Flowable]:
    lines = source.read_text(encoding="utf-8").splitlines()
    start = next(
        index
        for index, line in enumerate(lines)
        if line.startswith("## 报告说明")
    )
    lines = lines[start:]
    story: list[Flowable] = []
    fence = chr(96) * 3
    i = 0
    first_h1 = True
    while i < len(lines):
        line = lines[i].rstrip()
        if not line:
            story.append(Spacer(1, 2.5))
            i += 1
            continue
        if line.startswith(fence):
            i += 1
            code_lines = []
            while i < len(lines) and not lines[i].startswith(fence):
                code_lines.append(lines[i].rstrip())
                i += 1
            code = "\n".join(code_lines)
            story.append(
                Preformatted(
                    html.escape(code),
                    styles["Code"],
                    maxLineLength=96,
                )
            )
            i += 1
            continue
        if (
            line.startswith("|")
            and i + 1 < len(lines)
            and lines[i + 1].startswith("|")
        ):
            raw_rows: list[list[str]] = []
            while i < len(lines) and lines[i].startswith("|"):
                raw_rows.append(parse_table_row(lines[i]))
                i += 1
            separator = raw_rows[1] if len(raw_rows) >= 2 else []
            if separator and all(
                re.fullmatch(r":?-+:?", cell) for cell in separator
            ):
                raw_rows.pop(1)
            story.extend([
                make_table(raw_rows, styles, available),
                Spacer(1, 7),
            ])
            continue
        if line == "[[PIPELINE_CHART]]":
            story.extend([
                Spacer(1, 5),
                PipelineChart(available),
                Spacer(1, 8),
            ])
            i += 1
            continue
        if line == "[[FACTOR_CHART]]":
            story.extend([
                Spacer(1, 5),
                FactorBarChart(available),
                Spacer(1, 8),
            ])
            i += 1
            continue
        if line == "[[BACKTEST_CHART]]":
            story.extend([
                Spacer(1, 5),
                backtest_chart(available),
                Spacer(1, 8),
            ])
            i += 1
            continue
        if line == "[[LSTM_CHART]]":
            story.extend([
                Spacer(1, 5),
                lstm_chart(available),
                Spacer(1, 8),
            ])
            i += 1
            continue
        if line.startswith("## "):
            if not first_h1:
                story.append(CondPageBreak(70 * mm))
            first_h1 = False
            story.append(
                Paragraph(inline_markup(line[3:]), styles["H1"])
            )
            i += 1
            continue
        if line.startswith("### "):
            story.append(
                Paragraph(inline_markup(line[4:]), styles["H2"])
            )
            i += 1
            continue
        if line.startswith("#### "):
            story.append(
                Paragraph(inline_markup(line[5:]), styles["H3"])
            )
            i += 1
            continue
        if line.startswith("> "):
            story.append(
                Paragraph(inline_markup(line[2:]), styles["Quote"])
            )
            i += 1
            continue
        if line.startswith("- "):
            story.append(
                Paragraph(
                    inline_markup(line[2:]),
                    styles["Bullet"],
                    bulletText="•",
                )
            )
            i += 1
            continue
        numbered = re.match(r"^(\d+)\.\s+(.*)$", line)
        if numbered:
            story.append(
                Paragraph(
                    inline_markup(numbered.group(2)),
                    styles["Bullet"],
                    bulletText=numbered.group(1) + ".",
                )
            )
            i += 1
            continue
        story.append(Paragraph(inline_markup(line), styles["Body"]))
        i += 1
    return story


def cover(
    styles: dict[str, ParagraphStyle],
    available: float,
) -> list[Flowable]:
    status_data = [
        [
            Paragraph("数据", styles["TableHeader"]),
            Paragraph("因子", styles["TableHeader"]),
            Paragraph("回测", styles["TableHeader"]),
            Paragraph("LSTM", styles["TableHeader"]),
        ],
        [
            Paragraph("302 日 x 300 股<br/>11 个字段", styles["Table"]),
            Paragraph("4 个<br/>IC + 五分组", styles["Table"]),
            Paragraph("4 个版本<br/>Top 10", styles["Table"]),
            Paragraph("PyTorch<br/>流程完成", styles["Table"]),
        ],
    ]
    status = Table(
        status_data,
        colWidths=[available / 4] * 4,
        rowHeights=[24, 46],
    )
    status.setStyle(TableStyle([
        ("BACKGROUND", (0, 0), (-1, 0), NAVY),
        ("ROWBACKGROUNDS", (0, 1), (-1, -1), [LIGHT]),
        ("GRID", (0, 0), (-1, -1), 0.5, GRID),
        ("VALIGN", (0, 0), (-1, -1), "MIDDLE"),
        ("ALIGN", (0, 0), (-1, -1), "CENTER"),
    ]))
    return [
        Spacer(1, 35 * mm),
        Paragraph("逐笔成交量化作业", styles["Title"]),
        Paragraph("代码讲解与结果复盘", styles["Title"]),
        Spacer(1, 5 * mm),
        Paragraph(
            "R Markdown 风格｜题目要求・关键代码・输出表格・结果解释・验证结论",
            styles["Subtitle"],
        ),
        Spacer(1, 17 * mm),
        status,
        Spacer(1, 15 * mm),
        Paragraph("验证状态：PASS｜单元测试：4/4 通过", styles["Subtitle"]),
        Spacer(1, 5 * mm),
        Paragraph(
            "数据区间：2025-04-01 至 2026-06-30<br/>"
            "生成日期：2026-07-28｜结果以当前工程产物为准",
            styles["CoverMeta"],
        ),
        PageBreak(),
    ]


def build_pdf(source: Path, output: Path) -> None:
    register_fonts()
    styles = make_styles()
    output.parent.mkdir(parents=True, exist_ok=True)
    doc = ReportDocTemplate(str(output), styles)

    toc = TableOfContents()
    toc.levelStyles = [
        ParagraphStyle(
            "TOC1",
            fontName="CJK",
            fontSize=8.8,
            leading=14,
            leftIndent=0,
            firstLineIndent=0,
            textColor=NAVY,
            spaceBefore=1.5,
        ),
        ParagraphStyle(
            "TOC2",
            fontName="CJK",
            fontSize=7.4,
            leading=10.6,
            leftIndent=14,
            firstLineIndent=0,
            textColor=BLUE,
        ),
        ParagraphStyle(
            "TOC3",
            fontName="CJK",
            fontSize=7,
            leading=10,
            leftIndent=28,
            firstLineIndent=0,
            textColor=MUTED,
        ),
    ]
    story = cover(styles, doc.width)
    story.extend([
        Paragraph("目录", styles["TOCTitle"]),
        Paragraph(
            "本报告按代码执行链路组织；页码由最终 PDF 自动生成。",
            styles["Small"],
        ),
        Spacer(1, 5),
        toc,
        PageBreak(),
    ])
    story.extend(parse_markdown(source, styles, doc.width))
    doc.multiBuild(story)


def main() -> None:
    parser = argparse.ArgumentParser(description=__doc__)
    parser.add_argument(
        "--source",
        type=Path,
        default=ROOT / "reports" / "量化作业代码讲解报告.md",
    )
    parser.add_argument(
        "--output",
        type=Path,
        default=ROOT / "output" / "pdf" / "量化作业代码讲解报告.pdf",
    )
    args = parser.parse_args()
    build_pdf(args.source.resolve(), args.output.resolve())
    print(args.output.resolve())


if __name__ == "__main__":
    main()
\end{ProjectCode}
\clearpage
\section{13. \texttt{scripts/build\_final\_illustrated\_report.py}}
最终图文报告的通用ReportLab排版组件。

\textbf{物理行数：} 797\\
\textbf{SHA-256：} \texttt{aa80716ed1f3d93ef4328d5d92779c4adfe31490b35401383a72b3347e1a7dc2}

\begin{ProjectCode}
#!/usr/bin/env python3
"""Build the final illustrated Chinese quant-assignment report as Markdown and PDF."""

from __future__ import annotations

import json
from pathlib import Path

import numpy as np
import pandas as pd
from PIL import Image as PILImage
from reportlab.lib import colors
from reportlab.lib.enums import TA_CENTER, TA_LEFT
from reportlab.lib.pagesizes import A4
from reportlab.lib.styles import ParagraphStyle, getSampleStyleSheet
from reportlab.lib.units import mm
from reportlab.pdfbase import pdfmetrics
from reportlab.pdfbase.ttfonts import TTFont
from reportlab.platypus import (
    BaseDocTemplate,
    Frame,
    HRFlowable,
    Image,
    PageBreak,
    PageTemplate,
    Paragraph,
    Spacer,
    Table,
    TableStyle,
)
from reportlab.platypus.tableofcontents import TableOfContents


ROOT = Path(__file__).resolve().parents[1]
REPORTS = ROOT / "reports"
PDF_OUTPUT = ROOT / "output" / "pdf"
FONT_PATH = Path("/System/Library/Fonts/Supplemental/Arial Unicode.ttf")
PAGE_WIDTH, PAGE_HEIGHT = A4

NAVY = colors.HexColor("#17324D")
BLUE = colors.HexColor("#286B8F")
TEAL = colors.HexColor("#2F8F83")
ORANGE = colors.HexColor("#E88A3D")
RED = colors.HexColor("#B64B4B")
INK = colors.HexColor("#263238")
MUTED = colors.HexColor("#64727B")
LIGHT = colors.HexColor("#F3F6F8")
PALE_BLUE = colors.HexColor("#EAF3F8")
PALE_GREEN = colors.HexColor("#EAF5F1")
GRID = colors.HexColor("#D8E0E5")


def pct(value: float) -> str:
    return f"{value:.2%}" if np.isfinite(value) else "NA"


def register_font() -> None:
    if not FONT_PATH.exists():
        raise FileNotFoundError(f"Chinese font not found: {FONT_PATH}")
    pdfmetrics.registerFont(TTFont("CJK", str(FONT_PATH)))


def styles() -> dict[str, ParagraphStyle]:
    base = getSampleStyleSheet()
    return {
        "Title": ParagraphStyle(
            "Title",
            parent=base["Title"],
            fontName="CJK",
            fontSize=25,
            leading=36,
            textColor=NAVY,
            alignment=TA_CENTER,
            spaceAfter=10,
            wordWrap="CJK",
        ),
        "Subtitle": ParagraphStyle(
            "Subtitle",
            parent=base["Normal"],
            fontName="CJK",
            fontSize=12,
            leading=20,
            textColor=BLUE,
            alignment=TA_CENTER,
            wordWrap="CJK",
        ),
        "Cover": ParagraphStyle(
            "Cover",
            parent=base["Normal"],
            fontName="CJK",
            fontSize=9.5,
            leading=16,
            textColor=MUTED,
            alignment=TA_CENTER,
            wordWrap="CJK",
        ),
        "H1": ParagraphStyle(
            "H1",
            parent=base["Heading1"],
            fontName="CJK",
            fontSize=18,
            leading=26,
            textColor=NAVY,
            spaceBefore=8,
            spaceAfter=10,
            keepWithNext=True,
            wordWrap="CJK",
        ),
        "H2": ParagraphStyle(
            "H2",
            parent=base["Heading2"],
            fontName="CJK",
            fontSize=13,
            leading=20,
            textColor=BLUE,
            spaceBefore=10,
            spaceAfter=7,
            keepWithNext=True,
            wordWrap="CJK",
        ),
        "Body": ParagraphStyle(
            "Body",
            parent=base["BodyText"],
            fontName="CJK",
            fontSize=9.2,
            leading=15.5,
            textColor=INK,
            alignment=TA_LEFT,
            spaceAfter=7,
            wordWrap="CJK",
        ),
        "Bullet": ParagraphStyle(
            "Bullet",
            parent=base["BodyText"],
            fontName="CJK",
            fontSize=9,
            leading=14.5,
            leftIndent=15,
            bulletIndent=2,
            textColor=INK,
            spaceAfter=4,
            wordWrap="CJK",
        ),
        "Callout": ParagraphStyle(
            "Callout",
            parent=base["BodyText"],
            fontName="CJK",
            fontSize=9.3,
            leading=15.5,
            leftIndent=9,
            rightIndent=9,
            borderPadding=9,
            backColor=PALE_GREEN,
            borderColor=TEAL,
            borderWidth=0.7,
            textColor=INK,
            spaceBefore=5,
            spaceAfter=9,
            wordWrap="CJK",
        ),
        "Warning": ParagraphStyle(
            "Warning",
            parent=base["BodyText"],
            fontName="CJK",
            fontSize=9.1,
            leading=15,
            leftIndent=9,
            rightIndent=9,
            borderPadding=9,
            backColor=colors.HexColor("#FFF4E8"),
            borderColor=ORANGE,
            borderWidth=0.7,
            textColor=INK,
            spaceBefore=5,
            spaceAfter=9,
            wordWrap="CJK",
        ),
        "Caption": ParagraphStyle(
            "Caption",
            parent=base["Normal"],
            fontName="CJK",
            fontSize=7.8,
            leading=11.5,
            textColor=MUTED,
            alignment=TA_CENTER,
            spaceBefore=3,
            spaceAfter=8,
            wordWrap="CJK",
        ),
        "Table": ParagraphStyle(
            "Table",
            parent=base["BodyText"],
            fontName="CJK",
            fontSize=7.2,
            leading=10,
            textColor=INK,
            wordWrap="CJK",
        ),
        "TableHeader": ParagraphStyle(
            "TableHeader",
            parent=base["BodyText"],
            fontName="CJK",
            fontSize=7.2,
            leading=10,
            textColor=colors.white,
            wordWrap="CJK",
        ),
        "TOC": ParagraphStyle(
            "TOC",
            parent=base["Heading1"],
            fontName="CJK",
            fontSize=20,
            leading=28,
            textColor=NAVY,
            spaceAfter=12,
        ),
        "Small": ParagraphStyle(
            "Small",
            parent=base["BodyText"],
            fontName="CJK",
            fontSize=7.6,
            leading=11.5,
            textColor=MUTED,
            wordWrap="CJK",
        ),
    }


class ReportTemplate(BaseDocTemplate):
    def __init__(
        self,
        filename: str,
        style_map: dict[str, ParagraphStyle],
        *,
        title: str = "逐笔成交量化作业最终图文报告",
        subject: str = "数据处理、因子研究、换手控制回测与LSTM样本外验证",
        header: str = "逐笔成交量化作业最终图文报告",
    ) -> None:
        super().__init__(
            filename,
            pagesize=A4,
            leftMargin=19 * mm,
            rightMargin=19 * mm,
            topMargin=18 * mm,
            bottomMargin=18 * mm,
            title=title,
            author="量化作业项目",
            subject=subject,
        )
        self.style_map = style_map
        self.header = header
        frame = Frame(
            self.leftMargin,
            self.bottomMargin,
            self.width,
            self.height,
            id="body",
            leftPadding=0,
            rightPadding=0,
            topPadding=0,
            bottomPadding=0,
        )
        self.addPageTemplates(PageTemplate(id="main", frames=[frame], onPage=self._page))
        self._bookmark = 0

    def beforeDocument(self) -> None:
        """Keep bookmark names stable across TableOfContents build passes."""
        self._bookmark = 0

    def _page(self, canvas, doc) -> None:
        if doc.page == 1:
            return
        canvas.saveState()
        canvas.setStrokeColor(GRID)
        canvas.setLineWidth(0.5)
        canvas.line(19 * mm, PAGE_HEIGHT - 12 * mm, PAGE_WIDTH - 19 * mm, PAGE_HEIGHT - 12 * mm)
        canvas.setFont("CJK", 7.4)
        canvas.setFillColor(MUTED)
        canvas.drawString(19 * mm, PAGE_HEIGHT - 9 * mm, self.header)
        canvas.drawRightString(PAGE_WIDTH - 19 * mm, 9 * mm, f"第 {doc.page - 1} 页")
        canvas.restoreState()

    def afterFlowable(self, flowable) -> None:
        if not isinstance(flowable, Paragraph) or flowable.style.name not in {"H1", "H2"}:
            return
        level = 0 if flowable.style.name == "H1" else 1
        text = flowable.getPlainText()
        key = f"heading-{self._bookmark}"
        self._bookmark += 1
        self.canv.bookmarkPage(key)
        self.canv.addOutlineEntry(text, key, level=level, closed=False)
        if level == 0:
            self.notify("TOCEntry", (0, text, self.page - 1, key))


def p(text: str, style_map: dict[str, ParagraphStyle], style: str = "Body") -> Paragraph:
    return Paragraph(text, style_map[style])


def bullet(text: str, style_map: dict[str, ParagraphStyle]) -> Paragraph:
    return Paragraph(text, style_map["Bullet"], bulletText="•")


def table(
    rows: list[list[object]],
    style_map: dict[str, ParagraphStyle],
    widths: list[float] | None = None,
    font_size: float = 7.2,
) -> Table:
    converted = []
    for row_number, row in enumerate(rows):
        style = style_map["TableHeader"] if row_number == 0 else style_map["Table"]
        converted.append([Paragraph(str(cell), style) for cell in row])
    result = Table(converted, colWidths=widths, repeatRows=1, hAlign="LEFT")
    result.setStyle(TableStyle([
        ("BACKGROUND", (0, 0), (-1, 0), NAVY),
        ("TEXTCOLOR", (0, 0), (-1, 0), colors.white),
        ("FONTNAME", (0, 0), (-1, -1), "CJK"),
        ("FONTSIZE", (0, 0), (-1, -1), font_size),
        ("VALIGN", (0, 0), (-1, -1), "MIDDLE"),
        ("GRID", (0, 0), (-1, -1), 0.35, GRID),
        ("ROWBACKGROUNDS", (0, 1), (-1, -1), [colors.white, LIGHT]),
        ("LEFTPADDING", (0, 0), (-1, -1), 5),
        ("RIGHTPADDING", (0, 0), (-1, -1), 5),
        ("TOPPADDING", (0, 0), (-1, -1), 4),
        ("BOTTOMPADDING", (0, 0), (-1, -1), 4),
    ]))
    return result


def report_image(path: Path, width: float = 170 * mm, max_height: float = 205 * mm) -> Image:
    with PILImage.open(path) as image:
        pixel_width, pixel_height = image.size
    height = width * pixel_height / pixel_width
    if height > max_height:
        height = max_height
        width = height * pixel_width / pixel_height
    result = Image(str(path), width=width, height=height)
    result.hAlign = "CENTER"
    return result


def markdown_table(headers: list[str], rows: list[list[str]]) -> str:
    lines = [
        "| " + " | ".join(headers) + " |",
        "|" + "|".join("---" for _ in headers) + "|",
    ]
    lines.extend("| " + " | ".join(row) + " |" for row in rows)
    return "\n".join(lines)


def load_data() -> dict[str, object]:
    return {
        "factors": pd.read_csv(ROOT / "evaluation" / "factor_summary.csv"),
        "turnover": pd.read_csv(ROOT / "backtest" / "turnover_control_comparison.csv"),
        "sensitivity": pd.read_csv(ROOT / "backtest" / "cost_sensitivity.csv"),
        "ablation": pd.read_csv(ROOT / "backtest" / "factor_ablation.csv"),
        "lstm": pd.read_csv(ROOT / "lstm" / "model_comparison.csv"),
        "validation": json.loads(
            (ROOT / "reports" / "validation.json").read_text(encoding="utf-8")
        ),
        "neutralization": json.loads(
            (ROOT / "evaluation" / "neutralization_status.json").read_text(encoding="utf-8")
        ),
    }


def build_markdown(data: dict[str, object]) -> Path:
    factors = data["factors"]
    turnover = data["turnover"]
    lstm = data["lstm"]
    controlled = turnover.loc[turnover["turnover_policy"].eq("five_day_buffer20")]
    factor_rows = [
        [
            row.factor,
            f"{row.ic:.4f}",
            f"{row.rank_ic:.4f}",
            pct(row.long_short_total_return),
        ]
        for row in factors.itertuples(index=False)
    ]
    turnover_rows = [
        [
            row.mode,
            pct(row.avg_one_way_turnover),
            f"{row.total_cost:,.0f}",
            pct(row.total_return),
            pct(row.active_total_return),
            pct(row.annual_excess_return),
        ]
        for row in controlled.itertuples(index=False)
    ]
    lstm_rows = [
        [
            row.model,
            pct(row.accuracy),
            pct(row.balanced_accuracy),
            pct(row.f1),
            f"{row.auc:.4f}",
            f"{row.pr_auc:.4f}",
        ]
        for row in lstm.itertuples(index=False)
    ]
    content = f"""# 逐笔成交量化作业最终图文报告

## 摘要

项目覆盖 300 只股票、302 个交易日和 11 个基础字段，完成逐笔数据降采样、四因子研究、执行约束回测、换手控制和多股票 walk-forward LSTM。最终采用每 5 日调仓加 Top20 缓冲区，次日开盘口径成本后净收益为 15.83%，累计主动收益为 12.10%，年化超额收益为 11.00%。

## 数据处理

- 日频：11 张 302 x 300 宽表。
- 分钟频：3,322 张 253 x 300 表，包含 15:00。
- 日频 K 线长表：90,600 行。
- 复权价格：`Price / 100 * adjfactor`；成交额按未复权真实价格计算。

## 因子评价

{markdown_table(["因子", "IC", "Rank IC", "Q5-Q1累计收益"], factor_rows)}

![因子相关性](../evaluation/figures/factor_correlation_heatmap.png)

![成交额因子诊断](../evaluation/figures/amount_mean_sd_log_diagnostics.png)

![反转因子诊断](../evaluation/figures/reversal_5d_diagnostics.png)

## 回测与换手控制

信号在 t 日形成，在 t+1 日开盘或收盘成交。停牌不可交易，一字涨停不可买入，一字跌停不可卖出。成本包括卖出万分之五、双边 5bp 基础滑点和平方根冲击成本。

最终策略每 5 日调仓，保留仍在前 20 名的原持仓，再补足 10 只股票。

{markdown_table(["成交口径", "日均单边换手", "总成本", "净收益", "累计主动收益", "年化超额"], turnover_rows)}

![换手控制比较](../backtest/figures/turnover_control_comparison.png)

![换手控制净值](../backtest/figures/turnover_control_nav.png)

## LSTM样本外验证

LSTM 使用 6 只股票、140 个交易日和 3 折 expanding-window walk-forward，共产生 83,880 条样本外预测。

{markdown_table(["模型", "Accuracy", "Balanced Accuracy", "F1", "ROC AUC", "PR AUC"], lstm_rows)}

![模型比较](../lstm/figures/model_comparison.png)

![混淆矩阵](../lstm/figures/lstm_confusion_matrix.png)

## 结论

四因子组合具有成本前选股能力。每日换仓的高换手会使冲击成本吞噬信号收益；每 5 日调仓加 Top20 缓冲区把日均单边换手降至约 12.5%，并使次日开盘执行获得正净收益和正超额收益。LSTM 的排序指标高于逻辑回归，但上涨召回率仍低，需要进一步进行阈值和交易层检验。
"""
    path = REPORTS / "量化作业最终图文报告.md"
    path.write_text(content, encoding="utf-8")
    return path


def build_pdf(data: dict[str, object]) -> Path:
    register_font()
    style_map = styles()
    PDF_OUTPUT.mkdir(parents=True, exist_ok=True)
    output = PDF_OUTPUT / "量化作业最终图文报告.pdf"
    document = ReportTemplate(str(output), style_map)
    factors: pd.DataFrame = data["factors"]
    turnover: pd.DataFrame = data["turnover"]
    sensitivity: pd.DataFrame = data["sensitivity"]
    ablation: pd.DataFrame = data["ablation"]
    lstm: pd.DataFrame = data["lstm"]
    validation: dict = data["validation"]

    controlled = turnover.loc[turnover["turnover_policy"].eq("five_day_buffer20")]
    close_row = controlled.loc[controlled["mode"].eq("next_close_to_close")].iloc[0]
    open_row = controlled.loc[controlled["mode"].eq("next_open_to_open")].iloc[0]

    story = [
        Spacer(1, 34 * mm),
        p("逐笔成交量化作业", style_map, "Title"),
        p("最终图文研究报告", style_map, "Subtitle"),
        Spacer(1, 9 * mm),
        HRFlowable(width="62%", thickness=1.2, color=TEAL, hAlign="CENTER"),
        Spacer(1, 10 * mm),
        p("逐笔成交数据处理 · 四因子评价 · 换手控制回测 · Walk-forward LSTM", style_map, "Cover"),
        Spacer(1, 22 * mm),
        table(
            [
                ["研究范围", "核心产物", "验证状态"],
                ["300只股票 / 302日", "日频、分钟频、因子、回测、LSTM", "PASS"],
                ["90,600行日K", "3,322张分钟表", "12项单元测试通过"],
            ],
            style_map,
            widths=[52 * mm, 66 * mm, 50 * mm],
        ),
        Spacer(1, 30 * mm),
        p("生成日期：2026-07-30", style_map, "Cover"),
        PageBreak(),
        p("目录", style_map, "TOC"),
    ]
    toc = TableOfContents()
    toc.levelStyles = [ParagraphStyle(
        "TOCLevel1",
        fontName="CJK",
        fontSize=10,
        leading=17,
        leftIndent=0,
        firstLineIndent=0,
        textColor=INK,
        spaceBefore=4,
    )]
    story += [toc, PageBreak()]

    story += [
        p("摘要与核心结论", style_map, "H1"),
        p(
            "本项目建立了从逐笔成交数据到可执行策略与机器学习验证的完整链路。"
            "研究重点不是单独追求高回测收益，而是依次检查数据质量、前视偏差、"
            "交易限制、基准、换手、成本和严格样本外表现。",
            style_map,
        ),
        p(
            "正式推荐策略为每 5 个交易日调仓一次，并设置 Top20 排名缓冲区："
            "原持仓只要仍位于前 20 名就继续保留，再按综合因子得分补足 10 只股票。",
            style_map,
            "Callout",
        ),
        table(
            [
                ["最终策略指标", "次日收盘成交", "次日开盘成交"],
                ["日均单边换手", pct(close_row.avg_one_way_turnover), pct(open_row.avg_one_way_turnover)],
                ["成本后净收益", pct(close_row.total_return), pct(open_row.total_return)],
                ["累计主动收益", pct(close_row.active_total_return), pct(open_row.active_total_return)],
                ["年化超额收益", pct(close_row.annual_excess_return), pct(open_row.annual_excess_return)],
                ["信息比率", f"{close_row.information_ratio:.3f}", f"{open_row.information_ratio:.3f}"],
                ["最大回撤", pct(close_row.max_drawdown), pct(open_row.max_drawdown)],
            ],
            style_map,
            widths=[60 * mm, 54 * mm, 54 * mm],
        ),
        Spacer(1, 7 * mm),
        bullet("次日开盘成交取得 15.83% 净收益、12.10% 累计主动收益和 11.00% 年化超额收益。", style_map),
        bullet("次日收盘成交取得 1.37% 净收益，但累计主动收益为 -2.26%，仍略低于等权基准。", style_map),
        bullet("LSTM 的 ROC AUC 为 0.6465，高于逻辑回归的 0.6086，但上涨召回率仍只有 8.10%。", style_map),
        PageBreak(),
    ]

    story += [
        p("一、数据处理与质量控制", style_map, "H1"),
        p(
            "数据覆盖 2025 年 4 月 1 日至 2026 年 6 月 30 日，共 300 只混淆代码股票和 302 个交易日。"
            "程序从逐笔成交记录提取 OHLC、成交量、成交笔数、成交额及主买主卖量额，共 11 个字段。",
            style_map,
        ),
        table(
            [
                ["数据层", "规模", "关键处理"],
                ["日频宽表", "11张，每张302 x 300", "无成交日使用前一日收盘价"],
                ["分钟宽表", "3,322张，每张253 x 300", "无成交分钟使用此前收盘价"],
                ["日K长表", "90,600行", "date-code唯一且代码保留前导零"],
                ["复权价格", "Price / 100 x adjfactor", "成交额仍按实际未复权价格"],
            ],
            style_map,
            widths=[38 * mm, 54 * mm, 76 * mm],
        ),
        Spacer(1, 7 * mm),
        p("质量控制", style_map, "H2"),
        bullet("分钟索引完整包含集合竞价、连续竞价和 15:00，共 253 个时间点。", style_map),
        bullet("无成交时只使用已经发生的价格向后填充，不用未来价格填过去。", style_map),
        bullet("全量审计检查 OHLC 关系、有限值、宽表对齐、无成交填充和长表唯一性。", style_map),
        p(
            "行业和市值中性化接口已经实现，但题目数据没有真实行业、市值映射。"
            "本报告不使用代码前缀或流动性变量伪造风险暴露。",
            style_map,
            "Warning",
        ),
        PageBreak(),
    ]

    factor_rows = [["因子", "IC", "ICIR", "Rank IC", "Rank ICIR", "Q5-Q1累计收益"]]
    for row in factors.itertuples(index=False):
        factor_rows.append([
            row.factor,
            f"{row.ic:.4f}",
            f"{row.icir:.3f}",
            f"{row.rank_ic:.4f}",
            f"{row.rank_icir:.3f}",
            pct(row.long_short_total_return),
        ])
    story += [
        p("二、因子构建与评价", style_map, "H1"),
        p(
            "项目构建成交额多窗口均值离散度、10 日主动买卖不平衡冲击、10 日日内振幅和 5 日反转四个因子。"
            "每日横截面先做 1% 和 99% 缩尾，再标准化，以下一交易日收盘到收盘收益为评价目标。",
            style_map,
        ),
        table(factor_rows, style_map, widths=[43 * mm, 22 * mm, 24 * mm, 25 * mm, 27 * mm, 31 * mm], font_size=6.8),
        Spacer(1, 6 * mm),
        p(
            "成交额离散度和日内振幅表现为反向因子；5 日反转因子为正向因子。"
            "主动买卖不平衡冲击呈稳定负向预测，反映订单流异常后的短期反转。",
            style_map,
            "Callout",
        ),
        report_image(ROOT / "evaluation" / "figures" / "factor_correlation_heatmap.png", width=126 * mm, max_height=120 * mm),
        p("图 1：因子相关性热力图。最大绝对相关性约为 0.46，四个因子没有完全重复。", style_map, "Caption"),
        PageBreak(),
        p("因子分组与时序稳定性", style_map, "H2"),
        report_image(ROOT / "evaluation" / "figures" / "amount_mean_sd_log_diagnostics.png", width=170 * mm, max_height=65 * mm),
        p("图 2：成交额离散度因子。分组收益从 Q1 到 Q5 明显下降，滚动 Rank IC 主要为负。", style_map, "Caption"),
        report_image(ROOT / "evaluation" / "figures" / "intraday_range_10d_diagnostics.png", width=170 * mm, max_height=65 * mm),
        p("图 3：日内振幅因子。整体呈反向关系，但滚动 Rank IC 随时间波动。", style_map, "Caption"),
        PageBreak(),
        report_image(ROOT / "evaluation" / "figures" / "reversal_5d_diagnostics.png", width=170 * mm, max_height=65 * mm),
        p("图 4：5 日反转因子。Q5 相对 Q1 表现更强，平均 Rank IC 为正。", style_map, "Caption"),
        report_image(ROOT / "evaluation" / "figures" / "buy_sell_imbalance_surprise_10d_diagnostics.png", width=170 * mm, max_height=65 * mm),
        p("图 5：主动买卖不平衡冲击因子。负 Rank IC 表明异常主动买入后更容易出现短期反转。", style_map, "Caption"),
        PageBreak(),
    ]

    base = turnover.loc[turnover["turnover_policy"].eq("daily_top10")]
    base_close = base.loc[base["mode"].eq("next_close_to_close")].iloc[0]
    base_open = base.loc[base["mode"].eq("next_open_to_open")].iloc[0]
    story += [
        p("三、基准、成本与策略诊断", style_map, "H1"),
        p(
            "组合使用截至 t-1 日的 20 日历史 IC 均值确定因子方向和权重，t 日形成信号，"
            "并在 t+1 日开盘或收盘执行。基准为入场日可以买入股票的每日等权收益。",
            style_map,
        ),
        table(
            [
                ["每日Top10指标", "次日收盘成交", "次日开盘成交"],
                ["成本前收益", pct(base_close.gross_total_return), pct(base_open.gross_total_return)],
                ["基准收益", pct(base_close.benchmark_total_return), pct(base_open.benchmark_total_return)],
                ["日均单边换手", pct(base_close.avg_one_way_turnover), pct(base_open.avg_one_way_turnover)],
                ["总成本", f"{base_close.total_cost:,.0f}", f"{base_open.total_cost:,.0f}"],
                ["成本后净收益", pct(base_close.total_return), pct(base_open.total_return)],
            ],
            style_map,
            widths=[58 * mm, 55 * mm, 55 * mm],
        ),
        Spacer(1, 6 * mm),
        p(
            "每日 Top10 的成本前收益高于等权基准，说明综合信号并非完全无效。"
            "但超过 52% 的日均单边换手令基础滑点和冲击成本持续累积，成本后收益转负。",
            style_map,
            "Warning",
        ),
        report_image(ROOT / "backtest" / "figures" / "cost_and_return_breakdown.png", width=170 * mm, max_height=80 * mm),
        p("图 6：交易成本拆分与成本前、基准、成本后收益。冲击成本是主要成本项。", style_map, "Caption"),
        PageBreak(),
        p("成本敏感性", style_map, "H2"),
        p(
            "乐观情景仅保留卖出手续费；基准情景使用双边 5bp 滑点和平方根冲击；"
            "悲观情景使用双边 10bp 滑点，并把冲击参数提高 50%。",
            style_map,
        ),
        report_image(ROOT / "backtest" / "figures" / "cost_sensitivity.png", width=160 * mm, max_height=105 * mm),
        p("图 7：每日 Top10 的成本敏感性。策略经济价值对交易成本高度敏感。", style_map, "Caption"),
        p("因子消融", style_map, "H2"),
        report_image(ROOT / "backtest" / "figures" / "factor_ablation.png", width=160 * mm, max_height=90 * mm),
        p("图 8：一至四因子的逐步组合。增加因子没有单调改善成本后收益。", style_map, "Caption"),
        PageBreak(),
    ]

    turnover_rows = [["政策", "口径", "日均单边换手", "总成本", "净收益", "累计主动收益", "年化超额"]]
    for row in turnover.itertuples(index=False):
        turnover_rows.append([
            row.turnover_policy,
            "收盘" if row.mode == "next_close_to_close" else "开盘",
            pct(row.avg_one_way_turnover),
            f"{row.total_cost / 1_000_000:.2f}百万",
            pct(row.total_return),
            pct(row.active_total_return),
            pct(row.annual_excess_return),
        ])
    story += [
        p("四、换手控制与最终策略", style_map, "H1"),
        p(
            "换手控制分别检验调仓间隔与排名缓冲区。5 日调仓降低交易频率；Top20 缓冲区保留仍有较高排名的原持仓，"
            "避免边界名次的小幅波动触发无意义交易。",
            style_map,
        ),
        table(turnover_rows, style_map, widths=[32 * mm, 18 * mm, 27 * mm, 24 * mm, 22 * mm, 25 * mm, 24 * mm], font_size=6.4),
        Spacer(1, 6 * mm),
        report_image(ROOT / "backtest" / "figures" / "turnover_control_comparison.png", width=170 * mm, max_height=78 * mm),
        p("图 9：四种政策的日均单边换手与成本后总收益。5 日缓冲区同时取得最低换手和最高净收益。", style_map, "Caption"),
        PageBreak(),
        p("最终净值表现", style_map, "H2"),
        report_image(ROOT / "backtest" / "figures" / "turnover_control_nav.png", width=158 * mm, max_height=168 * mm),
        p("图 10：每日 Top10、5 日缓冲区和等权基准净值。次日开盘口径的最终策略跑赢基准。", style_map, "Caption"),
        p(
            "换手控制把日均单边换手下降约 76%，总成本下降约三分之二。"
            "次日开盘口径获得正净收益和正超额收益；次日收盘口径接近盈亏平衡但仍略低于基准。",
            style_map,
            "Callout",
        ),
        PageBreak(),
    ]

    lstm_rows = [["模型", "Accuracy", "Balanced Acc.", "Precision", "Recall", "F1", "ROC AUC", "PR AUC"]]
    for row in lstm.itertuples(index=False):
        lstm_rows.append([
            row.model,
            pct(row.accuracy),
            pct(row.balanced_accuracy),
            pct(row.precision),
            pct(row.recall),
            pct(row.f1),
            f"{row.auc:.4f}",
            f"{row.pr_auc:.4f}",
        ])
    story += [
        p("五、多股票 Walk-forward LSTM", style_map, "H1"),
        p(
            "分钟模型使用 6 只股票、140 个交易日和过去 20 分钟的五个特征预测下一分钟涨跌。"
            "验证采用 3 折 expanding-window walk-forward，每折只用训练期估计标准化参数，验证期选模，测试窗口互不重叠。",
            style_map,
        ),
        table(lstm_rows, style_map, widths=[30 * mm, 20 * mm, 23 * mm, 20 * mm, 18 * mm, 17 * mm, 20 * mm, 20 * mm], font_size=6.3),
        Spacer(1, 6 * mm),
        report_image(ROOT / "lstm" / "figures" / "model_comparison.png", width=160 * mm, max_height=92 * mm),
        p("图 11：83,880 条样本外预测的模型比较。LSTM 的 ROC AUC、PR AUC 和 F1 高于逻辑回归。", style_map, "Caption"),
        PageBreak(),
        p("分类误差结构", style_map, "H2"),
        report_image(ROOT / "lstm" / "figures" / "lstm_confusion_matrix.png", width=120 * mm, max_height=115 * mm),
        p("图 12：LSTM 样本外混淆矩阵。真正例 2,113，假负例 23,969。", style_map, "Caption"),
        report_image(ROOT / "lstm" / "figures" / "fold_model_metrics.png", width=150 * mm, max_height=72 * mm),
        p("图 13：各测试折的 ROC AUC 和 PR AUC。所有折均保持严格时间顺序。", style_map, "Caption"),
        p(
            "LSTM 的排序能力优于简单模型，但 0.5 阈值下上涨召回率只有 8.10%。"
            "模型概率还需要通过阈值选择、概率分组收益和交易成本进行经济价值验证。",
            style_map,
            "Warning",
        ),
        PageBreak(),
    ]

    story += [
        p("六、验证、局限与结论", style_map, "H1"),
        p("自动验证结果", style_map, "H2"),
        table(
            [
                ["检查项", "结果"],
                ["全量验证状态", validation["status"]],
                ["日频字段", str(validation["daily_fields"])],
                ["分钟表", f"{validation['minute_files']:,}"],
                ["回测交易审计", f"{validation['backtest_trade_rows']:,}条"],
                ["换手控制变体", str(validation["turnover_control_variants"])],
                ["LSTM样本外样本", f"{validation['lstm_oos_samples']:,}"],
                ["单元测试", "12项全部通过"],
            ],
            style_map,
            widths=[82 * mm, 82 * mm],
        ),
        Spacer(1, 8 * mm),
        p("研究结论", style_map, "H2"),
        bullet("数据处理、复权、缺失填充和输出结构均通过全量审计。", style_map),
        bullet("四因子组合存在成本前选股信息，但每日调仓无法覆盖换手和市场冲击。", style_map),
        bullet("5 日调仓加 Top20 缓冲区显著降低换手；次日开盘执行取得正净收益和正超额收益。", style_map),
        bullet("次日收盘执行仍略低于等权基准，策略结果对成交时点保持敏感。", style_map),
        bullet("LSTM 在排序指标上优于逻辑回归，但上涨召回率较低，尚不能直接等同于交易收益。", style_map),
        p("主要局限", style_map, "H2"),
        bullet("股票代码混淆且没有真实行业、市值数据，风险中性化接口未在正式结果中启用。", style_map),
        bullet("换手政策只比较预设的四种规则，仍需在更长样本和不同市场阶段验证稳定性。", style_map),
        bullet("LSTM 只覆盖 6 只股票和 140 日，下一步应扩展股票数并把概率转成含成本的交易规则。", style_map),
        p(
            "综合而言，本项目的核心成果不是得到一条无条件成立的高收益曲线，而是形成了一套能够识别前视偏差、"
            "拆解交易成本、控制换手并进行严格样本外验证的可复现研究流程。",
            style_map,
            "Callout",
        ),
        Spacer(1, 8 * mm),
        p("主要输出文件", style_map, "H2"),
        p(
            "回测汇总：backtest/turnover_control_comparison.csv<br/>"
            "换手控制净值：backtest/turnover_control_results.csv<br/>"
            "模型比较：lstm/model_comparison.csv<br/>"
            "全量审计：reports/validation.json<br/>"
            "复现入口：scripts/analyze_factors_backtest.py、scripts/train_lstm_pytorch.py",
            style_map,
            "Small",
        ),
    ]

    document.multiBuild(story)
    return output


def main() -> None:
    data = load_data()
    markdown = build_markdown(data)
    pdf = build_pdf(data)
    print(markdown)
    print(pdf)


if __name__ == "__main__":
    main()
\end{ProjectCode}
\clearpage
\section{14. \texttt{scripts/build\_factor\_optimization\_report.py}}
生成七因子扩展与收益优化图文报告。

\textbf{物理行数：} 513\\
\textbf{SHA-256：} \texttt{6b59696b7b7a753013ee7774eed5dd5acef990db48ac5127c31460762daf7529}

\begin{ProjectCode}
#!/usr/bin/env python3
"""Build the illustrated seven-factor and turnover-optimized assignment report."""

from __future__ import annotations

import json
from pathlib import Path

import pandas as pd
from reportlab.lib.units import mm
from reportlab.platypus import HRFlowable, PageBreak, Spacer
from reportlab.platypus.tableofcontents import TableOfContents
from reportlab.lib.styles import ParagraphStyle

from build_final_illustrated_report import (
    INK,
    PDF_OUTPUT,
    REPORTS,
    ROOT,
    TEAL,
    ReportTemplate,
    bullet,
    markdown_table,
    p,
    pct,
    register_font,
    report_image,
    styles,
    table,
)


NEW_FACTORS = [
    "avg_trade_size_surprise_20d",
    "price_volume_pressure_10d",
    "gap_intraday_divergence_5d",
]
RECOMMENDED_POLICY = "five_day_buffer60_top30"


def load_data() -> dict[str, object]:
    return {
        "metadata": pd.read_csv(ROOT / "factors" / "factor_metadata.csv"),
        "factor_summary": pd.read_csv(ROOT / "evaluation" / "factor_summary.csv"),
        "optimized": pd.read_csv(ROOT / "backtest" / "optimized_strategy_metrics.csv"),
        "sensitivity": pd.read_csv(ROOT / "backtest" / "cost_sensitivity.csv"),
        "ablation": pd.read_csv(ROOT / "backtest" / "factor_ablation.csv"),
        "lstm": pd.read_csv(ROOT / "lstm" / "model_comparison.csv"),
        "validation": json.loads(
            (ROOT / "reports" / "validation.json").read_text(encoding="utf-8")
        ),
    }


def row_for(frame: pd.DataFrame, sample: str, mode: str) -> pd.Series:
    return frame.loc[frame["sample"].eq(sample) & frame["mode"].eq(mode)].iloc[0]


def new_factor_rows(data: dict[str, object]) -> list[list[str]]:
    metadata: pd.DataFrame = data["metadata"]
    summary: pd.DataFrame = data["factor_summary"]
    merged = metadata.merge(summary, on="factor", how="left")
    merged = merged.loc[merged["factor"].isin(NEW_FACTORS)].set_index("factor").loc[NEW_FACTORS]
    return [
        [
            row.name_zh,
            row.formula,
            f"{row.ic:.4f}",
            f"{row.rank_ic:.4f}",
            f"{row.rank_icir:.3f}",
            pct(row.long_short_total_return),
        ]
        for row in merged.itertuples(index=False)
    ]


def all_factor_rows(data: dict[str, object]) -> list[list[str]]:
    metadata: pd.DataFrame = data["metadata"]
    summary: pd.DataFrame = data["factor_summary"]
    merged = metadata.merge(summary, on="factor", how="left")
    return [
        [
            row.name_zh,
            row.kind,
            f"{row.ic:.4f}",
            f"{row.rank_ic:.4f}",
            f"{row.rank_icir:.3f}",
            pct(row.long_short_total_return),
        ]
        for row in merged.itertuples(index=False)
    ]


def build_markdown(data: dict[str, object]) -> Path:
    optimized: pd.DataFrame = data["optimized"]
    full_open = row_for(optimized, "full", "next_open_to_open")
    full_close = row_for(optimized, "full", "next_close_to_close")
    holdout = row_for(optimized, "holdout20", "next_open_to_open")
    factor_rows = all_factor_rows(data)
    content = f"""# 七因子扩展与收益优化图文报告

## 结论摘要

正式策略使用七个因子、Top30 持仓、每 5 日调仓和 Top60 持有缓冲。信号在 t 日形成，在 t+1 日成交；因子权重只使用截至 t-1 的 20 日历史 IC。次日开盘口径在手续费、双边滑点、平方根冲击成本、停牌和涨跌停限制全部生效后，累计收益为 **{pct(full_open.total_return)}**，年化收益为 **{pct(full_open.annual_return)}**，最大回撤为 **{pct(full_open.max_drawdown)}**，夏普率为 **{full_open.sharpe:.3f}**。

完整样本达到 25% 目标，但最后 20% 留出期的绝对收益为 {pct(holdout.total_return)}，同期基准为 {pct(holdout.benchmark_total_return)}，主动收益为 {pct(holdout.active_total_return)}。因此 25% 是历史完整样本结果，不是未来保证。

## 七个因子的评价

{markdown_table(["因子", "类别", "IC", "Rank IC", "Rank ICIR", "Q5-Q1累计收益"], factor_rows)}

三个新增因子是 20 日单笔成交规模异常、10 日价量压力、5 日隔夜与日内收益背离。它们不复制参考 PDF 或汇报 PPT 已列出的因子。

![因子相关性](../evaluation/figures/factor_correlation_heatmap.png)

## 完整样本风险收益

| 指标 | 次日收盘 | 次日开盘 |
|---|---:|---:|
| 成本前累计收益 | {pct(full_close.gross_total_return)} | {pct(full_open.gross_total_return)} |
| 成本后累计收益 | {pct(full_close.total_return)} | {pct(full_open.total_return)} |
| 年化收益 | {pct(full_close.annual_return)} | {pct(full_open.annual_return)} |
| 年化超额收益 | {pct(full_close.annual_excess_return)} | {pct(full_open.annual_excess_return)} |
| 年化波动率 | {pct(full_close.annual_volatility)} | {pct(full_open.annual_volatility)} |
| 最大回撤 | {pct(full_close.max_drawdown)} | {pct(full_open.max_drawdown)} |
| 夏普率 | {full_close.sharpe:.3f} | {full_open.sharpe:.3f} |
| 信息比率 | {full_close.information_ratio:.3f} | {full_open.information_ratio:.3f} |
| 日均单边换手 | {pct(full_close.avg_one_way_turnover)} | {pct(full_open.avg_one_way_turnover)} |

![优化策略净值和回撤](../backtest/figures/optimized_strategy_nav_drawdown.png)

![研究期与留出期](../backtest/figures/optimized_period_split.png)

## 解释边界

- 25% 目标在完整样本次日开盘口径达到，且悲观成本情景的累计收益仍为 27.37%。
- 最后 20% 留出期跑赢下跌的等权基准，但绝对收益仍为负。
- 结果不能视为收益保证；下一步应使用独立的新时间区间继续前推检验。
"""
    path = REPORTS / "七因子扩展与收益优化图文报告.md"
    path.write_text(content, encoding="utf-8")
    return path


def build_pdf(data: dict[str, object]) -> Path:
    register_font()
    style_map = styles()
    PDF_OUTPUT.mkdir(parents=True, exist_ok=True)
    output = PDF_OUTPUT / "七因子扩展与收益优化图文报告.pdf"
    document = ReportTemplate(
        str(output),
        style_map,
        title="七因子扩展与收益优化图文报告",
        subject="三个新增因子、真实成本回测、年化收益、回撤与留出期验证",
        header="七因子扩展与收益优化图文报告",
    )

    metadata: pd.DataFrame = data["metadata"]
    summary: pd.DataFrame = data["factor_summary"]
    optimized: pd.DataFrame = data["optimized"]
    sensitivity: pd.DataFrame = data["sensitivity"]
    ablation: pd.DataFrame = data["ablation"]
    lstm: pd.DataFrame = data["lstm"]
    validation: dict = data["validation"]

    full_close = row_for(optimized, "full", "next_close_to_close")
    full_open = row_for(optimized, "full", "next_open_to_open")
    research = row_for(optimized, "research80", "next_open_to_open")
    holdout = row_for(optimized, "holdout20", "next_open_to_open")

    story = [
        Spacer(1, 33 * mm),
        p("七因子扩展与收益优化", style_map, "Title"),
        p("逐笔成交量化作业图文报告", style_map, "Subtitle"),
        Spacer(1, 9 * mm),
        HRFlowable(width="62%", thickness=1.2, color=TEAL, hAlign="CENTER"),
        Spacer(1, 10 * mm),
        p("三个新增因子 · 真实成本回测 · 年化回报与回撤 · 留出期披露", style_map, "Cover"),
        Spacer(1, 20 * mm),
        table(
            [
                ["研究范围", "正式策略", "验证状态"],
                ["300只股票 / 302日", "Top30 / 5日调仓 / Top60缓冲", validation["status"]],
                ["七个日频因子", "t信号 / t+1成交 / 历史IC", "13项测试通过"],
            ],
            style_map,
            widths=[52 * mm, 72 * mm, 44 * mm],
        ),
        Spacer(1, 26 * mm),
        p("生成日期：2026-08-02", style_map, "Cover"),
        PageBreak(),
        p("目录", style_map, "TOC"),
    ]
    toc = TableOfContents()
    toc.levelStyles = [
        ParagraphStyle(
            "TOCLevel1",
            fontName="CJK",
            fontSize=10,
            leading=17,
            leftIndent=0,
            firstLineIndent=0,
            textColor=INK,
            spaceBefore=4,
        )
    ]
    story += [toc, PageBreak()]

    story += [
        p("一、核心结论", style_map, "H1"),
        p(
            "正式策略使用七个因子、Top30 持仓、每 5 个交易日调仓和 Top60 持有缓冲。"
            "因子权重只使用截至 t-1 的 20 日历史 IC，信号在 t 日形成，在 t+1 日成交。",
            style_map,
        ),
        p(
            f"次日开盘口径在全部交易成本和约束生效后取得 {pct(full_open.total_return)} 累计收益、"
            f"{pct(full_open.annual_return)} 年化收益和 {pct(full_open.max_drawdown)} 最大回撤，"
            "完整样本达到 25% 的目标。",
            style_map,
            "Callout",
        ),
        table(
            [
                ["完整样本指标", "次日收盘成交", "次日开盘成交"],
                ["成本前累计收益", pct(full_close.gross_total_return), pct(full_open.gross_total_return)],
                ["成本后累计收益", pct(full_close.total_return), pct(full_open.total_return)],
                ["年化收益率", pct(full_close.annual_return), pct(full_open.annual_return)],
                ["基准累计收益", pct(full_close.benchmark_total_return), pct(full_open.benchmark_total_return)],
                ["累计主动收益", pct(full_close.active_total_return), pct(full_open.active_total_return)],
                ["年化超额收益", pct(full_close.annual_excess_return), pct(full_open.annual_excess_return)],
                ["年化波动率", pct(full_close.annual_volatility), pct(full_open.annual_volatility)],
                ["最大回撤", pct(full_close.max_drawdown), pct(full_open.max_drawdown)],
                ["夏普率", f"{full_close.sharpe:.3f}", f"{full_open.sharpe:.3f}"],
                ["信息比率", f"{full_close.information_ratio:.3f}", f"{full_open.information_ratio:.3f}"],
                ["日均单边换手", pct(full_close.avg_one_way_turnover), pct(full_open.avg_one_way_turnover)],
                ["总交易成本", f"{full_close.total_cost:,.0f}", f"{full_open.total_cost:,.0f}"],
            ],
            style_map,
            widths=[62 * mm, 53 * mm, 53 * mm],
        ),
        Spacer(1, 6 * mm),
        p(
            "收益提高来自降低换手和扩大组合容量，不来自同日未来 IC、删掉亏损日期或取消交易成本。"
            "但完整样本包含策略选择过程，因此仍需结合最后 20% 留出期判断稳健性。",
            style_map,
            "Warning",
        ),
        PageBreak(),
    ]

    new_merged = metadata.merge(summary, on="factor", how="left")
    new_merged = new_merged.loc[new_merged["factor"].isin(NEW_FACTORS)].set_index("factor").loc[NEW_FACTORS]
    new_rows = [["新增因子", "公式", "IC", "Rank IC", "Rank ICIR", "Q5-Q1累计收益"]]
    new_rows.extend(new_factor_rows(data))
    story += [
        p("二、三个新增因子", style_map, "H1"),
        p(
            "参考 PDF 和汇报 PPT 中已经出现成交额离散度、动量或反转、主买比例、振幅或收盘位置、"
            "以及 Amihud 非流动性等思路。本次新增因子改用成交笔数、价量联动和开盘—收盘分解，避免直接复制。",
            style_map,
        ),
        table(new_rows, style_map, widths=[30 * mm, 62 * mm, 16 * mm, 20 * mm, 21 * mm, 23 * mm], font_size=6.2),
        Spacer(1, 7 * mm),
        bullet("20 日单笔成交规模异常：衡量平均每笔成交金额相对自身历史的异常程度。", style_map),
        bullet("10 日价量压力：衡量收益方向是否获得成交量变化确认。", style_map),
        bullet("5 日隔夜—日内收益背离：衡量开盘定价与盘中价格修正之间的差异。", style_map),
        p(
            "所有滚动窗口只使用当日及更早数据；因子与下一交易日收益配对只发生在评价阶段。"
            "三个新增因子在综合信号中乘以 0.1 的保守缩放，避免新增变量因一次样本内尝试就主导组合。",
            style_map,
            "Callout",
        ),
        PageBreak(),
    ]

    for index, (factor, caption) in enumerate(
        [
            ("avg_trade_size_surprise_20d", "单笔成交规模异常的五分组、累计净值和滚动 Rank IC。"),
            ("price_volume_pressure_10d", "价量压力因子的分组收益与时序稳定性。"),
            ("gap_intraday_divergence_5d", "隔夜—日内收益背离因子，Rank IC 为正且多空累计收益较强。"),
        ],
        start=1,
    ):
        factor_name = new_merged.loc[factor, "name_zh"]
        story += [
            p(f"新增因子诊断：{factor_name}", style_map, "H2"),
            report_image(
                ROOT / "evaluation" / "figures" / f"{factor}_diagnostics.png",
                width=170 * mm,
                max_height=150 * mm,
            ),
            p(f"图 {index}：{caption}", style_map, "Caption"),
        ]
        if index < 3:
            story.append(PageBreak())
    story.append(PageBreak())

    all_rows = [["因子", "类别", "IC", "Rank IC", "Rank ICIR", "Q5-Q1累计收益"]]
    all_rows.extend(all_factor_rows(data))
    story += [
        p("三、七因子整体评价", style_map, "H1"),
        table(all_rows, style_map, widths=[42 * mm, 22 * mm, 20 * mm, 23 * mm, 25 * mm, 30 * mm], font_size=6.5),
        Spacer(1, 6 * mm),
        report_image(
            ROOT / "evaluation" / "figures" / "factor_correlation_heatmap.png",
            width=137 * mm,
            max_height=128 * mm,
        ),
        p("图 4：七因子横截面相关性。最大绝对相关性约为 0.54，没有完全重复的因子。", style_map, "Caption"),
        p(
            "因子评价使用下一交易日收盘到收盘收益，并在每日横截面做 1%/99% 缩尾和标准化。"
            "行业、市值中性化接口已经实现，但匿名代码无法匹配真实暴露，因此正式结果没有使用虚构代理变量。",
            style_map,
            "Warning",
        ),
        PageBreak(),
    ]

    story += [
        p("四、组合构建与交易约束", style_map, "H1"),
        p(
            "组合由七个标准化因子按历史 20 日 IC 加权。原有四个因子使用完整尺度，三个新增因子使用 0.1 倍尺度。"
            "最终持有综合排名前 30 的股票，每 5 日调仓；已有持仓只要仍在前 60 名便继续保留。",
            style_map,
        ),
        table(
            [
                ["模块", "正式设定"],
                ["信息时点", "t 日信号；仅使用截至 t-1 已实现的历史 IC"],
                ["成交时点", "t+1 日开盘或收盘，分别回测"],
                ["持仓与换手", "Top30；每5日调仓；Top60持有缓冲"],
                ["停牌", "成交量或成交额为零不可交易"],
                ["涨跌停", "一字涨停不能买；一字跌停不能卖"],
                ["成本", "卖出5bp；双边滑点各5bp；平方根冲击成本"],
                ["基准", "入场日可交易股票的每日等权收益"],
            ],
            style_map,
            widths=[43 * mm, 125 * mm],
        ),
        Spacer(1, 8 * mm),
        report_image(
            ROOT / "backtest" / "figures" / "turnover_control_comparison.png",
            width=170 * mm,
            max_height=90 * mm,
        ),
        p("图 5：不同换手政策的日均单边换手与成本后收益。Top30 的 5 日缓冲方案被选为正式策略。", style_map, "Caption"),
        PageBreak(),
    ]

    story += [
        p("五、收益、年化回报与回撤", style_map, "H1"),
        report_image(
            ROOT / "backtest" / "figures" / "optimized_strategy_nav_drawdown.png",
            width=170 * mm,
            max_height=128 * mm,
        ),
        p("图 6：正式策略与等权基准净值，以及两个成交口径的回撤曲线。", style_map, "Caption"),
        report_image(
            ROOT / "backtest" / "figures" / "optimized_strategy_risk_metrics.png",
            width=170 * mm,
            max_height=82 * mm,
        ),
        p("图 7：累计收益、年化收益、年化超额、波动率、最大回撤、夏普率和信息比率。", style_map, "Caption"),
        PageBreak(),
    ]

    open_sensitivity = sensitivity.loc[sensitivity["mode"].eq("next_open_to_open")]
    sensitivity_rows = [["成本情景", "累计收益", "年化收益", "最大回撤", "夏普率", "总成本"]]
    for row in open_sensitivity.itertuples(index=False):
        sensitivity_rows.append([
            row.scenario,
            pct(row.total_return),
            pct(row.annual_return),
            pct(row.max_drawdown),
            f"{row.sharpe:.3f}",
            f"{row.total_cost:,.0f}",
        ])
    story += [
        p("六、成本敏感性与因子消融", style_map, "H1"),
        p("次日开盘口径的三档成本结果如下。", style_map),
        table(sensitivity_rows, style_map, widths=[30 * mm, 27 * mm, 27 * mm, 27 * mm, 24 * mm, 33 * mm]),
        Spacer(1, 7 * mm),
        report_image(ROOT / "backtest" / "figures" / "cost_sensitivity.png", width=158 * mm, max_height=86 * mm),
        p("图 8：悲观成本情景的累计收益仍为 27.37%，但年化收益降至 24.72%。", style_map, "Caption"),
        report_image(ROOT / "backtest" / "figures" / "factor_ablation.png", width=158 * mm, max_height=86 * mm),
        p("图 9：从一个因子逐步增加至七个因子的执行约束回测。七因子组合改善了回撤和夏普率。", style_map, "Caption"),
        PageBreak(),
    ]

    story += [
        p("七、研究期与最后 20% 留出期", style_map, "H1"),
        table(
            [
                ["次日开盘口径", "研究期前80%", "留出期后20%"],
                ["日期", f"{research.start_date}—{research.end_date}", f"{holdout.start_date}—{holdout.end_date}"],
                ["成本后累计收益", pct(research.total_return), pct(holdout.total_return)],
                ["等权基准收益", pct(research.benchmark_total_return), pct(holdout.benchmark_total_return)],
                ["累计主动收益", pct(research.active_total_return), pct(holdout.active_total_return)],
                ["年化收益", pct(research.annual_return), pct(holdout.annual_return)],
                ["年化超额收益", pct(research.annual_excess_return), pct(holdout.annual_excess_return)],
                ["最大回撤", pct(research.max_drawdown), pct(holdout.max_drawdown)],
                ["夏普率", f"{research.sharpe:.3f}", f"{holdout.sharpe:.3f}"],
            ],
            style_map,
            widths=[58 * mm, 55 * mm, 55 * mm],
        ),
        Spacer(1, 8 * mm),
        report_image(ROOT / "backtest" / "figures" / "optimized_period_split.png", width=158 * mm, max_height=105 * mm),
        p("图 10：研究期与留出期的策略、基准和主动收益。", style_map, "Caption"),
        p(
            "留出期策略绝对收益为负，但跌幅小于股票池等权基准，因此主动收益仍为正。"
            "这意味着组合在该阶段有相对选股价值，却不能证明稳定的绝对盈利能力。",
            style_map,
            "Warning",
        ),
        PageBreak(),
    ]

    open_ablation = ablation.loc[ablation["mode"].eq("next_open_to_open")]
    ablation_rows = [["因子数", "组合", "累计收益", "年化收益", "最大回撤", "夏普率"]]
    for row in open_ablation.itertuples(index=False):
        ablation_rows.append([
            str(row.n_factors),
            row.ablation_step,
            pct(row.total_return),
            pct(row.annual_return),
            pct(row.max_drawdown),
            f"{row.sharpe:.3f}",
        ])
    story += [
        p("八、补充结果与验证", style_map, "H1"),
        p("次日开盘口径因子消融明细", style_map, "H2"),
        table(ablation_rows, style_map, widths=[15 * mm, 55 * mm, 27 * mm, 27 * mm, 27 * mm, 22 * mm], font_size=6.6),
        Spacer(1, 8 * mm),
        p("分钟级模型", style_map, "H2"),
    ]
    lstm_rows = [["模型", "Accuracy", "Balanced Acc.", "F1", "ROC AUC", "PR AUC"]]
    for row in lstm.itertuples(index=False):
        lstm_rows.append([
            row.model,
            pct(row.accuracy),
            pct(row.balanced_accuracy),
            pct(row.f1),
            f"{row.auc:.4f}",
            f"{row.pr_auc:.4f}",
        ])
    story += [
        table(lstm_rows, style_map, widths=[38 * mm, 26 * mm, 30 * mm, 24 * mm, 25 * mm, 25 * mm]),
        Spacer(1, 7 * mm),
        p(
            "LSTM 使用 3 折 expanding-window walk-forward，共 83,880 条样本外预测。"
            "其 ROC AUC 高于逻辑回归，但上涨召回率只有 8.10%，不能直接等同于可交易收益。",
            style_map,
        ),
        p("自动验证", style_map, "H2"),
        table(
            [
                ["检查项", "结果"],
                ["全量验证", validation["status"]],
                ["日频字段", str(validation["daily_fields"])],
                ["分钟表", f"{validation['minute_files']:,}"],
                ["回测交易审计", f"{validation['backtest_trade_rows']:,} 条"],
                ["换手政策变体", str(validation["turnover_control_variants"])],
                ["LSTM样本外预测", f"{validation['lstm_oos_samples']:,}"],
                ["单元测试", "13 项全部通过"],
            ],
            style_map,
            widths=[82 * mm, 82 * mm],
        ),
        PageBreak(),
    ]

    story += [
        p("九、结论与解释边界", style_map, "H1"),
        bullet("新增三个因子均使用与参考材料不同的变量组合，并完成 IC、Rank IC、分组收益和相关性评价。", style_map),
        bullet("Top30、5 日调仓和 Top60 持有缓冲将日均单边换手降到约 8.3%。", style_map),
        bullet("次日开盘口径完整样本成本后累计收益为 32.74%，年化收益为 29.51%，最大回撤为 -11.27%。", style_map),
        bullet("悲观成本情景累计收益为 27.37%，说明完整样本的 25% 结论并非只在零成本下成立。", style_map),
        bullet("最后 20% 留出期绝对收益为 -2.61%，但相对基准主动收益为 +5.41%。", style_map),
        p(
            "最重要的解释边界是：策略选择和参数比较使用了这段历史数据，因而完整样本 32.74% 仍带有研究选择效应。"
            "报告不把它表述为未来保证。更严格的下一步是锁定代码与参数，在完全未参与研究的新日期上进行 walk-forward 前推。",
            style_map,
            "Warning",
        ),
        p("主要复现文件", style_map, "H2"),
        p(
            "因子代码：pipeline_code/construct_factors.py<br/>"
            "回测与绘图：scripts/analyze_factors_backtest.py<br/>"
            "七因子指标：evaluation/factor_summary.csv<br/>"
            "优化策略指标：backtest/optimized_strategy_metrics.csv<br/>"
            "逐日净值：backtest/optimized_strategy_results.csv<br/>"
            "全量审计：reports/validation.json",
            style_map,
            "Small",
        ),
    ]

    document.multiBuild(story)
    return output


def main() -> None:
    data = load_data()
    markdown = build_markdown(data)
    pdf = build_pdf(data)
    print(markdown)
    print(pdf)


if __name__ == "__main__":
    main()
\end{ProjectCode}
\clearpage
\section{15. \texttt{scripts/build\_factor\_formula\_report.py}}
生成七因子公式与图像分析报告。

\textbf{物理行数：} 376\\
\textbf{SHA-256：} \texttt{25560e458d58e0d77ac4362098e2aa227dae70d703e90942b5251ece86e042b8}

\begin{ProjectCode}
#!/usr/bin/env python3
"""Build a factor-focused illustrated PDF with formulas and diagnostics."""

from __future__ import annotations

from pathlib import Path

import pandas as pd
from reportlab.lib.styles import ParagraphStyle
from reportlab.lib.units import mm
from reportlab.platypus import HRFlowable, PageBreak, Spacer
from reportlab.platypus.tableofcontents import TableOfContents

from build_final_illustrated_report import (
    INK,
    PDF_OUTPUT,
    ROOT,
    TEAL,
    ReportTemplate,
    bullet,
    p,
    pct,
    register_font,
    report_image,
    styles,
    table,
)


FACTOR_ORDER = [
    "amount_mean_sd_log",
    "buy_sell_imbalance_surprise_10d",
    "intraday_range_10d",
    "reversal_5d",
    "avg_trade_size_surprise_20d",
    "price_volume_pressure_10d",
    "gap_intraday_divergence_5d",
]

FACTOR_EXPLANATIONS = {
    "amount_mean_sd_log": {
        "meaning": "比较成交额在 1、5、10、20 日窗口上的均值差异，衡量成交活跃度在不同周期之间是否出现结构性分化。",
        "variables": "amount 为当日实际成交额；MA_k 表示过去 k 个交易日移动均值；Std 为四个窗口值的标准差。",
    },
    "buy_sell_imbalance_surprise_10d": {
        "meaning": "衡量当日主动买卖不平衡相对过去 10 日常态的冲击，用于识别短期订单流过度反应。",
        "variables": "主动不平衡为 (buy_volume-sell_volume)/(buy_volume+sell_volume)；历史均值只使用 t-1 及以前的数据。",
    },
    "intraday_range_10d": {
        "meaning": "用最高价与最低价之差除以收盘价，衡量日内波动幅度，再用 10 日均值降低单日噪声。",
        "variables": "high、low、close 分别为复权后的最高价、最低价和收盘价。",
    },
    "reversal_5d": {
        "meaning": "对过去 5 日累计收益取负，检验短期上涨后回落、短期下跌后修复的反转效应。",
        "variables": "close_t 为当日复权收盘价；close_t-5 为 5 个交易日前复权收盘价。",
    },
    "avg_trade_size_surprise_20d": {
        "meaning": "用成交额除以成交笔数估计平均单笔规模，再计算其相对过去 20 日的异常程度。",
        "variables": "trade_count 为成交笔数；zscore_20d 使用过去 20 日均值和标准差，不使用未来信息。",
    },
    "price_volume_pressure_10d": {
        "meaning": "把当日收益与成交量对数变化相乘，判断价格运动是否获得成交量变化确认，再取 10 日均值。",
        "variables": "return_1d 为当日收盘收益；diff(log(1+volume)) 为成交量对数的一阶变化。",
    },
    "gap_intraday_divergence_5d": {
        "meaning": "比较隔夜收益和日内收益，刻画开盘定价与盘中价格修正方向是否出现持续背离。",
        "variables": "隔夜收益为 open_t/close_t-1-1；日内收益为 close_t/open_t-1；最后取 5 日均值。",
    },
}


def factor_direction(rank_ic: float) -> str:
    if abs(rank_ic) < 0.01:
        return "近似中性"
    return "正向" if rank_ic > 0 else "反向"


def load_data() -> tuple[pd.DataFrame, pd.DataFrame]:
    metadata = pd.read_csv(ROOT / "factors" / "factor_metadata.csv")
    summary = pd.read_csv(ROOT / "evaluation" / "factor_summary.csv")
    factors = metadata.merge(summary, on="factor", how="inner")
    factors = factors.set_index("factor").loc[FACTOR_ORDER].reset_index()
    metrics = pd.read_csv(ROOT / "backtest" / "optimized_strategy_metrics.csv")
    return factors, metrics


def formula_display(factor: str, fallback: str) -> str:
    formulas = {
        "amount_mean_sd_log": "F_t = log(1 + Std(MA_1(A_t), MA_5(A_t), MA_10(A_t), MA_20(A_t)))",
        "buy_sell_imbalance_surprise_10d": "F_t = OI_t - MA_10(OI_(t-1), ..., OI_(t-10)), OI_t=(BuyVolume_t-SellVolume_t)/(BuyVolume_t+SellVolume_t)",
        "intraday_range_10d": "F_t = MA_10((High_t - Low_t) / Close_t)",
        "reversal_5d": "F_t = -(Close_t / Close_(t-5) - 1)",
        "avg_trade_size_surprise_20d": "F_t = ZScore_20(log(1 + Amount_t / max(TradeCount_t, 1)))",
        "price_volume_pressure_10d": "F_t = MA_10(Return_t * diff(log(1 + Volume_t)))",
        "gap_intraday_divergence_5d": "F_t = MA_5[(Open_t / Close_(t-1) - 1) - (Close_t / Open_t - 1)]",
    }
    return formulas.get(factor, fallback)


def build_pdf() -> Path:
    register_font()
    style_map = styles()
    factors, metrics = load_data()
    PDF_OUTPUT.mkdir(parents=True, exist_ok=True)
    output = PDF_OUTPUT / "七因子公式与图像分析报告.pdf"
    document = ReportTemplate(
        str(output),
        style_map,
        title="七因子公式与图像分析报告",
        subject="七个日频因子的公式、解释、评价指标、分组收益和诊断图",
        header="七因子公式与图像分析报告",
    )

    full_close = metrics.loc[
        metrics["sample"].eq("full") & metrics["mode"].eq("next_close_to_close")
    ].iloc[0]
    full_open = metrics.loc[
        metrics["sample"].eq("full") & metrics["mode"].eq("next_open_to_open")
    ].iloc[0]
    holdout_open = metrics.loc[
        metrics["sample"].eq("holdout20") & metrics["mode"].eq("next_open_to_open")
    ].iloc[0]

    story = [
        Spacer(1, 34 * mm),
        p("七因子公式与图像分析", style_map, "Title"),
        p("逐笔成交量化作业专题报告", style_map, "Subtitle"),
        Spacer(1, 9 * mm),
        HRFlowable(width="62%", thickness=1.2, color=TEAL, hAlign="CENTER"),
        Spacer(1, 10 * mm),
        p("因子定义 · 数学公式 · IC/Rank IC · 五分组 · 净值与回撤", style_map, "Cover"),
        Spacer(1, 22 * mm),
        table(
            [
                ["研究对象", "因子数量", "评价目标"],
                ["300只股票 / 302个交易日", "7个日频因子", "下一交易日收盘到收盘收益"],
                ["横截面缩尾与标准化", "IC及Rank IC体系", "五分组与Q5-Q1多空收益"],
            ],
            style_map,
            widths=[58 * mm, 52 * mm, 58 * mm],
        ),
        Spacer(1, 27 * mm),
        p("版式参考：用户提供的 A-2.pdf；图像与数值均由本项目数据重新生成。", style_map, "Cover"),
        p("生成日期：2026-08-03", style_map, "Cover"),
        PageBreak(),
        p("目录", style_map, "TOC"),
    ]
    toc = TableOfContents()
    toc.levelStyles = [
        ParagraphStyle(
            "TOCLevel1",
            fontName="CJK",
            fontSize=10,
            leading=17,
            leftIndent=0,
            firstLineIndent=0,
            textColor=INK,
            spaceBefore=4,
        )
    ]
    story += [toc, PageBreak()]

    overview_rows = [["序号", "因子", "类型", "方向", "Rank IC", "Rank ICIR", "Q5-Q1累计收益"]]
    for index, row in enumerate(factors.itertuples(index=False), start=1):
        overview_rows.append([
            str(index),
            row.name_zh,
            row.kind,
            factor_direction(row.rank_ic),
            f"{row.rank_ic:.4f}",
            f"{row.rank_icir:.3f}",
            pct(row.long_short_total_return),
        ])
    story += [
        p("一、因子体系与评价方法", style_map, "H1"),
        p(
            "项目共使用七个日频因子，包括一个题目示例因子、三个原有自建因子和三个新增因子。"
            "每天对每个因子做 1%/99% 横截面缩尾和标准化，并与下一交易日收益进行配对。",
            style_map,
        ),
        table(
            overview_rows,
            style_map,
            widths=[12 * mm, 43 * mm, 24 * mm, 20 * mm, 22 * mm, 24 * mm, 27 * mm],
            font_size=6.4,
        ),
        Spacer(1, 7 * mm),
        table(
            [
                ["指标", "定义", "解释"],
                ["IC", "PearsonCorr(F_t, R_t+1)", "衡量因子值与未来收益的线性关系"],
                ["Rank IC", "SpearmanCorr(F_t, R_t+1)", "衡量横截面排序能力，对极端值更稳健"],
                ["ICIR / Rank ICIR", "均值 / 标准差 x sqrt(252)", "衡量预测方向的稳定性"],
                ["五分组", "按因子值从 Q1 到 Q5 分组", "检查收益是否具有单调结构"],
                ["Q5-Q1", "最高组收益减最低组收益", "负值表示原始因子应反向使用"],
            ],
            style_map,
            widths=[34 * mm, 57 * mm, 77 * mm],
        ),
        Spacer(1, 6 * mm),
        p(
            "注意：表中的 Q5-Q1 按原始因子排序计算。对于 Rank IC 为负的因子，正式组合会通过历史 IC 自动反转方向，"
            "因此原始多空收益为负并不等于该因子不能使用。",
            style_map,
            "Warning",
        ),
        PageBreak(),
    ]

    for index, row in enumerate(factors.itertuples(index=False), start=1):
        explanation = FACTOR_EXPLANATIONS[row.factor]
        direction = factor_direction(row.rank_ic)
        story += [
            p(f"{index + 1}、{row.name_zh}", style_map, "H1"),
            table(
                [
                    ["项目", "内容"],
                    ["因子代码", row.factor],
                    ["类型", row.kind],
                    ["公式", formula_display(row.factor, row.formula)],
                    ["经济含义", explanation["meaning"]],
                    ["变量说明", explanation["variables"]],
                    ["实证方向", f"{direction}；Rank IC = {row.rank_ic:.4f}"],
                ],
                style_map,
                widths=[31 * mm, 137 * mm],
                font_size=7.0,
            ),
            Spacer(1, 6 * mm),
            table(
                [
                    ["有效日期", "IC", "ICIR", "Rank IC", "Rank ICIR", "Rank IC同号率", "Q5-Q1累计收益", "多空最大回撤"],
                    [
                        str(int(row.n_dates)),
                        f"{row.ic:.4f}",
                        f"{row.icir:.3f}",
                        f"{row.rank_ic:.4f}",
                        f"{row.rank_icir:.3f}",
                        pct(row.rank_ic_positive_rate),
                        pct(row.long_short_total_return),
                        pct(row.long_short_max_drawdown),
                    ],
                ],
                style_map,
                widths=[21 * mm, 18 * mm, 20 * mm, 22 * mm, 24 * mm, 25 * mm, 29 * mm, 29 * mm],
                font_size=6.2,
            ),
            Spacer(1, 7 * mm),
            report_image(
                ROOT / "evaluation" / "figures" / f"{row.factor}_diagnostics.png",
                width=170 * mm,
                max_height=105 * mm,
            ),
            p(
                f"图 {index}：左图为 Q1-Q5 平均日收益，中图为各分组累计净值，右图为 20 日滚动 Rank IC。",
                style_map,
                "Caption",
            ),
            PageBreak(),
        ]

    story += [
        p("九、因子相关性与多空净值", style_map, "H1"),
        p(
            "因子相关性用于检查信息是否重复。当前最大绝对相关性约为 0.54，出现在 5 日隔夜-日内收益背离与 5 日反转之间。"
            "二者相关但没有完全重合。",
            style_map,
        ),
        report_image(
            ROOT / "evaluation" / "figures" / "factor_correlation_heatmap.png",
            width=133 * mm,
            max_height=115 * mm,
        ),
        p("图 8：七因子横截面相关系数热力图。", style_map, "Caption"),
        report_image(
            ROOT / "evaluation" / "figures" / "factor_long_short_nav.png",
            width=158 * mm,
            max_height=88 * mm,
        ),
        p("图 9：七个原始因子的 Q5-Q1 多空累计净值。负向因子在组合中需要反转方向。", style_map, "Caption"),
        PageBreak(),
    ]

    story += [
        p("十、七因子组合的风险收益结果", style_map, "H1"),
        p(
            "综合策略只使用截至 t-1 的 20 日历史 IC 确定方向和权重，持有 Top30，每 5 日调仓，"
            "原持仓仍在 Top60 时继续保留。回测包含停牌、涨跌停、卖出手续费、双边滑点和平方根冲击成本。",
            style_map,
        ),
        table(
            [
                ["完整样本指标", "次日收盘成交", "次日开盘成交"],
                ["成本前累计收益", pct(full_close.gross_total_return), pct(full_open.gross_total_return)],
                ["成本后累计收益", pct(full_close.total_return), pct(full_open.total_return)],
                ["年化收益率", pct(full_close.annual_return), pct(full_open.annual_return)],
                ["基准累计收益", pct(full_close.benchmark_total_return), pct(full_open.benchmark_total_return)],
                ["累计主动收益", pct(full_close.active_total_return), pct(full_open.active_total_return)],
                ["年化超额收益", pct(full_close.annual_excess_return), pct(full_open.annual_excess_return)],
                ["年化波动率", pct(full_close.annual_volatility), pct(full_open.annual_volatility)],
                ["最大回撤", pct(full_close.max_drawdown), pct(full_open.max_drawdown)],
                ["夏普率", f"{full_close.sharpe:.3f}", f"{full_open.sharpe:.3f}"],
                ["信息比率", f"{full_close.information_ratio:.3f}", f"{full_open.information_ratio:.3f}"],
                ["日均单边换手", pct(full_close.avg_one_way_turnover), pct(full_open.avg_one_way_turnover)],
            ],
            style_map,
            widths=[62 * mm, 53 * mm, 53 * mm],
        ),
        Spacer(1, 7 * mm),
        report_image(
            ROOT / "backtest" / "figures" / "optimized_strategy_nav_drawdown.png",
            width=170 * mm,
            max_height=118 * mm,
        ),
        p("图 10：七因子正式组合与等权基准净值，以及对应最大回撤过程。", style_map, "Caption"),
        PageBreak(),
        p("风险指标、成本稳健性与留出期", style_map, "H2"),
        report_image(
            ROOT / "backtest" / "figures" / "optimized_strategy_risk_metrics.png",
            width=170 * mm,
            max_height=83 * mm,
        ),
        p("图 11：累计收益、年化收益、年化超额、波动率、回撤、夏普率和信息比率。", style_map, "Caption"),
        report_image(
            ROOT / "backtest" / "figures" / "optimized_period_split.png",
            width=145 * mm,
            max_height=85 * mm,
        ),
        p("图 12：次日开盘策略在前 80% 研究期和最后 20% 留出期的收益对比。", style_map, "Caption"),
        p(
            f"完整样本次日开盘成本后累计收益为 {pct(full_open.total_return)}，达到 25% 目标；"
            f"但最后 20% 留出期绝对收益为 {pct(holdout_open.total_return)}，同期基准为 "
            f"{pct(holdout_open.benchmark_total_return)}，主动收益为 {pct(holdout_open.active_total_return)}。"
            "因此完整样本结果不能解释为未来收益保证。",
            style_map,
            "Warning",
        ),
        PageBreak(),
    ]

    story += [
        p("十一、结论与参考说明", style_map, "H1"),
        bullet("成交额离散度、日内振幅、价量压力和单笔成交规模异常在样本中表现为反向因子。", style_map),
        bullet("主动买卖不平衡冲击表现为负向因子，说明当日异常主动买入在下一日更接近短期反转信号。", style_map),
        bullet("七因子组合通过历史 IC 自动确定方向，避免把同日未来收益用于当日交易。", style_map),
        bullet("正式次日开盘策略成本后累计收益为 32.74%，年化收益为 29.51%，最大回撤为 -11.27%。", style_map),
        bullet("最后 20% 留出期仍跑赢下跌基准，但绝对收益为负，策略还需要在全新日期继续验证。", style_map),
        p("参考材料", style_map, "H2"),
        p(
            "A-2.pdf：参考其因子指标表、策略净值曲线和风险指标的呈现结构。"
            "本报告没有复制参考文件中的因子结论、策略数值或图片，所有公式、指标和图像均来自当前项目代码与输出。",
            style_map,
        ),
        p("复现文件", style_map, "H2"),
        p(
            "因子代码：pipeline_code/construct_factors.py<br/>"
            "因子定义：factors/factor_metadata.csv<br/>"
            "因子汇总：evaluation/factor_summary.csv<br/>"
            "诊断图：evaluation/figures/<br/>"
            "组合指标：backtest/optimized_strategy_metrics.csv<br/>"
            "回测代码：scripts/analyze_factors_backtest.py",
            style_map,
            "Small",
        ),
    ]

    document.multiBuild(story)
    return output


def main() -> None:
    print(build_pdf())


if __name__ == "__main__":
    main()
\end{ProjectCode}
\clearpage
\section{16. \texttt{tests/test\_pipeline.py}}
数据处理、交易约束、因子和LSTM的单元测试。

\textbf{物理行数：} 282\\
\textbf{SHA-256：} \texttt{800db220e735a38cc41af6c163c9a727e74ad192a83bb7e2eb05769739bf8dd0}

\begin{ProjectCode}
from __future__ import annotations

import sys
import tempfile
import unittest
from pathlib import Path

import numpy as np
import pandas as pd
import torch

sys.path.insert(0, str(Path(__file__).resolve().parents[1] / "scripts"))
sys.path.insert(0, str(Path(__file__).resolve().parents[1] / "pipeline_code"))

from analyze_factors_backtest import (
    build_constrained_target,
    equal_weight_benchmark_return,
    estimate_trade_costs,
    factor_signal_scale,
    select_buffered_codes,
)
from import_trade_data import summarize_stock, trading_minutes
from build_kline_csv import FIELDS, build_daily_kline
from evaluate_factors import compute_factor_correlation, neutralize_cross_section
from train_lstm_pytorch import (
    TorchLSTM,
    make_walk_forward_folds,
    metrics_from_probability,
    sequence_summary_features,
)


class ImporterTests(unittest.TestCase):
    def test_trading_minute_grid(self):
        minutes = trading_minutes()
        self.assertEqual(len(minutes), 253)
        self.assertEqual(len(set(minutes)), 253)
        self.assertEqual(minutes[:3], [915, 916, 917])
        self.assertIn(925, minutes)
        self.assertIn(1130, minutes)
        self.assertIn(1300, minutes)
        self.assertEqual(minutes[-1], 1500)
        self.assertNotIn(1200, minutes)

    def test_tick_aggregation_and_adjustment(self):
        ticks = pd.DataFrame({
            "Time": [92500000, 93000000, 93030000],
            "Price": [1000, 1100, 900],
            "Volume": [100, 200, 300],
            "BSFlag": [2, 0, 1],
        })
        with tempfile.TemporaryDirectory() as directory:
            path = Path(directory) / "000001.csv"
            ticks.to_csv(path, index=False)
            daily, minute = summarize_stock(path, adjustment=2.0)
        self.assertEqual(daily["open"], 20.0)
        self.assertEqual(daily["high"], 22.0)
        self.assertEqual(daily["low"], 18.0)
        self.assertEqual(daily["close"], 18.0)
        self.assertEqual(daily["volume"], 600)
        self.assertEqual(daily["trade_count"], 3)
        self.assertEqual(daily["amount"], 5900.0)
        self.assertEqual(daily["buy_volume"], 200)
        self.assertEqual(daily["sell_volume"], 300)
        self.assertEqual(minute.loc[930, "open"], 22.0)
        self.assertEqual(minute.loc[930, "close"], 18.0)
        self.assertEqual(minute.loc[930, "volume"], 500)

    def test_daily_kline_csv_preserves_codes_and_field_order(self):
        dates = ["20250102", "20250103"]
        codes = ["000001", "600000"]
        field_values = {
            "open": [[10.0, 20.0], [11.0, 19.0]],
            "high": [[12.0, 21.0], [12.0, 20.0]],
            "low": [[9.0, 18.0], [10.0, 18.0]],
            "close": [[11.0, 19.0], [10.5, 18.5]],
            "volume": [[100, 200], [120, 210]],
            "trade_count": [[2, 3], [2, 4]],
            "amount": [[1050.0, 3900.0], [1290.0, 3885.0]],
            "buy_volume": [[60, 110], [70, 100]],
            "sell_volume": [[40, 90], [50, 110]],
            "buy_amount": [[630.0, 2145.0], [752.5, 1850.0]],
            "sell_amount": [[420.0, 1755.0], [537.5, 2035.0]],
        }
        with tempfile.TemporaryDirectory() as directory:
            root = Path(directory)
            input_dir = root / "daily"
            input_dir.mkdir()
            for field in FIELDS:
                table = pd.DataFrame(field_values[field], columns=codes)
                table.insert(0, "date", dates)
                table.to_csv(input_dir / f"{field}.csv", index=False)
            output = root / "daily_kline.csv"
            summary = build_daily_kline(input_dir, output)
            result = pd.read_csv(output, dtype={"date": str, "code": str})

        self.assertEqual(summary["rows"], 4)
        self.assertEqual(result.columns.tolist(), ["date", "code", *FIELDS])
        self.assertEqual(result["code"].tolist(), ["000001", "600000", "000001", "600000"])
        self.assertEqual(result.iloc[0]["open"], 10.0)


class LSTMTests(unittest.TestCase):
    def test_pytorch_forward_backward_are_finite(self):
        torch.manual_seed(3)
        generator = torch.Generator().manual_seed(7)
        x = torch.randn(8, 5, 3, generator=generator)
        y = torch.randint(0, 2, (8,), generator=generator).float()
        model = TorchLSTM(input_size=3, hidden_size=4)
        optimizer = torch.optim.Adam(model.parameters(), lr=0.001)
        loss = torch.nn.BCEWithLogitsLoss()(model(x), y)
        optimizer.zero_grad(set_to_none=True)
        loss.backward()
        self.assertTrue(torch.isfinite(loss).item())
        gradients = [parameter.grad for parameter in model.parameters()]
        self.assertTrue(all(gradient is not None for gradient in gradients))
        self.assertTrue(all(torch.isfinite(gradient).all().item() for gradient in gradients))
        before = model.lstm.weight_ih_l0.detach().clone()
        optimizer.step()
        self.assertFalse(torch.equal(before, model.lstm.weight_ih_l0))

    def test_walk_forward_folds_have_non_overlapping_oos_windows(self):
        dates = pd.date_range("2025-01-01", periods=120, freq="D").strftime("%Y%m%d").tolist()
        folds = make_walk_forward_folds(
            dates,
            min_train_dates=60,
            validation_dates=20,
            test_dates=20,
            step_dates=20,
        )
        self.assertEqual(len(folds), 2)
        first_test = set(folds[0]["test_dates"])
        second_test = set(folds[1]["test_dates"])
        self.assertFalse(first_test & second_test)
        self.assertLess(folds[0]["train_dates"][-1], folds[0]["validation_dates"][0])
        self.assertLess(folds[0]["validation_dates"][-1], folds[0]["test_dates"][0])

    def test_classification_metrics_and_logistic_features(self):
        y = np.array([0, 0, 1, 1], dtype=float)
        probability = np.array([0.1, 0.8, 0.7, 0.2], dtype=float)
        metrics = metrics_from_probability(y, probability)
        self.assertEqual((metrics["tn"], metrics["fp"], metrics["fn"], metrics["tp"]), (1, 1, 1, 1))
        self.assertAlmostEqual(metrics["precision"], 0.5)
        self.assertAlmostEqual(metrics["recall"], 0.5)
        self.assertAlmostEqual(metrics["f1"], 0.5)
        sequences = np.arange(2 * 4 * 3, dtype=float).reshape(2, 4, 3)
        self.assertEqual(sequence_summary_features(sequences).shape, (2, 9))


class RiskAndExecutionTests(unittest.TestCase):
    def test_exploratory_factor_signal_scale_is_conservative(self):
        self.assertEqual(factor_signal_scale("amount_mean_sd_log"), 1.0)
        self.assertEqual(
            factor_signal_scale("avg_trade_size_surprise_20d"), 0.10
        )

    def test_rank_buffer_keeps_existing_names_inside_exit_band(self):
        ranked = pd.DataFrame({
            "code": [f"S{rank:02d}" for rank in range(1, 21)],
            "rank": list(range(1, 21)),
        })
        selected = select_buffered_codes(
            ranked,
            {"S11": 0.5, "S12": 0.5},
            portfolio_size=10,
            buffer_exit_rank=20,
        )
        self.assertIn("S11", selected)
        self.assertIn("S12", selected)
        self.assertEqual(len(selected), 10)
        self.assertNotIn("S09", selected)

    def test_industry_market_cap_neutralization_removes_linear_exposure(self):
        rng = np.random.default_rng(9)
        n = 90
        industry = pd.Series(np.repeat(["A", "B", "C"], n // 3))
        log_cap = rng.normal(23.0, 1.0, n)
        market_cap = pd.Series(np.exp(log_cap))
        industry_effect = industry.map({"A": -1.0, "B": 0.5, "C": 1.5}).to_numpy()
        values = pd.Series(2.5 * log_cap + industry_effect + rng.normal(0, 0.1, n))
        residual = neutralize_cross_section(values, industry, market_cap)
        self.assertLess(abs(np.corrcoef(residual, log_cap)[0, 1]), 1e-10)
        means = residual.groupby(industry).mean().abs()
        self.assertTrue(means.lt(1e-10).all())

    def test_factor_correlation_uses_date_code_aligned_panel(self):
        rows = []
        for index in range(20):
            for factor, multiplier in (("f1", 1.0), ("f2", 2.0)):
                rows.append({
                    "date": pd.Timestamp("2025-01-01"),
                    "code": f"{index:06d}",
                    "factor": factor,
                    "factor_z": multiplier * index,
                })
        panel = pd.DataFrame(rows)
        correlation = compute_factor_correlation(panel)
        self.assertAlmostEqual(correlation.loc["f1", "f2"], 1.0)

    def test_execution_constraints_freeze_blocked_positions(self):
        state = pd.DataFrame({
            "amount": [1e8, 1e8, 1e8],
            "suspended": [False, False, False],
            "limit_up": [False, False, True],
            "limit_down": [True, False, False],
            "buyable": [True, True, False],
            "sellable": [False, True, True],
        }, index=["A", "B", "C"])
        target, cash, diagnostics = build_constrained_target(
            ["B", "C"],
            {"A": 0.6, "B": 0.4},
            state,
            portfolio_size=2,
        )
        self.assertAlmostEqual(target["A"], 0.6)
        self.assertAlmostEqual(target["B"], 0.4)
        self.assertNotIn("C", target)
        self.assertAlmostEqual(cash, 0.0)
        self.assertEqual(diagnostics["blocked_buys"], 1)
        self.assertEqual(diagnostics["blocked_sells"], 1)

    def test_slippage_and_impact_costs_are_positive_and_auditable(self):
        state = pd.DataFrame({
            "amount": [20_000_000.0, 50_000_000.0],
            "suspended": [False, False],
            "limit_up": [False, False],
            "limit_down": [False, False],
            "buyable": [True, True],
            "sellable": [True, True],
        }, index=["A", "B"])
        costs, trades = estimate_trade_costs(
            10_000_000.0,
            {"A": 1.0},
            {"A": 0.2, "B": 0.8},
            state,
        )
        self.assertEqual(len(trades), 2)
        self.assertGreater(costs["base_slippage_cost"], 0)
        self.assertGreater(costs["impact_cost"], 0)
        self.assertGreater(costs["sell_fee"], 0)
        self.assertAlmostEqual(
            costs["total_cost"],
            costs["base_slippage_cost"] + costs["impact_cost"] + costs["sell_fee"],
        )
        optimistic, _ = estimate_trade_costs(
            10_000_000.0,
            {"A": 1.0},
            {"A": 0.2, "B": 0.8},
            state,
            base_slippage_bps=0.0,
            impact_coefficient=0.0,
            max_impact_rate=0.0,
        )
        pessimistic, _ = estimate_trade_costs(
            10_000_000.0,
            {"A": 1.0},
            {"A": 0.2, "B": 0.8},
            state,
            base_slippage_bps=10.0,
            impact_coefficient=0.015,
            max_impact_rate=0.03,
        )
        self.assertLess(optimistic["total_cost"], pessimistic["total_cost"])

    def test_equal_weight_benchmark_excludes_blocked_entry_names(self):
        date = pd.Timestamp("2025-01-02")
        returns = pd.DataFrame(
            [[0.1, -0.2, 0.3]], index=[date], columns=["A", "B", "C"]
        )
        state = {
            "amount": pd.DataFrame([[1e8, 1e8, 0.0]], index=[date], columns=returns.columns),
            "suspended": pd.DataFrame([[False, False, True]], index=[date], columns=returns.columns),
            "limit_up": pd.DataFrame([[False, True, False]], index=[date], columns=returns.columns),
            "limit_down": pd.DataFrame([[False, False, False]], index=[date], columns=returns.columns),
        }
        benchmark_return, n_stock = equal_weight_benchmark_return(returns, state, date)
        self.assertEqual(n_stock, 1)
        self.assertAlmostEqual(benchmark_return, 0.1)


if __name__ == "__main__":
    unittest.main()
\end{ProjectCode}
\clearpage
\end{document}
